\documentclass[11pt]{article}
% \usepackage[T1]{fontenc}
\usepackage{lmodern}
\usepackage{microtype}
\usepackage{booktabs}
\usepackage{amssymb, amsmath, amsthm}
\usepackage{mathtools}
\usepackage[colorlinks=true, linkcolor=blue, citecolor=blue, urlcolor=blue]{hyperref}


\usepackage{enumitem}
\usepackage{xcolor}
\usepackage{verbatim}
\usepackage[nameinlink,capitalise]{cleveref}
\usepackage{setspace}








%%%--------
%%%-- Custom Packages
%%%--------
\usepackage{flexart}
\usepackage[many]{tcolorbox} %box environment
\usepackage[normalem]{ulem} %custom underline

% Bar width and spacing
\newlength{\framelinewidth}
\setlength{\framelinewidth}{0.6mm}
\newlength{\marginsep}
\setlength{\marginsep}{3pt}
\newlength{\defparindent}
\setlength{\defparindent}{15pt}
\setlength{\parindent}{\defparindent}
\newlength{\boxparskip}
\setlength{\boxparskip}{\parskip}

%Define colors
% https://davidmathlogic.com/colorblind/#%23038DFC-%232B65C6-%2322AA22-%23003300
\definecolor{aicolor}{HTML}{038DFC} 
\definecolor{editedaicolor}{HTML}{2B65C6}
\definecolor{humancolor}{HTML}{22AA22}  
\definecolor{aiassistedcolor}{HTML}{003300}
\definecolor{defaultcolback}{HTML}{FFFFFF}  % white
\newcommand{\opacityFrame}{0.5} %Opacity Percentage for the Frame
\newcommand{\opacityText}{0.7} %Opacity Percentage for the Label Text

%Define footenote commands: Creates an aux file for the footnotes inside the boxes
\newif\iffootaux
\newcounter{footaux}% for \csname
\newcounter{auxpage}% to handle frontmatter etc.
\setcounter{auxpage}{1}

\makeatletter

\def\tcb@footnote@setup{
\let\footnotetext=\footaux
\let\footnote=\footnoteaux
}
  
\def\footnoteaux{\@ifnextchar[\@xfootnoteaux{\stepcounter{footnote}%
  \footnotemark[\thefootnote]\@footaux[\thefootnote]}}
  
\def\@xfootnoteaux[#1]{\footnotemark[#1]\@footaux[#1]}

\def\footaux{\@ifnextchar[\@footaux{\@footaux[\thefootnote]}}
\def\@footaux[#1]#2{\protected@write{\@auxout}{\let\theauxpage\relax}{\string\newfootnote{\theauxpage}{#1}{#2}}}

\newcommand{\newfootnote}[3]{% #1 = \theauxpage, #2 = \thefoontnote, #3 = text
  \ifnum#1=1\relax
    \footnotetext[#2]{#3}%
  \else
    \stepcounter{footaux}%
    \expandafter\xdef\csname footauxpage\thefootaux\endcsname{#1}%
    \expandafter\xdef\csname footauxnumber\thefootaux\endcsname{#2}%
    \expandafter\protected@xdef\csname footauxtext\thefootaux\endcsname{#3}%
  \fi}

% to run at the end of every page (debug version)
\newcommand{\addfootaux}{\loop
  \@ifundefined{footauxpage\thefootaux}{\footauxfalse}{%
    \count1=\csname footauxpage\thefootaux\endcsname\relax
    \advance\count1 by -1
    \ifnum\count1=\value{auxpage}%
      \expandafter\let\expandafter\footauxnumber\csname footauxnumber\thefootaux\endcsname
      \expandafter\let\expandafter\footauxtext\csname footauxtext\thefootaux\endcsname
      \footnotetext[\footauxnumber]{\footauxtext}%
      \footauxtrue
    \else
      \footauxfalse
    \fi}
  \iffootaux
    \global\expandafter\let\csname footauxpage\thefootaux\endcsname\relax
    \global\expandafter\let\csname footauxnumber\thefootaux\endcsname\relax
    \global\expandafter\let\csname footauxtext\thefootaux\endcsname\relax
    \stepcounter{footaux}%
  \repeat
\stepcounter{auxpage}}
\makeatother

\AtBeginDocument{\setcounter{footaux}{1}}% reset counter after \newfootnote
\AddToHook{shipout/after}{\addfootaux}


%Define box styles
\makeatletter
\tcbset{
    roundbox/.style={
    % Box Options
    rounded corners,
    arc=3mm,
    boxrule=\framelinewidth
  },
  rectanglebox/.style={
    % Box Options
    sharp corners,  
    boxrule=\framelinewidth
  },
    linemargin/.style={
    % Box Options
    % Make frame with 1 line
    sharp corners,   
    leftrule=\framelinewidth,    % Thickness of the left line
    rightrule=-0.1mm,            % Removes right frame
    toprule=-0.1mm,              % Removes top frame
    bottomrule=-0.125mm,         % Removes bottom frame
    % % make white background
    colback=white,  
    % Oversize makes bigger, so that text alings with rest of the document
    oversize,
    % box margins 
    left=\marginsep,                                
    right=\marginsep,
    top=0.5pt,
    bottom=0.5pt,
    before skip=0.5\baselineskip,
    after skip=0.5\baselineskip
  },
  alignoptions/.style={
    % box margins 
    left=\marginsep, 
    right=\marginsep,
    top=0.5pt,
    bottom=0.5pt,
    before skip=0.5\baselineskip,
    after skip=\parskip,
    before upper={
    \parindent=\defparindent
    \setlength{\parskip}{\boxparskip}
    }
  },
  %Format for footnotes to appear in the correct page
  footnoteformat/.style={
  before upper app=\tcb@footnote@setup
  }
}
\makeatother

\tcbset{
  enhancebreak/.style={
    % Box Options
    enhanced,
    breakable,    
    skin first=enhanced,
    skin middle=enhanced,
    skin last=enhanced,
    % Oversize makes bigger, so that text alings with rest of the document
    oversize,
  },
}

% Define Tcolorbox environments
%% Args: #1 additional options, #2 color, #3 Side text label
\newtcolorbox{labeledblock}[3][]{
    colframe=#2,opacityframe=\opacityFrame,
    enhancebreak, alignoptions, linemargin, 
    colback=white,
    overlay={
    \node[
    #2, font=\tiny\bfseries\sffamily, 
    rotate=90, anchor=east, align=right,
    text opacity=\opacityText
      ]
      at ([xshift=-\marginsep-\framelinewidth/2, yshift=3pt]frame.north west) {#3};
    }, footnoteformat,
    #1
}

\newtcolorbox{humanblock}[1][]{
    colframe=humancolor, opacityframe=\opacityFrame,
    enhancebreak, alignoptions, linemargin,
    colback=white,
    overlay={
    \node[
    humancolor, font=\tiny\bfseries\sffamily, 
    rotate=90, anchor=east, align=right, 
    text opacity=\opacityText
      ]
      at ([xshift=-\marginsep-\framelinewidth/2, yshift=3pt]frame.north west) {Human};
    }, footnoteformat,
    #1
}

\newtcolorbox{aiassistedblock}[1][]{
    colframe=aiassistedcolor, opacityframe=\opacityFrame,
    enhancebreak, alignoptions, linemargin,
    colback=white,
    overlay={
    \node[aiassistedcolor, font=\tiny\bfseries\sffamily, rotate=90, anchor=east, align=right,
      text opacity=\opacityText
      ]
      at ([
      xshift=-\marginsep-\framelinewidth/2-4pt, yshift=3pt
      ]frame.north west) {Human: \\ AI-assisted};
    }, footnoteformat,
    #1
}

\newtcolorbox{aiblock}[1][]{
    colframe=aicolor, opacityframe=\opacityFrame,
    enhancebreak, alignoptions, linemargin, 
    colback=white,
    overlay={
    \node[aicolor, font=\tiny\bfseries\sffamily, rotate=90, anchor=east, align=right,
      text opacity=\opacityText
      ]
      at ([xshift=-\marginsep-\framelinewidth/2, yshift=3pt]frame.north west) {AI};
    }, footnoteformat,
    #1
}

\newtcolorbox{editedaiblock}[1][]{
    colframe=editedaicolor, opacityframe=\opacityFrame,
    enhancebreak, alignoptions, linemargin,
    colback=white,
    overlay={
    \node[
    editedaicolor, font=\tiny\bfseries\sffamily, 
    rotate=90, anchor=east, align=right, 
    text opacity=\opacityText
      ]
      at ([xshift=-\marginsep-\framelinewidth/2, yshift=3pt]frame.north west) {Edited AI};
    }, footnoteformat,
    #1
}


%%%---------------------- Underline Commands -------------------------

% Dash pattern:
\gdef\dashpattern{\kern.05em\lower.3ex\vbox{\hrule width.2em\kern.1ex}}

\def\humanul{\textcolor{humancolor}\bgroup\markoverwith{\kern.05em\lower.3ex\vbox{\hrule width.2em\kern.1ex}\kern.05em}\ULon}

\def\aiul{\textcolor{aicolor}\bgroup\markoverwith{\kern.05em\lower.3ex\vbox{\hrule width.2em\kern.1ex}\kern.05em}\ULon}
\newcommand{\Muse}{with {\itshape Muse Spark}}


%% Commands
%%%_____________________________________


\numberwithin{equation}{section}
\newtheorem{theorem}{Theorem}[section]
\newtheorem{proposition}[theorem]{Proposition}
\newtheorem{lemma}[theorem]{Lemma}
\newtheorem{corollary}[theorem]{Corollary}
\theoremstyle{definition}
\newtheorem{definition}[theorem]{Definition}
\theoremstyle{remark}
\newtheorem{remark}[theorem]{Remark}

%%%%%%%%%%%%%%


%%% Macros
%%%_____________________________________


\newcommand{\R}{\mathbb R}
\newcommand{\E}{\mathbb E}
\newcommand{\Pp}{\mathbb P}
\newcommand{\1}{\boldsymbol 1}
\newcommand{\cA}{\mathcal A}
\newcommand{\cB}{\mathcal B}
\newcommand{\cF}{\mathcal F}
\newcommand{\cQ}{\mathcal Q}
\newcommand{\cT}{\mathcal T}
\newcommand{\PiZero}{\Pi_n}
\newcommand{\Sym}{\operatorname{Sym}}
\newcommand{\Tr}{\operatorname{Tr}}
\newcommand{\supp}{\operatorname{supp}}
\newcommand{\rank}{\operatorname{rank}}
\newcommand{\diag}{\operatorname{diag}}
\newcommand{\res}{\operatorname{res}}
\newcommand{\rvec}{\res}
\newcommand{\rentry}[1]{\res_{#1}}
\newcommand{\op}{\operatorname{op}}
\newcommand{\Frob}{F}
\newcommand{\whp}{\text{with high probability }}
%%%%%%%%%%%%%%%%%%



%%%%%%%%%%%%%%%%%%%
%%% ---- Flexart Options
%%%%%%%%%%%%%%%%%%%





\title[Strict Threshold for Gaussian Ellipsoid Fitting]{The Strict Threshold for Gaussian Ellipsoid Fitting}


\primaryauthors{
\MakeUppercase{Aykut Arslan}
\\
\Muse
}

\newcommand{\secondaryAuthorsLabel}{Mathematical review by}

\secondaryauthors{
\MakeUppercase{Babak Modami} 
\and
\MakeUppercase{Alexander Roitershtein}
\and
\MakeUppercase{Mark Sepanski}
\and
\MakeUppercase{Grigory Sokolov}
}

\flexset{
authorwidth      = 0.7\textwidth,
abovesecondary   = 0.3em,
secondarylabel   = {\secondaryAuthorsLabel},
partrulelength   = 0.75\textwidth,
}

\subjclass[2020]{60B20, 52A20, 90C22}



\flexabstract{%
Let $x_1,\ldots,x_n$ be independent $N(0,I_d)$ vectors.
We study whether there exists a positive semidefinite matrix $S$
such that $x_i^T S x_i=d$ for every $1\le i\le n.$ 
We prove a sharp threshold at $n\sim d^2/4$. 
\par
If $\limsup_{d\to\infty}\frac{n(d)}{d^2}<\frac14,$
then with probability tending to one such an $S$ exists, 
with $S$ in fact positive definite. 
\par
If $\liminf_{d\to\infty}\frac{n(d)}{d^2}>\frac14,$
then with probability tending to one no such $S$ exists. 
We make no claim when $n(d)/d^2\to1/4$.
}












\begin{document}


\maketitle

\vspace{-0.5em}


\setcounter{tocdepth}{1} 
\tableofcontents


\vspace{2em}

%%%%%%%%%%%%%
%%%%%%%%%%%%%
%%%%%%%%%%%%%

%%%%%%%%%%%%%%%%%%%%
\section{Introduction and statement of the main result} 
%%%%%%%%%%%%%%%%%%%%

\begin{humanblock}


We consider the following problem. Given $n$ random points in $\R^d$, when do they lie on the boundary of an ellipsoid? For small $n$ and a standard Gaussian setting such an event has a high likelihood, while for sufficiently large $n$ it is unlikely to occur. The main regime of interest is therefore near the phase transition, which turns out to be $n\asymp d^2$ \cite{BandeiraMaillard2025, MaillardKunisky2024, SaundersonParriloWillsky2013}.

We study the random setting in which
\[
  x_1,\ldots,x_n\sim N(0,I_d)
\]
are independent. Since $\|x_i\|_2^2$ is typically close to $d$, the sphere already gives an approximate fit. The question is whether it can be deformed into a centered ellipsoid so that all of the points lie exactly on its boundary.

Saunderson, Parrilo, and Willsky \cite{SaundersonParriloWillsky2013} conjectured that this problem has a sharp transition at $n=d^2/4$. More precisely, for every fixed $\varepsilon>0$, they predicted feasibility of such an ellipsoid with high probability when
\[
  n\le\left(\frac14-\varepsilon\right)d^2,
\]
and infeasibility with high probability when
\[
  n\ge\left(\frac14+\varepsilon\right)d^2.
\]

The value $1/4$ comes from the geometry of the positive semidefinite cone. One way to measure the size of this cone is to take a random Gaussian symmetric matrix, project it onto the cone, and compute the expected squared size of that projection. For $d\times d$ matrices, this quantity is $d(d+1)/4$, making $1/4$ the natural candidate for the asymptotic threshold.

% The constant $1/4$ reflects the statistical dimension of the positive semidefinite cone, which is asymptotically $d^2/4$. The ellipsoid-fitting problem arose in minimum trace factor analysis and vector discrepancy, and related formulations have also appeared in independent component analysis and sum-of-squares.



Earlier work established feasibility in increasingly large ranges of $n$, eventually reaching $n\le cd^2$ for a
sufficiently small constant $c>0$. Maillard and Kunisky \cite{MaillardKunisky2024} predicted the sharp transition at $1/4$ using the non-rigorous replica method of statistical physics. Bandeira and Maillard \cite{BandeiraMaillard2025} later proved the $1/4$ threshold for an approximate version of the ellipsoid-fitting problem. The corresponding exact fitting problem remained open.






We prove the exact threshold with fixed margin. No assertion is made for the critical case when $n/d^2\to1/4$.

\begin{theorem}[Strict-threshold theorem]\label{thm:main}
  Let $n=n(d)$ be any sequence of positive
  integers. For each $d$, let
  $x_1,\ldots,x_n$ be independent
  $N(0,I_d)$ vectors and define
  \[
    \cF_d=\{S\succeq0:S^\top = S \text{~and~}
    x_i^TSx_i=d\text{ for all }1\le i\le n\}.
  \]
  Then, the following statements hold.
  \begin{enumerate}[label=\textup{(\roman*)}]
    \item If $\limsup_{d\to\infty}n/d^2<1/4$,
    then
    \[
      \Pp\bigl(\text{there exists }S\succ0
      \text{ with }x_i^TSx_i=d \text{ for every }
      i\bigr)\longrightarrow1.
    \]
    \item If $\liminf_{d\to\infty}n/d^2>1/4$,
    then $\Pp(\cF_d\ne\varnothing)\to0$.
  \end{enumerate}
  
\end{theorem}
\end{humanblock}


\section*{Concurrent work}

\begin{humanblock}
While completing this work, we became aware of the concurrent work of Misiakiewicz and Wen \cite{MisiakiewiczWen2026}, posted on August 10, 2026. 

  
The two works prove the same sharp threshold, but use different proof architectures. Misiakiewicz and Wen build directly on the Gaussian-equivalence framework of Bandeira--Maillard. Below the threshold, they use a head--tail decomposition of the dual vector, with exact treatment of the sparse head and Gaussian comparison for the low-influence tail. Above the threshold, they combine a low-rank/diffuse decomposition with conditional Gaussianization and a projected Gordon escape argument.


Our proof follows a different route. Below the threshold, we begin with the approximate fitting result of Bandeira--Maillard and perturb it to an exact fit using a spectrally constrained robust-loss construction, a sparse/diffuse decomposition of the coefficient vector, and a dyadic multiscale estimate. Above the threshold, we also use Gordon's Gaussian min--max theorem, but in a different role. We first establish the Gaussian margin and then transfer it to centered
rank-one measurements through a leave-one-row Gibbs interpolation, in the spirit of rigorous cavity methods
\cite{AizenmanSimsStarr2003, Talagrand2011}. A Grassmann-net conditioning argument and a finite large-eigenvalue decomposition then extend the diffuse estimate to the full positive semidefinite cone.
 
\end{humanblock}


\section*{Statement on AI Use}

\begin{humanblock}

Meta AI played a substantial collaborative role in the research leading
to this manuscript. The proof was developed through many rounds of
human-directed interaction rather than a single autonomous AI run.
Meta AI proposed proof strategies, derived intermediate lemmas and
calculations, explored alternative approaches, produced candidate
arguments, responded to mathematical objections, and revised or
replaced arguments when gaps were identified. Many proposed arguments
did not survive further checking and were rejected, substantially
changed, or replaced.

To make this process transparent, the manuscript uses 
{
\bfseries\sffamily \color{aicolor!70}AI
}
to mark material primarily generated by Meta AI and 
{
\bfseries\sffamily \color{humancolor!70}Human
}
to mark material primarily authored by human researchers. These labels describe the provenance of the written material, not the level of mathematical verification.

Review and collaboration with other mathematicians were essential throughout this process. The author discussed the developing proof with reviewers and collaborators, shared drafts for examination, received mathematical objections and suggestions, and used this feedback to identify parts of the argument requiring further work. In several cases, this feedback led to new rounds of mathematical analysis, clarification, and revision.

The author selected the problem, set the research goals, decided
which proof directions to pursue or abandon, evaluated competing
arguments, and determined which parts of the proof required further
investigation. Meta AI outputs were not accepted simply because it described an argument as correct or complete. AI-generated
claims of correctness were treated as research suggestions rather
than mathematical verification.

Meta AI also contributed to the preparation and revision of the mathematical exposition. The discovery, derivation, criticism, organization, and writing of parts of this manuscript emerged from a collaborative process involving Meta AI, the author, and reviewers. Because these contributions were closely intertwined throughout the project, no numerical percentages of contribution are assigned.

The author takes responsibility for all mathematical claims in this
manuscript. The final arguments are presented on the basis of the
author's judgment together with extensive human review and
collaboration. AI-generated arguments or assessments of correctness
are not treated as substitutes for mathematical verification.

\end{humanblock}





%%%%%%%%%%%%%%%%%%%%
\section{Overview and notation}
\label{deterministic-framework}
%%%%%%%%%%%%%%%%%%%%
\begin{humanblock}
    
% We work with the notation introduced above.
% The main asymptotic is $d\to\infty$ with
% $n=n(d)$. All other parameters are regarded as
% fixed before $d\to\infty$. In the sequel we will
% introduce a fixed gap to $1/4$, clipping levels,
% net radii, and the diffuse cap $\tau$, among
% others. Constants may depend on these auxiliary
% parameters but not on $d$ or $n$; their values
% are specified when first used.


For $k\ge1$, let $\Sym_k$ denote the space of real symmetric $k\times k$ matrices, equipped with the Frobenius inner product
\[
  \langle A,B\rangle=\Tr(AB).
\]
Let $\|A\|_F$, $\|A\|_{op}$, and $\|A\|_*$ denote the
Frobenius, operator, and nuclear norms, respectively, with
\[
  \|A\|_*=\sup_{\|W\|_{op}\le1}\langle W,A\rangle,
\]
and write $A\succeq0$ when $A$ is positive semidefinite.
Let $\1_n$ denote the all-ones vector and
\[
  \PiZero=I_n-\frac1n\1_n\1_n^T
\]
the centering projector. Let $\supp$ denote support, $h=\lfloor d/2\rfloor$, and let $O_{\Pp}$ and $o_{\Pp}$  denote the usual probabilistic Landau symbols \cite{vaart}. Let $N(T,\varepsilon B)$ denote covering numbers and $w(T)$ Gaussian width \cite{Vershynin2018}.

Let $G$ be a centered Gaussian random matrix in $\Sym_p$ such that
\[
  \E\langle G,A\rangle\langle G,B\rangle
  =\langle A,B\rangle
  =\Tr(AB)
\]
for all $A,B\in\Sym_p$, and let
\[
  \mathcal P_p=\{D\in\Sym_p:D\succeq0,\|D\|_F=1\}.
\]




For $A\in\Sym_d$, let
\[
  \cA_d(A)_i=x_i^TAx_i,\qquad
  \cQ_d(A)_i=\frac{x_i^TAx_i-\Tr A}{\sqrt2}.
\]
For $S\in\Sym_d$, let
\[
  \rentry{i}(S)=\frac{x_i^TSx_i-d}{\sqrt d},\qquad
  \cT_d(S)=\frac1{\sqrt d}\PiZero\cA_d(S).
\]
For $y\perp\1_n$, let
\[
  \cT_d^*y=\frac1{\sqrt d}\sum_iy_ix_ix_i^T
  =\frac1{\sqrt d}\sum_iy_i(x_ix_i^T-I_d).
\]


For each $i$, let $a_i,b_i\in\mathbb R^h$ denote the first and last $h$ coordinate blocks of $x_i$. The vectors are mutually independent and each has distribution $N(0,I_h)$. If $d$ is odd, the middle coordinate is discarded. Let
\[
  \cB y=\frac1{\sqrt d}\sum_iy_ia_ib_i^T.
\]

We next briefly outline our proof strategy. Let $S\succeq0$, $S\neq0$, and suppose
$\PiZero\cQ_d(S)=0$. Then $x_i^TSx_i=c$ is
independent of $i$. If in addition $n\ge d$ and
$\sum_ix_ix_i^T\succ0$, then $c>0$ and
$(d/c)S\in\cF_d$. When $n\ge d$, Gaussian vectors
span $\R^d$ a.s., so $\sum_ix_ix_i^T\succ0$ a.s.
Thus above $1/4$, after discarding finitely many
$d$, it suffices to show the centered kernel has
no nonzero PSD matrix.

% The following terminology is defined formally
% later; we give a brief guide.

% \begin{description}
%   \item[Huber selector] Minimizer of
%     $(1/n)\sum_i\phi_{\rm Hub}(\rentry{i}(S))
%     +R_d(S)$ over $\lambda_-I\preceq S\preceq
%     \lambda_+I$, where $R_d=\tau_d\|S-I\|_F^2/(2d)$,
%     $\tau_d=d^{-2}$,
%     \cref{the-huber-selector-and-its-residual-energy}.
%   \item[Large/small] For a coefficient vector $y$,
%     $y_i^{\rm large}=y_i\mathbf1_{\{|y_i|>t\}}$ and
%     $y_i^{\rm small}=y_i\mathbf1_{\{|y_i|\le t\}}$.
%   \item[Cap/flat] $v$ is flat if
%     $\sqrt d\|v\|_\infty\le\theta\|v\|_2$, i.e.\@
%     no coordinate dominates.
%   \item[Dyadic record] For $u$, let
%     $B_t=\{i:2^{-t-1}\lambda_{\max}<|u_i|
%     \le2^{-t}\lambda_{\max}\}$; select next record
%     when $|B_{t_{r+1}}|\ge q|B_{t_r}|$, $q>1$.
%   \item[Effective support] $R(u)
%     =\|u\|_1^2/\|u\|_2^2$.
%   \item[Diffuse] $D\succeq0$ with $\|D\|_F=1$ and
%     $\|D\|_{op}\le\tau$.
%   \item[Peeling] $S=H+R$ with $HR=0$,
%     $\rank(H)\le K$, either $\|R\|_F\le\delta$ or
%     $R/\|R\|_F$ diffuse, $K=O_{\tau,\delta}(1)$.
%   \item[Cavity Gibbs] Log-sum-exp over diffuse
%     after deleting row $j$, independent of row $j$
%     given others,
%     \cref{a-common-leave-one-row-cavity-law}.
%   \item[Nuisance] $L\subset\R^n$ with
%     $\dim L=O_K(1)$, spanned by $\1_n$ and residual
%     trace plus leading eigenspace, deterministic
%     after conditioning.
% \end{description}

For the lower bound, fix $\alpha_0$ with
$\limsup n/d^2<\alpha_0<1/4$ and enlarge to
$n=\lfloor\alpha_0d^2\rfloor$ by adding
independent $N(0,I_d)$ vectors; a fit for the
larger sample induces one for the smaller. For
the upper bound, fix $1/4<\alpha_1<\liminf n/d^2$
and keep the first $\lfloor\alpha_1d^2\rfloor$
points; they remain i.i.d.\ and
$\cF_{n_2}\subseteq\cF_{n_1}$ for $n_1\le n_2$, so
infeasibility for $\lfloor\alpha_1d^2\rfloor$
points implies that for $n$ points. Hence
$n=\Theta(d^2)$ suffices. Below $1/4$, we start
from the approximate fit in
\cref{the-published-approximate-comparator} and
upgrade it with the Huber selector to
$o_{\Pp}(d)$ $\ell_2$ and $O_{\Pp}(\sqrt{\log d})$
$\ell_\infty$ residuals; we prove
$\|\cT_d^*y\|_*$ large via determinant small-ball
for flat and dyadic multiscale for sparse, mixed
by thresholding, and duality gives $o_{\Pp}(1)$
correction $S_d^*+H\succ0$.

Above $1/4$, let $\mathcal P_p=\{D\succeq0:\|D\|_F=1\}$ with
Gaussian width $w(\mathcal P_p)^2\le p(p+1)/4$,
so $(w/p)^2\to1/4$. Let $A:\Sym_p\to\R^r$, $(AD)_j=\langle G_j,D\rangle$, where
$G_1,\ldots,G_r$ are independent centered Gaussian random matrices in $\Sym_p$ with the covariance specified above, and $r\ge(1/4+\varepsilon)p^2$. Then Gordon's
escape-through-a-mesh \cite{Gordon1988} gives
$\inf_{D\in\mathcal P_p}\|AD\|_2\ge\gamma p$
with high probability. The free energy $F_X$ in
\cref{deterministic-nets-and-free-energy-inequalities}
stabilizes at positive level; row-deleted
interpolation in
\cref{a-common-leave-one-row-cavity-law} transfers
this Gaussian bound to rank-one rows
$Z_i=(w_iw_i^\top-I_p)/\sqrt2$ for
$\|D\|_{op}\le\tau$, and peeling
$S=H+R$, $\rank(H)\le K$, in
\cref{fixed-rank-peeling-and-unsat-above-one-quarter}
covers the full PSD cone via fixed-rank and
diffuse cases.
\end{humanblock}



% %%%%%%%%%%%%%%%%%%%% -Part-
% \part{Feasibility below the
%   threshold}\label{part-feasibility-below-the-threshold}
% %%%%%%%%%%%%%%%%%%%%


%%%%%%%%%%%%%%%%%%%% 
\section{Previous work: Approximate-fitting
theorem}\label{the-published-approximate-comparator}
%%%%%%%%%%%%%%%%%%%%
\begin{humanblock}

    


We use the groundbreaking approximate-fitting
theorem of Bandeira and Maillard
\cite[Theorem 1.4]{BandeiraMaillard2025}  which asserts that if
$\limsup n/d^2<1/4$, then there are constants
\[
  0<\lambda_-\le\lambda_+<\infty
\]
such that, for every $1\le r<4/3$ and $n=n(d)$,
\[
  \min_{\lambda_-I\preceq S\preceq\lambda_+I}
  \frac1n\sum_{i=1}^n
  \left|\sqrt d\left(\frac{x_i^TSx_i}{d}-1\right)
  \right|^r\xrightarrow{\Pp}0 \text{ as $d \to\infty$.}
\]
The residual inside the absolute value is exactly
$\rentry{i}(S)$. Taking $r=1$ gives a measurable
approximate solution $\widetilde S_d$ in the fixed
spectral interval
\[
  \mathcal S_d=
  \{S:\lambda_-I\preceq S\preceq\lambda_+I\}
\]
such that
\[
  \frac1n\sum_i|\rentry{i}(\widetilde S_d)|
  \xrightarrow{\Pp}0.
\]
The comparator may depend on the full sample. It
is used only in pathwise objective comparisons,
so no independence is required.

\end{humanblock}

%%%%%%%%%%%%%%%%%%%%
\section{A convex robust-loss estimator and its
residual bounds}\label{the-huber-selector-and-its-residual-energy}
%%%%%%%%%%%%%%%%%%%%

\begin{aiblock}  

Let
\[
  \phi_{\rm Hub}(t)=
  \begin{cases}
    t^2/2,&|t|\le1,\\
    |t|-1/2,&|t|>1,
  \end{cases}
\]
and let
\[
  R_d(S)=\frac{\tau_d}{2d}\|S-I\|_F^2,
  \ \tau_d=d^{-2}.
\]
Define the unique selector
\[
  S_d^*=\operatorname*{argmin}_{S\in\mathcal S_d}
  \left\{\frac1n\sum_i\phi_{\rm Hub}(\rentry{i}(S))
  +R_d(S)\right\}.
\]
For all sufficiently large $d$, one has $n\geq2$. For such
$d$ and each $1\leq j\leq n$, define the deleted selector
\[
  S_d^{(-j)}=\operatorname*{argmin}_{S\in\mathcal S_d}
  \left\{\frac1{n-1}\sum_{i\ne j}
  \phi_{\rm Hub}(\rentry{i}(S))+R_d(S)\right\}.
\]
Compactness gives existence, ridge gives
uniqueness and measurability. In particular
$S_d^{(-j)}$ is independent of $x_j$.

\begin{proposition}
\label{main-proposition-3-1}
  If $n/d^2\to\alpha_0$ for fixed
  $0<\alpha_0<1/4$, then
  \[
    \frac1n\sum_i\rentry{i}(S_d^*)^2
    \xrightarrow{\Pp}0,
  \]
  \[
    \|\PiZero\rvec(S_d^*)\|_2=o_{\Pp}(d),
  \]
  and
  \[
    \|\PiZero\rvec(S_d^*)\|_\infty
    =O_{\Pp}(\sqrt{\log d}).
  \]
\end{proposition}

\begin{proof}
  Put
  \[
    H_d(S)=\frac1n\sum_i\phi_{\rm Hub}(\rentry{i}(S)).
  \]
  Since $\phi_{\rm Hub}(t)\le|t|$ and
  $R_d(\widetilde S_d)$ is $O(d^{-2})$, optimality
  gives $H_d(S_d^*)\xrightarrow{\Pp}0$.
  Set $a_j=\rentry{j}(S_d^*)$,
  $b_j=\rentry{j}(S_d^{(-j)})$. Write
  $G_{-j}(S)=\frac1n\sum_{i\ne j}\phi_{\rm Hub}(\rentry{i}(S))+\frac{n-1}{n}R_d(S)$.
  The deleted selector minimizes $G_{-j}$. Comparing
  the two exact optimality inequalities gives
  $\phi_{\rm Hub}(a_j)+R_d(S_d^*)\le\phi_{\rm Hub}(b_j)+R_d(S_d^{(-j)})$.
  Let
  $b_0=\max\{|\lambda_--1|,|\lambda_+-1|\}$,
  $\epsilon_d=\frac{b_0^2\tau_d}{2}$.
  Then $\phi_{\rm Hub}(a_j)\le\phi_{\rm Hub}(b_j)+\epsilon_d$.
  Comparison of
  $G_{-j}$ minimizers gives
  $|a_j|\le|b_j|+e_d$. Put
  $e_d=\phi_{\rm Hub}^{-1}(\epsilon_d)$. For all
  sufficiently large $d$, one has $b_0/d\leq1$, and
  hence
  \[
    e_d=\frac{b_0}{d}.
  \]

  For $M\ge2e_d$,
  \[
    \frac1n\sum_j a_j^2\mathbf1_{\{|a_j|>M\}}
    \le\frac2n\sum_j b_j^2\mathbf1_{\{|b_j|>M/2\}}+2e_d^2.
  \]
  We next prove uniform integrability of the
  deleted residuals. Put
  \[
    \widehat\Sigma_{-j}=\frac1{n-1}\sum_{i\ne j}x_ix_i^T,
    \quad C_d=\max_j\|\widehat\Sigma_{-j}-I\|_{op}.
  \]
  Deleted optimality against $\widetilde S_d$ gives,
  simultaneously for every $j$,
  \[
    \frac1{n-1}\sum_{i\ne j}\phi_{\rm Hub}(\rentry{i}(S_d^{(-j)}))
    +R_d(S_d^{(-j)})
    \le\frac{n}{n-1}H_d(\widetilde S_d)+R_d(\widetilde S_d)=:r_d.
  \]
  The right side does not depend on $j$ and satisfies
  $r_d\xrightarrow{\Pp}0$. Using
  \[
    |t|\le\sqrt{2\phi_{\rm Hub}(t)}+2\phi_{\rm Hub}(t),
  \]
  and Jensen gives, uniformly in $j$,
  \[
    \frac1{n-1}\sum_{i\ne j}|\rentry{i}(S_d^{(-j)})|\le a_d,
  \]
  where $a_d=\sqrt{2r_d}+2r_d$, so $a_d\xrightarrow{\Pp}0$.
  Since $0\preceq S_d^{(-j)}\preceq\lambda_+I$,
  \[
    \frac1{n-1}\sum_{i\ne j}\rentry{i}(S_d^{(-j)})
    =\frac{\operatorname{Tr}(S_d^{(-j)}\widehat\Sigma_{-j})-d}{\sqrt d}.
  \]
  Therefore
  \[
    \max_j\left|\frac{\operatorname{Tr}S_d^{(-j)}}{d}-1\right|
    \le\lambda_+C_d+\frac{a_d}{\sqrt d}.
  \]
  Let $X$ be $n$ by $d$ with rows $x_i^T$. Gaussian
  singular-value \cite[Thm II.13]{DavidsonSzarek2001}
  with $t=\sqrt d$ puts every singular value of
  $X/\sqrt n$ between $1-2\sqrt{d/n}$ and
  $1+2\sqrt{d/n}$ except $2e^{-d/2}$. Hence
  $\|n^{-1}X^TX-I\|_{op}=O_{\Pp}(d^{-1/2})$. Chi-square
  Chernoff and union over $n=O(d^2)$ give
  $\max_j\|x_j\|_2^2=O_{\Pp}(d)$. Finally,
  \[
    \widehat\Sigma_{-j}-I=\frac{n}{n-1}(\widehat\Sigma-I)
    +\frac{I-x_jx_j^T}{n-1}
  \]
  gives $C_d=O_{\Pp}(d^{-1/2})$. Hence
  \[
    \mu_j=\frac{\operatorname{Tr}S_d^{(-j)}-d}{\sqrt d}
  \]
  satisfies $\max_j|\mu_j|=O_{\Pp}(1)$.

  Conditionally on $x_{-j}$,
  \[
    b_j=Z_j+\mu_j,\quad
    Z_j=\frac{x_j^TS_d^{(-j)}x_j-\operatorname{Tr}S_d^{(-j)}}{\sqrt d}.
  \]
  Hanson--Wright \cite[Thm 1.1]{RudelsonVershynin2013}
  gives $c_{\rm HW}=c_{\rm HW}(\lambda_+)>0$ with
  \[
    \Pp(|Z_j|>t\mid x_{-j})\le2\exp[-c_{\rm HW}\min\{t^2,t\sqrt d\}].
  \]
  For $S\in\mathcal S_d$, $x\sim N(0,I_d)$ set
  $Z_S=(x^TSx-\operatorname{Tr}S)/\sqrt d$. For fixed $K$,
  \[
    \rho_{\lambda_+,K}(L)=\sup_{S\in\mathcal S_d,|\mu|\le K}
    \E[(Z_S+\mu)^2\mathbf1_{\{|Z_S+\mu|>L\}}]\to0
  \]
  as $L\to\infty$. Let $E_{d,K}=\{\max_j|\mu_j|\le K\}$.
  Then
  \[
    \E\left[1_{E_{d,K}}\frac1n\sum_j b_j^2\mathbf1_{\{|b_j|>L\}}\right]
    \le\rho_{\lambda_+,K}(L).
  \]
  Markov and full-to-deleted comparison give empirical
  UI of $a_j^2$.

  For fixed $M\ge1$, $t^2\mathbf1_{\{|t|\le M\}}\le2M\phi_{\rm Hub}(t)$,
  so
  \[
    \frac1n\sum_j a_j^2\mathbf1_{\{|a_j|\le M\}}\le2M H_d(S_d^*)\xrightarrow{\Pp}0.
  \]
  First choose $K$ with $\Pp(E_{d,K}^c)$ small, then $M$
  with $\rho_{\lambda_+,K}(M/2)$ small, then $d\to\infty$.
  This proves $\frac1n\sum_j\rentry{j}(S_d^*)^2\xrightarrow{\Pp}0$.
  Cauchy--Schwarz and contractivity give
  \[
    \frac{\|\PiZero\rvec(S_d^*)\|_2}{d}\le\frac{\sqrt n}{d}
    \left(\frac1n\sum_i\rentry{i}(S_d^*)^2\right)^{1/2}\xrightarrow{\Pp}0.
  \]
  Finally $t=C\sqrt{\log d}$, $n\le d^2$ eventually, union bound
  \[
    \Pp(\max_j|Z_j|>C\sqrt{\log d})\le2d^2\exp[-c_{\rm HW}C^2\log d]
  \]
  for large $d$. Choosing $C$ large gives
  $\max_j|Z_j|=O_{\Pp}(\sqrt{\log d})$, hence
  $\max_j|\rentry{j}(S_d^*)|=O_{\Pp}(\sqrt{\log d})$. Since
  $\|\PiZero z\|_\infty\le2\|z\|_\infty$, L$\infty$ assertion follows.
\end{proof}
\end{aiblock}



%%%%%%%%%%%%%%%%%%%%
\section{The arbitrary-center small-ball
estimate}\label{expert-determinant-module}
%%%%%%%%%%%%%%%%%%%%

\begin{aiblock}

The next result is the quantitative small-ball
estimate used on the small-coordinate part of a
coefficient vector. We prove it here before it is
used. The matrix center is arbitrary and
deterministic. Later, after conditioning and
discretization, this center will represent the
contribution of the large coordinates.

%%%%%%%%%%%%%%%%%%%%
\subsection{Capped determinant moments}
%%%%%%%%%%%%%%%%%%%%

We begin with the negative determinant estimate used
in the bilinear small-ball argument. The proof
includes the extremization over all admissible
coefficient vectors.

\begin{lemma}[Capped simplex extremizer]
  \label{lem:det-capped-extremizer}
  Let $N\ge1$, let $0<a<1$, and assume that
  \[
    \mathcal W_{N,a}=\left\{w\in[0,a]^N:\sum_{i=1}^Nw_i=1\right\}
  \]
  is nonempty. Put
  \[
    k=\lfloor a^{-1}\rfloor
  \]
  and
  \[
    \rho=1-ka.
  \]
  Thus $0\le\rho<a$. After adding or deleting zero
  coordinates, define
  \[
    w^*=(a,\ldots,a,\rho,0,\ldots,0),
  \]
  where $a$ occurs $k$ times and the coordinate $\rho$
  is omitted when it is zero. Feasibility of
  $\mathcal W_{N,a}$ ensures that this vector has at
  most $N$ nonzero coordinates. Let
  $g_1,\ldots,g_N$ be independent $N(0,I_r)$ vectors.
  For every $s>0$,
  \[
    \E\det\left(\sum_{i=1}^Nw_ig_ig_i^T\right)^{-s} \le
    \E\det\left(\sum_{i=1}^Nw_i^*g_ig_i^T\right)^{-s}
  \]
  for every $w\in\mathcal W_{N,a}$, where both sides are
  understood as extended nonnegative expectations.
\end{lemma}

\begin{proof}
  For $\varepsilon>0$, define
  \[
    \Phi_\varepsilon(w)= \E\det\left(\varepsilon
    I_r+\sum_{i=1}^Nw_ig_ig_i^T\right)^{-s}.
  \]
  The map
  \[
    A\longmapsto\det(A)^{-s}=\exp[-s\log\det A]
  \]
  is convex on the positive definite cone. Indeed,
  $-\log\det A$ is convex, and the exponential is convex
  and increasing. Hence $\Phi_\varepsilon$ is convex in
  $w$. It is also invariant under coordinate
  permutations because the $g_i$ are identically
  distributed.

  We next verify the required majorization directly.
  Arrange the coordinates of $w$ in decreasing order.
  For every $1\le j\le N$,
  \[
    \sum_{i=1}^jw_i\le\min\{ja,1\}.
  \]
  The right side is the sum of the first $j$
  coordinates of $w^*$. Thus $w^*$ majorizes $w$. By the
  finite-dimensional majorization theorem, there is a
  doubly stochastic matrix $D$ such that $w=Dw^*$. By the
  Birkhoff theorem, $D$ is a convex combination of
  permutation matrices. Convexity and permutation
  invariance therefore give
  \[
    \Phi_\varepsilon(w)\le\Phi_\varepsilon(w^*).
  \]
  As $\varepsilon$ decreases to zero, both integrands
  increase pointwise. Monotone convergence proves the
  asserted extended-valued inequality.
\end{proof}

\begin{lemma}[Real Wishart negative moment]
  \label{lem:det-wishart-moment}
  Let $G$ be an $r$ by $k$ real matrix with independent
  $N(0,1)$ entries and let $W=GG^T$. If $s>0$ and
  \[
    k>r+2s-1,
  \]
  then
  \[
    \E\det(W)^{-s} =2^{-rs}\prod_{j=1}^r
    \frac{\Gamma((k-j+1)/2-s)}{\Gamma((k-j+1)/2)}.
  \]
  If $k\le r+2s-1$, the negative moment is infinite.
\end{lemma}

\begin{proof}
  If $k<r$, then $\rank(W)=\rank(G)\leq k<r$, so
  $\det(W)=0$ almost surely and $\E\det(W)^{-s}=+\infty$
  for every $s>0$. Assume henceforth $k\geq r$.

  Apply Gram-Schmidt to the rows of $G$. For $1\leq
  j\leq r\leq k$, the squared length of the component of
  row $j$ orthogonal to the preceding $j-1$ rows has a
  chi-square law with $k-j+1$ degrees of freedom.
  Gaussian orthogonal invariance also makes these squared
  lengths independent. Consequently,
  \[
    \det(W)\stackrel{d}=\prod_{j=1}^rX_j,
  \]
  where the $X_j$ are independent and $X_j$ has a
  chi-square law with $k-j+1$ degrees of freedom.

  If $X$ has a chi-square law with $m$ degrees of
  freedom, direct integration of its density gives
  \[
    \E X^{-s} =\frac1{2^{m/2}\Gamma(m/2)}
    \int_0^\infty x^{m/2-s-1}e^{-x/2}dx
    =2^{-s}\frac{\Gamma(m/2-s)}{\Gamma(m/2)}.
  \]
  This integral is finite exactly when $m>2s$.
  Multiplying the $r$ independent moments proves the
  formula and the finiteness condition.
\end{proof}

\begin{proposition}[Uniform capped determinant moment]
  \label{prop:det-capped-moment}
  Fix $0<\theta<1/\sqrt2$. There is $C_\theta<\infty$
  such that the following holds for every $r,N\ge1$. Let
  $g_1,\ldots,g_N$ be independent $N(0,I_r)$ vectors and
  let $z\in\R^N$ satisfy
  \[
    \|z\|_2=1
  \]
  and
  \[
    \sqrt r\|z\|_\infty\le\theta.
  \]
  Then, with
  \[
    W_z=\sum_{i=1}^Nz_i^2g_ig_i^T,
  \]
  one has
  \[
    \E\det(W_z)^{-r/2}\le C_\theta^{r^2}.
  \]
  One may take
  \[
    C_\theta=(1-2\theta^2)^{-1/2}
  \]
  for $r\ge2$, after increasing the constant to cover
  $r=1$.
\end{proposition}

\begin{proof}
  Set
  \[
    a=\frac{\theta^2}{r},
    k=\left\lfloor\frac r{\theta^2}\right\rfloor,
    \rho=1-ka.
  \]
  The weights $w_i=z_i^2$ belong to $\mathcal W_{N,a}$.
  By \cref{lem:det-capped-extremizer}, their negative
  determinant moment is at most the moment at the capped
  vector $w^*$. For independent Gaussian vectors
  $g_1,\ldots,g_{k+1}$,
  \[
    \sum_iw_i^*g_ig_i^T =a\sum_{i=1}^kg_ig_i^T+\rho
    g_{k+1}g_{k+1}^T \succeq a\sum_{i=1}^kg_ig_i^T.
  \]
  When $\rho=0$, the last rank-one term is absent.
  Monotonicity of the determinant on the positive
  definite cone gives
  \[
    \E\det(W_z)^{-r/2} \le
    a^{-r^2/2}\E\det(W_{r,k})^{-r/2},
  \]
  where $W_{r,k}$ is a real Wishart matrix with $k$
  degrees of freedom.

  Put
  \[
    \delta_\theta=\theta^{-2}-2>0.
  \]
  The floor in the definition of $k$ gives
  \[
    k\ge\frac r{\theta^2}-1
  \]
  and therefore
  \[
    k-2r+1\ge\delta_\theta r>0.
  \]
  In particular, $k>2r-1$, so
  \cref{lem:det-wishart-moment} applies with $s=r/2$. It
  gives
  \[
    \E\det(W_{r,k})^{-r/2}
    =2^{-r^2/2}\prod_{j=1}^r
    \frac{\Gamma(y_j)}{\Gamma(y_j+r/2)},
  \]
  where
  \[
    y_j=\frac{k-j+1-r}{2}.
  \]
  The same floor estimate yields
  \[
    y_j\ge\frac{k-2r+1}{2}\ge\frac{\delta_\theta r}{2}.
  \]

  Assume first that $r\ge2$. If $X$ has the gamma
  distribution with shape $y_j$ and unit scale, then
  \[
    \frac{\Gamma(y_j+r/2)}{\Gamma(y_j)}=\E X^{r/2}.
  \]
  Since $r/2\ge1$, Jensen's inequality gives
  \[
    \E X^{r/2}\ge(\E X)^{r/2}=y_j^{r/2}.
  \]
  Thus
  \[
    \frac{\Gamma(y_j)}{\Gamma(y_j+r/2)} \le
    y_j^{-r/2} \le \left(\frac2{\delta_\theta
    r}\right)^{r/2}.
  \]
  Multiplication over $j$ now gives
  \[
    \E\det(W_z)^{-r/2} \le
    \left(\frac1{2a}\right)^{r^2/2}
    \left(\frac2{\delta_\theta r}\right)^{r^2/2} =
    \left(\frac1{\theta^2\delta_\theta}\right)^{r^2/2} =
    (1-2\theta^2)^{-r^2/2}.
  \]
  This calculation includes the floor in $k$ and does not
  use an asymptotic replacement of $ak$ by one.

  For $r=1$, the exact expression from
  \cref{lem:det-wishart-moment} is finite because
  $k=\lfloor\theta^{-2}\rfloor\ge2$. It is a finite number
  depending only on $\theta$. Increasing $C_\theta$ to
  cover this one case completes the proof.
\end{proof}


%%%%%%%%%%%%%%%%%%%%
\subsection{Arbitrary-center Frobenius small balls}
%%%%%%%%%%%%%%%%%%%%

\begin{proposition}[Determinant small ball]
  \label{prop:det-frobenius-small-ball}
  Fix $0<\theta<1/\sqrt2$. There is $C_\theta<\infty$ such
  that the following holds. Let $g_i,k_i$ be mutually
  independent $N(0,I_r)$ vectors, let $z\in\R^N$ satisfy
  \[
    \|z\|_2=1
  \]
  and
  \[
    \sqrt r\|z\|_\infty\le\theta,
  \]
  and put
  \[
    Y_z=\frac1{\sqrt r}\sum_{i=1}^Nz_ig_ik_i^T.
  \]
  For every deterministic $M\in\R^{r\times r}$ and every
  $t>0$,
  \[
    \Pp(\|Y_z-M\|_F\le t\sqrt r) \le (C_\theta t)^{r^2}.
  \]
  The same bound holds conditionally when $M$ and $z$ are
  measurable with respect to a sigma-field independent of
  all $g_i,k_i$, provided the two displayed conditions on
  $z$ hold almost surely.
\end{proposition}

\begin{proof}
  Condition on $g_1,\ldots,g_N$ and write
  \[
    W_z=\sum_{i=1}^Nz_i^2g_ig_i^T.
  \]
  The columns of $Y_z$ are conditionally independent
  Gaussian vectors with covariance $W_z/r$. On the event
  that $W_z$ is positive definite, their joint density is
  bounded by
  \[
    (2\pi)^{-r^2/2}\det(W_z/r)^{-r/2}.
  \]
  The moment estimate in \cref{prop:det-capped-moment}
  implies that $W_z$ is positive definite almost surely:
  positive probability of $\det W_z=0$ would make
  $\E\det(W_z)^{-r/2}$ infinite.

  The Euclidean volume of a ball of radius $R$ in
  dimension $m$ is at most
  \[
    \left(\frac{C R}{\sqrt m}\right)^m.
  \]
  Apply this with $m=r^2$ and $R=t\sqrt r$. The
  conditional density bound and the volume estimate give,
  uniformly in $M$,
  \[
    \Pp(\|Y_z-M\|_F\le t\sqrt r\mid g_1,\ldots,g_N)
    \le(Ct)^{r^2}\det(W_z)^{-r/2}.
  \]
  Taking expectation and applying
  \cref{prop:det-capped-moment} proves the claim. For the
  conditional statement, condition first on the
  sigma-field that determines $M$ and $z$ and repeat the
  same argument pointwise.
\end{proof}

\begin{corollary}[Scaled Frobenius small ball]
  \label{cor:det-scaled-frobenius}
  Fix $0<\theta<1/\sqrt2$. Let $u\ne0$, let
  $\sigma=\|u\|_2$, and assume
  \[
    \sqrt r\frac{\|u\|_\infty}{\sigma}\le\theta.
  \]
  For mutually independent $N(0,I_r)$ vectors $g_i,k_i$,
  define
  \[
    M_u=\frac1{\sqrt d}\sum_i u_ig_ik_i^T.
  \]
  Then, uniformly over independent centers $C$ and every
  $t>0$,
  \[
    \Pp\left(\|M_u-C\|_F\le t\sigma\frac r{\sqrt d}\right)
    \le (C_\theta t)^{r^2}.
  \]
  The same statement holds conditionally when $u$ and $C$
  are measurable with respect to a sigma-field independent
  of the displayed Gaussian vectors.
\end{corollary}

\begin{proof}
  Apply \cref{prop:det-frobenius-small-ball} to
  $z=u/\sigma$ and use
  \[
    M_u=\sigma\sqrt{\frac rd}Y_z.
  \]
\end{proof}

%%%%%%%%%%%%%%%%%%%%
\subsection{Nuclear-ball covering and the block estimate}
%%%%%%%%%%%%%%%%%%%%

Let $B_F$ denote the unit ball in Frobenius norm. For a
bounded set $T$, let $N(T,\varepsilon B_F)$ be the least
number of translates of $\varepsilon B_F$ required to
cover $T$.

\begin{lemma}[Nuclear-ball entropy]\label{lem:det-nuclear-cover}
  Let
  \[
    K_{r,R}=\{A\in\R^{r\times r}:\|A\|_*\le R\}.
  \]
  There is a numerical constant $C$ such that, for every
  $\varepsilon>0$,
  \[
    \log N(K_{r,R},\varepsilon B_F) \le
    C\frac{R^2r}{\varepsilon^2}.
  \]
  In particular, for $R=\nu r$ and
  $\varepsilon=t\sqrt r$,
  \[
    \log N(K_{r,\nu r},t\sqrt r B_F) \le C\frac{\nu^2}{t^2}r^2.
  \]
\end{lemma}

\begin{proof}
  Let $G$ be an $r$ by $r$ matrix with independent
  $N(0,1)$ entries. Nuclear and operator norm duality
  gives
  \[
    \E\sup_{A\in K_{r,R}}\langle G,A\rangle
    =R\E\|G\|_{op}.
  \]
  The Gaussian comparison inequality applied to the
  process $y^TGx$, with $x,y$ on the Euclidean unit
  sphere, gives
  \[
    \E\|G\|_{op}\le2\sqrt r.
  \]
  Thus the Gaussian width of $K_{r,R}$ is at most
  $2R\sqrt r$.

  Choose a maximal $\varepsilon$-separated subset of
  $K_{r,R}$ in Frobenius norm. It is an
  $\varepsilon$-net. Sudakov's minoration in the
  finite-set form \cite[Section 2.1]{Fernique1974} says
  that an $\varepsilon$-separated set of cardinality $M$
  in the canonical metric of a Gaussian process has
  expected supremum at least
  $c\varepsilon\sqrt{\log M}$. Applied here, it gives
  \[
    \varepsilon\sqrt{\log N(K_{r,R},\varepsilon B_F)}
    \le C\E\sup_{A\in K_{r,R}}\langle G,A\rangle.
  \]
  Squaring and using the width bound proves
  \[
    \log N(K_{r,R},\varepsilon B_F) \le
    C\frac{R^2r}{\varepsilon^2}.
  \]
  The stated specialization follows by substitution.
\end{proof}

\begin{theorem}[Arbitrary-center nuclear small ball]
  \label{thm:det-nuclear-small-ball}
  Fix $0<\theta<1/\sqrt2$. There are constants
  $c_\theta,\nu_\theta>0$ such that the following holds
  under the hypotheses of
  \cref{prop:det-frobenius-small-ball}. For every
  deterministic $M\in\R^{r\times r}$,
  \[
    \Pp(\|Y_z-M\|_*\le\nu_\theta r) \le e^{-c_\theta r^2}.
  \]
  The same conclusion holds conditionally when $M$ and $z$
  are measurable with respect to a sigma-field independent
  of the driver vectors and the stated conditions on $z$
  hold almost surely.
\end{theorem}

\begin{proof}
  Let $C_\theta$ be the constant in
  Lemma~\ref{prop:det-frobenius-small-ball}. Choose $t_\theta>0$
  so small that
  \[
    \log(C_\theta t_\theta)\le-2.
  \]
  By Lemma~\ref{lem:det-nuclear-cover}, the nuclear ball
  $K_{r,\nu r}$ has a Frobenius $t_\theta\sqrt r$-net of
  cardinality at most
  \[
    \exp\left(C\frac{\nu^2}{t_\theta^2}r^2\right).
  \]
  If $\|Y_z-M\|_*\le\nu r$, then for one center $A$ in
  this net,
  \[
    \|Y_z-(M+A)\|_F\le t_\theta\sqrt r.
  \]
  The union bound and
  \cref{prop:det-frobenius-small-ball} therefore give
  \[
    \Pp(\|Y_z-M\|_*\le\nu r) \le \exp\left[
    \left(C\frac{\nu^2}{t_\theta^2}
    +\log(C_\theta t_\theta)\right)r^2 \right].
  \]
  Choose $\nu=\nu_\theta>0$ so small that
  \[
    C\frac{\nu_\theta^2}{t_\theta^2}\le1.
  \]
  The exponent is then at most $-r^2$, and the theorem
  follows after changing the constant. The conditional
  statement follows because both the Frobenius small-ball
  estimate and the covering number are uniform in the
  center.
\end{proof}

\begin{corollary}[Scaled coordinate block]
  \label{cor:det-scaled-block}
  Fix $0<\theta<1/\sqrt2$. Let $u\ne0$ be deterministic,
  let $\sigma=\|u\|_2$, and assume
  \[
    \sqrt r\frac{\|u\|_\infty}{\sigma}\le\theta.
  \]
  For mutually independent $N(0,I_r)$ vectors $g_i,k_i$,
  define
  \[
    M_u=\frac1{\sqrt d}\sum_i u_ig_ik_i^T.
  \]
  If $C$ is deterministic, or is measurable with respect
  to a sigma-field independent of the displayed Gaussian
  vectors, then
  \[
    \Pp\left( \|M_u-C\|_*\le
    \nu_\theta\sigma\frac{r^{3/2}}{\sqrt d} \right) \le
    e^{-c_\theta r^2}.
  \]
\end{corollary}

\begin{proof}
  Apply \cref{thm:det-nuclear-small-ball} to
  $z=u/\sigma$ and use
  \[
    M_u=\sigma\sqrt{\frac rd}Y_z.
  \]
  Rescale the center by the same deterministic factor.
\end{proof}    
\end{aiblock}



%%%%%%%%%%%%%%%%%%%%
\section{A localized uniform estimate for small
  coefficients}\label{expert-localized-module}
%%%%%%%%%%%%%%%%%%%%

\begin{aiblock}
We now pass from one deterministic vector
to uniform over small-coordinate vectors and
admissible centers. With common radius
$\zeta$, each $V_0$ has small-ball
$\lesssim(C_\theta\zeta)^{h^2}$ while
$|\mathcal N_V|\lesssim(C/\zeta)^n$; the
union is
$\exp[-(h^2-n)\log(1/\zeta)+O(d^2)]$. Hence
we need the strict gap $h^2-n\ge cd^2$ from
$\alpha_0<1/4$ to make fixed small $\zeta$
beat entropy.

% BEGIN INLINED FILE: localized_tube_module.tex
%%%%%%%%%%%%%%%%%%%%
\subsection{Uniform small-ball estimate after a
  large/small coordinate split}
%%%%%%%%%%%%%%%%%%%%

The proposition below says the following. Fix a
support $I$ of at most $Kd$ large coefficients.
The contribution from $I$ may be arbitrary and
is treated as a random center after freezing
$\mathcal G_I=\sigma(a_i,b_i:i\in I)$: then
$\mathcal B U_0$ and $N_0-\mathcal B U_0$ are
$\mathcal G_I$-measurable while the $I^c$
pairs stay independent $N(0,I_h)$, so
$\|\mathcal B V_0-(N_0-\mathcal B U_0)\|_F
\le3\zeta\sqrt d$ has arbitrary-center
small-ball $(C_\theta\zeta)^{h^2}$. The
outside coefficients have comparable $\ell_2$
scale and no single coordinate too large.
Uniformly over all such choices the sum
cannot have small nuclear norm. The word
``localized'' refers only to the fixed upper
bounds on the two classes.


Let $h=\lfloor d/2\rfloor$ and
\[
  \cB y=\frac1{\sqrt d}\sum_{i=1}^ny_ia_ib_i^T,
\]
where all $a_i,b_i$ are independent $N(0,I_h)$ vectors.
For $L>0$, define
\[
  E_d(L)=\{\|\cB\|_{2\to F}\le L\sqrt d\}.
\]

\begin{proposition}[Uniform conditional small-ball
  estimate]
  \label{main-proposition-5-1}
  Fix $\alpha_0<1/4$, $0<\theta<1/\sqrt2$, and positive
  constants $A,B,\beta,K,L$. There are constants
  $c,\gamma>0$, depending only on these fixed parameters,
  such that the following holds whenever
  $n\le\alpha_0d^2$.

  Let $F_{d,L}$ be the event that there are $q>0$,
  $I\subset[n]$, and $U,V\in\R^n$ satisfying
  \[
    |I|\le Kd,
  \]
  \[
    \supp U\subset I,
  \]
  \[
    \supp V\subset I^c,
  \]
  \[
    \|U\|_1\le Aq\sqrt d,
  \]
  \[
    \beta q\le\|V\|_2\le Bq,
  \]
  \[
    \sqrt d\|V\|_\infty\le\theta\|V\|_2,
  \]
  but
  \[
    \|\cB(U+V)\|_*<cqd.
  \]
  Then, for all sufficiently large $d$,
  \[
    \Pp(F_{d,L}\cap E_d(L))\le e^{-\gamma d^2}.
  \]
\end{proposition}

\begin{proof}
  The assumptions and conclusion are homogeneous in
  $(q,U,V)$. Dividing $U$ and $V$ by $q$ reduces the
  proof exactly to $q=1$. Thus no discretization of the
  scale $q$ is needed.

  Fix $I\subset[n]$ with $s=|I|\le Kd$, and define
  \[
    \mathcal U_I=\{U:\supp U\subset I,\|U\|_1\le
    A\sqrt d\}
  \]
  and
  \[
    \mathcal V_I=\{V:\supp V\subset I^c,\beta\le
    \|V\|_2\le B,
  \]
  \[
    \sqrt d\|V\|_\infty\le\theta\|V\|_2\}.
  \]
  Choose
  \[
    0<\zeta<\min\{1,\beta/6\}.
  \]

  Let $\mathcal N_U(I)$ be an $\ell_1$ net of
  $\mathcal U_I$ at radius $\zeta/L$. A volumetric
  estimate in the $s$-dimensional normed space gives
  \[
    |\mathcal N_U(I)| \le
    \left(1+\frac{2AL\sqrt d}{\zeta}\right)^s.
  \]
  Let $\mathcal N_V(I)$ be a maximal internal Euclidean
  net of $\mathcal V_I$ at radius $\zeta/L$. Since
  $\mathcal V_I$ lies in the Euclidean ball of radius
  $B$ in dimension $n-s$, the packing estimate gives
  \[
    |\mathcal N_V(I)| \le
    \left(1+\frac{2BL}{\zeta}\right)^{n-s}.
  \]
  Every center $V_0\in\mathcal N_V(I)$ belongs to
  $\mathcal V_I$. In particular,
  \[
    \beta\le\|V_0\|_2\le B
  \]
  and
  \[
    \sqrt d\|V_0\|_\infty\le\theta\|V_0\|_2.
  \]

  We next construct a deterministic matrix-center net.
  Put
  \[
    \mathcal K_c=\{N\in\R^{h\times h}:\|N\|_*\le cd\}.
  \]
  Applying the nuclear-ball covering estimate
  Lemma~\ref{lem:det-nuclear-cover} with $r=h$, $R=cd$, and
  $\varepsilon=\zeta\sqrt d$ gives a deterministic
  Frobenius net $\mathcal N_*(c)$ of $\mathcal K_c$ at
  that radius such that
  \[
    \log|\mathcal N_*(c)| \le
    C\frac{c^2}{\zeta^2}d^2.
  \]
  All three nets are fixed before the Gaussian samples
  are observed.

  Suppose that a bad pair $(U,V)$ exists on $E_d(L)$.
  Choose $U_0\in\mathcal N_U(I)$ and
  $V_0\in\mathcal N_V(I)$ such that
  \[
    \|U-U_0\|_1\le\zeta/L
  \]
  and
  \[
    \|V-V_0\|_2\le\zeta/L.
  \]
  On $E_d(L)$,
  \[
    \|\cB(U-U_0)\|_F \le L\sqrt d\|U-U_0\|_2
    \le\zeta\sqrt d
  \]
  and
  \[
    \|\cB(V-V_0)\|_F\le\zeta\sqrt d.
  \]
  Since $\|\cB(U+V)\|_*<cd$, there is
  $N_0\in\mathcal N_*(c)$ such that
  \[
    \|\cB(U+V)-N_0\|_F\le\zeta\sqrt d.
  \]
  It follows that
  \[
    \|\cB V_0-(N_0-\cB U_0)\|_F\le3\zeta\sqrt d.
  \]

  Fix a deterministic tuple $(I,U_0,V_0,N_0)$ and let
  \[
    \mathcal G_I=\sigma(a_i,b_i:i\in I).
  \]
  The matrix $N_0-\cB U_0$ is $\mathcal G_I$ measurable.
  All pairs $(a_i,b_i)$ with $i\in I^c$ are independent
  of $\mathcal G_I$. Put
  \[
    \sigma_0=\|V_0\|_2
  \]
  and
  \[
    z=V_0/\sigma_0.
  \]
  Then
  \[
    \|z\|_2=1
  \]
  and
  \[
    \sqrt d\|z\|_\infty\le\theta.
  \]

  We record the normalization from
  \cref{prop:det-frobenius-small-ball}. Since $h/d\to1/2$,
  the proposition, with $r=h$, implies after rescaling
  that there is a constant $C_\theta<\infty$ such that
  \[
    \Pp\left( \left\|\frac1{\sqrt d}\sum_i
    z_ia_ib_i^T-M\right\|_F \le t\sqrt d \right)
    \le(C_\theta t)^{h^2}
  \]
  for every deterministic $M\in\R^{h\times h}$ and every
  $t>0$. Indeed,
  \[
    \frac1{\sqrt d}\sum_i z_ia_ib_i^T
    =\sqrt{\frac hd}Y_z,
  \]
  the coefficient cap for $Y_z$ is
  \[
    \sqrt h\|z\|_\infty\le\theta\sqrt{h/d}<\theta,
  \]
  and a Frobenius radius $t\sqrt d$ becomes a radius
  \[
    t\frac d h\sqrt h\le3t\sqrt h
  \]
  for all sufficiently large $d$. The constant factor is
  absorbed into $C_\theta$. The same estimate holds
  conditionally for a center measurable with respect to a
  sigma-field independent of the driver vectors.

  After division by $\sigma_0$, the fixed-tuple bad event
  implies
  \[
    \left\| \frac1{\sqrt d}\sum_{i\in I^c}z_ia_ib_i^T
    -\frac{N_0-\cB U_0}{\sigma_0} \right\|_F \le
    t_0\sqrt d,
  \]
  where
  \[
    t_0=\frac{3\zeta}{\sigma_0}
    \le\frac{3\zeta}{\beta}<1.
  \]
  The conditional arbitrary-center small-ball estimate
  gives, almost surely in $\mathcal G_I$,
  \[
    \Pp(\text{the fixed tuple is bad}\mid\mathcal G_I)
    \le \left(\frac{3C_\theta\zeta}{\beta}\right)^{h^2}.
  \]
  The bound is uniform in the realized center.
  Integration over $\mathcal G_I$ gives the same
  unconditional estimate.

  It remains to pay the deterministic entropy. The
  support count satisfies
  \[
    \log N_{supp}=O_K(d\log d).
  \]
  To see this, if $Kd\le n/2$, then
  \[
    \log\sum_{s\le Kd}\binom ns \le
    Kd\log\left(\frac{en}{Kd}\right)+O_K(d).
  \]
  If $Kd>n/2$, then $n<2Kd$ and the bound
  $2^n=e^{O_K(d)}$ applies. The head entropy satisfies
  \[
    \log N_{head} \le
    Kd\log\left(1+\frac{2AL\sqrt d}{\zeta}\right)
    =O_{A,K,L,\zeta}(d\log d).
  \]
  The tail and matrix-center entropies satisfy
  \[
    \log N_{tail} \le
    n\log\left(1+\frac{2BL}{\zeta}\right)
  \]
  and
  \[
    \log N_{center} \le C\frac{c^2}{\zeta^2}d^2.
  \]
  Put
  \[
    \Xi_d=h^2\log\left(\frac{3C_\theta\zeta}{\beta}\right)
    +n\log\left(1+\frac{2BL}{\zeta}\right).
  \]
  The union bound gives
  \[
    \log\Pp(F_{d,L}\cap E_d(L)) \le\Xi_d
    +C\frac{c^2}{\zeta^2}d^2
    +O_{A,K,L,\zeta}(d\log d).
  \]

  Set
  \[
    g=\frac14-\alpha_0>0.
  \]
  For all sufficiently large $d$,
  \[
    h^2-n\ge\frac g2d^2.
  \]
  There is a constant $C_0<\infty$, depending only on
  $\theta,\beta,B,L$, such that
  \[
    \Xi_d\le-(h^2-n)\log(1/\zeta)+C_0d^2.
  \]
  First choose the fixed $\zeta$ so small that
  \[
    \frac g2\log(1/\zeta)\ge4C_0+4.
  \]
  With this $\zeta$ fixed, choose $c>0$ so small that
  \[
    C\frac{c^2}{\zeta^2} \le \frac g8\log(1/\zeta).
  \]
  Finally take $d$ large enough that
  \[
    O_{A,K,L,\zeta}(d\log d) \le
    \frac g8d^2\log(1/\zeta).
  \]
  The resulting exponent is at most $-\gamma d^2$ for a
  fixed $\gamma>0$. This proves the proposition.
\end{proof}    
\end{aiblock}


% END INLINED FILE: localized_tube_module.tex

%%%%%%%%%%%%%%%%%%%%
\section{A multiscale estimate for the large
  coordinates}\label{expert-multiscale-module}
%%%%%%%%%%%%%%%%%%%%

\begin{aiblock}

The small-coordinate estimate does not control a vector
concentrated on a few coordinates. The following section
treats those vectors by sorting the coefficients into
dyadic magnitude classes. We retain a later class only
when its support is geometrically larger than the
previously retained support. This is only a deterministic
bookkeeping rule. The section proves a simultaneous lower
bound for every large-coordinate vector, including vectors
with arbitrary dynamic range.

% BEGIN INLINED FILE: multiscale_coefficients_module.tex
%%%%%%%%%%%%%%%%%%%%
\subsection{Statement and
  notation}\label{statement-and-notation}
%%%%%%%%%%%%%%%%%%%%

This section proves \cref{the-clipped-head-theorem}. We
first state the estimate in the local notation used
throughout the proof; the proof is completed at the end of
this section in Section \ref{whole-class-deterministic-union}.
The result is uniform over all coefficient vectors with at
most $Kd$ nonzero coordinates, including vectors whose
support, signs, and magnitudes are selected after the
Gaussian sample is observed.

Let $h=\lfloor d/2\rfloor$. Let
\[
  a_1,\ldots,a_n,b_1,\ldots,b_n
\]
be mutually independent $N(0,I_h)$ vectors, and define
\[
  \mathcal{B}u=\frac1{\sqrt d}\sum_{i=1}^n u_i a_i
  b_i^T.
\]
For $I\subset[n]$, let $A_I$ and $B_I$ be the $h$ by
$|I|$ matrices whose columns are respectively
$(a_i)_{i\in I}$ and $(b_i)_{i\in I}$, in increasing
index order.

\begin{theorem}\label{main-theorem-6-1}
  For every fixed $0<K,A<\infty$, there is $c_{K,A}>0$
  such that, uniformly over every integer sequence
  $n=n(d)\le Ad^2$,
  \[
    \Pp\left(\inf_{\substack{u\ne0\\
    |\operatorname{supp}u|\le
    Kd}}\frac{\|\mathcal{B}u\|_*}{\sqrt d\|u\|_1}\ge
    c_{K,A}\right)\longrightarrow1.
  \]
  If $K_+=\max\{K,1\}$, then there is a constant
  $c_{\rm ms}(A)>0$, depending only on $A$, such that
  \[
    c_{K,A}\ge\frac{c_{\rm ms}(A)}{\sqrt{K_+}(1+\log
    K_+)}.
  \]
  The event is simultaneous over supports and
  coefficients chosen after all Gaussian samples are
  observed. No sign, flatness, centering, or
  dynamic-range condition is imposed on $u$.
\end{theorem}

%%%%%%%%%%%%%%%%%%%%
\subsection{The arbitrary-center block
  estimate}\label{the-arbitrary-center-block-estimate}
%%%%%%%%%%%%%%%%%%%%

\begin{lemma}\label{record-lemma-2-1}
  Let $r\ge1$ and $m\ge1$ be integers and let
  $z\in{\mathbb R}^{m}$ be deterministic. Fix
  $\theta<1/\sqrt2$. There are $c_{\rm sb}$,
  $\nu_{\rm sb}>0$ with the following property. Let
  $g_1,\ldots,g_m,k_1,\ldots,k_m$ be mutually
  independent $N(0,I_r)$ vectors and suppose
  \[
    \|z\|_2=1,\ \sqrt r\|z\|_\infty\le\theta.
  \]
  Set
  \[
    Y_z=\frac1{\sqrt r}\sum_{i=1}^{m} z_i g_i k_i^T.
  \]
  For every deterministic $C\in{\mathbb R}^{r\times r}$,
  \[
    \Pp(\|Y_z-C\|_*\le\nu_{\rm sb}r)\le
    e^{-c_{\rm sb}r^2}.
  \]
  The same statement holds conditionally if $C$ is
  measurable with respect to a sigma-field independent
  of all $g_i,k_i$.
\end{lemma}

\begin{proof}
  This is the arbitrary-center nuclear small-ball
  estimate from \cref{thm:det-nuclear-small-ball},
  applied in dimension $r$. Its conditional form follows
  because that theorem is uniform in every center
  measurable with respect to an independent sigma-field.
\end{proof}

\begin{lemma}[Rectangular
  pinching]\label{record-lemma-2-2}
  Let $L_1,\ldots,L_J$ and $R_1,\ldots,R_J$ be pairwise
  orthogonal row and column projections. Then, for every
  matrix $M$,
  \[
    \|M\|_*\ge\sum_{j=1}^J\|L_jMR_j\|_*.
  \]
\end{lemma}

\begin{proof}
  For a matrix $A$, define its partial polar factor by
  \[
    \operatorname{pol}(A)=A(A^TA)^{\dagger/2},
  \]
  where the pseudoinverse power is zero on the kernel. It
  satisfies $\|\operatorname{pol}(A)\|_{op}\le1$ and
  $\langle\operatorname{pol}(A),A\rangle=\|A\|_*$. Let
  $W_j=\operatorname{pol}(L_jMR_j)$, extended by zero
  outside the indicated row and column subspaces, so that
  $W_j=L_jW_jR_j$. Thus $\|W_j\|_{op}\le1$ and
  \[
    \langle W_j,L_jMR_j\rangle=\|L_jMR_j\|_*.
  \]
  Put $W=\sum_jW_j$. Orthogonality gives
  \[
    W^TW=\sum_jR_jW_j^TL_jW_jR_j\preceq\sum_jR_j\preceq
    I.
  \]
  Hence $\|W\|_{op}\le1$. Nuclear and operator norm
  duality gives
  \[
    \|M\|_*\ge\langle W,M\rangle=\sum_j\|L_jMR_j\|_*.
  \]
  This proves the inequality.
\end{proof}


%%%%%%%%%%%%%%%%%%%%
\subsection{Two elementary events that hold for
  every small support}\label{two-uniform-elementary-events}
%%%%%%%%%%%%%%%%%%%%

\begin{lemma}[Uniform bound for individual rank-one
  terms]\label{record-lemma-3-1}
  There is a fixed $C_A<\infty$ such that
  \[
    \Pp(E_{atom}^c)\to0,
  \]
  where
  \[
    E_{atom}=\left\{\max_{i\le n}\frac{\|a_i\|_2\|b_i\|_2}
    {\sqrt d}\le C_A\sqrt d\right\}.
  \]
  On $E_{atom}$,
  \[
    \|\mathcal{B}z\|_*\le C_A\sqrt d\|z\|_1
  \]
  for every $z\in{\mathbb R}^n$ simultaneously.
\end{lemma}

\begin{proof}
  Gaussian norm concentration gives
  \[
    \Pp(\|a_i\|_2>C\sqrt d)+\Pp(\|b_i\|_2>C\sqrt
    d)\le4e^{-cd}.
  \]
  Since $n\le Ad^2$, a union bound proves
  $\Pp(E_{atom}^c)\le4Ad^2e^{-cd}\to0$. On this event,
  \[
    \|\mathcal{B}z\|_* \le \frac1{\sqrt d}\sum_i|z_i|
    \|a_i b_i^T\|_* = \frac1{\sqrt d}\sum_i|z_i|
    \|a_i\|_2\|b_i\|_2 \le C_A\sqrt d\|z\|_1.
  \]
\end{proof}

\begin{lemma}[Uniform singular-value bound on small
  supports]\label{record-lemma-3-2}
  There are absolute $c_s,\kappa>0$ such that, with
  \[
    s_{\rm sp}=\left\lfloor\kappa\frac{d}{\log d}
    \right\rfloor,
  \]
  the event
  \[
    E_{\rm sp}=\left\{\min\{\sigma_{\min}(A_I),
    \sigma_{\min}(B_I)\}\ge c_s\sqrt d\text{ for
    all }I\subset[n]\text{ with }1\le|I|\le
    s_{\rm sp}\right\}
  \]
  satisfies $\Pp(E_{\rm sp}^c)\le e^{-c_sd}$ for all
  large $d$.
\end{lemma}

\begin{proof}
  For a fixed $I\subset[n]$ of size $1\le s\le
  s_{\rm sp}$, the Gaussian singular-value inequality
  \cite[Theorem II.13]{DavidsonSzarek2001} gives
  \[
    \Pp(\sigma_{\min}(A_I)\le\sqrt h-\sqrt s-t)
    \le e^{-t^2/2}.
  \]
  Choose $\kappa$ small enough that
  $s_{\rm sp}\le h/16$ and take $t=\sqrt h/4$. Then
  \[
    \sqrt h-\sqrt s-t\ge\frac12\sqrt h\ge c_s\sqrt d.
  \]
  Therefore, for each fixed support $I$, the
  probability that either $A_I$ or $B_I$ violates the
  preceding lower bound is at most
  \[
    2e^{-h/32}\le2e^{-c_{\rm sv}d},
  \]
  where $c_{\rm sv}=1/96$. If $n\ge2s_{\rm sp}$, then
  \[
    \log\sum_{s\le s_{\rm sp}}\binom ns\le
    s_{\rm sp}\log\left(\frac{en}{s_{\rm sp}}\right)
    +O(s_{\rm sp})\le2\kappa d
  \]
  for all large $d$, since $n\le Ad^2$. If
  $n<2s_{\rm sp}$, the number of supports is at most
  $2^n\le e^{2s_{\rm sp}\log2}$. Choose
  $\kappa<c_{\rm sv}/8$. Union over both Gaussian
  matrices and every support gives
  \[
    \Pp(E_{\rm sp}^c)\le e^{-c_sd}
  \]
  for all large $d$, after adjusting $c_s>0$.
\end{proof}

%%%%%%%%%%%%%%%%%%%%
\subsection{Dyadic decomposition with geometrically
  increasing support sizes}\label{deterministic-record-decomposition}
%%%%%%%%%%%%%%%%%%%%

The next construction is entirely deterministic. Its
purpose is to replace an arbitrary coefficient vector,
possibly with an enormous dynamic range, by finitely
many groups whose coefficients are comparable in
magnitude and whose support sizes grow geometrically. We
call a selected dyadic level a support record; this word
has no meaning beyond the selection rule in the next
paragraph.

Fix $q>1$. For $u=0$ set $V=W=0$. For $u\ne0$ put
$\lambda_{\max}=\|u\|_\infty$ and define
\[
  B_t=\{i:2^{-t-1}\lambda_{\max}<|u_i|
  \le2^{-t}\lambda_{\max}\},\quad t\ge0.
\]
Empty levels are ignored. Let $t_1$ be the first
nonempty level. Given $t_r$, let $t_{r+1}$ be the
first later nonempty level satisfying
\[
  |B_{t_{r+1}}|\ge q|B_{t_r}|,
\]
if it exists; otherwise the construction stops and the
last group contains all levels after its selected level.
Put
\[
  S_r=|B_{t_r}|,\quad
  \lambda_r=2^{-t_r-1}\lambda_{\max},
\]
and $G_r=\{t:t_r\le t<t_{r+1}\}$ with
$t_{R+1}=\infty$ for the last record $R$. Every
nonempty nonselected level $t\in G_r\setminus\{t_r\}$
satisfies $|B_t|<qS_r$.

Fix an integer $T\ge1$. In group $r$, retain only
levels whose offset $j=t-t_r$ lies in
$\{0,\dots,T-1\}$. Let $V_r$ be the restriction of $u$
to $\bigcup_{j=0}^{T-1}B_{t_r+j}\cap G_r$, let $W_r$
be the restriction to $\bigcup_{j\ge T}B_{t_r+j}\cap
G_r$, and put
\[
  V=\sum_rV_r,\quad W=u-V=\sum_rW_r,
\]
\[
  k_r=|\operatorname{supp}V_r|,\quad
  m_r=\|V_r\|_1,\quad Q_r=\|V_r\|_2^2,\quad
  R_r=\frac{m_r^2}{Q_r}.
\]
Supports of distinct $V_r$ are disjoint and $m_r>0$,
$Q_r>0$ because $B_{t_r}\subset\operatorname{supp}V_r$.

Every nonselected level in group $r$ has fewer than
$qS_r$ entries. The base level has entries in
$(\lambda_r,2\lambda_r]$, while a level at offset $j$
has entries at most $2\lambda_r2^{-j}$. Therefore
\[
  S_r\le k_r\le TqS_r,
\]
\[
  \lambda_rS_r\le m_r\le2\lambda_rS_r
  +2q\lambda_rS_r\sum_{j\ge1}2^{-j}
  \le4q\lambda_rS_r,
\]
and
\[
  \lambda_r^2S_r\le Q_r\le4\lambda_r^2S_r
  +4q\lambda_r^2S_r\sum_{j\ge1}4^{-j}
  \le\frac{16q}{3}\lambda_r^2S_r.
\]
Consequently,
\[
  \frac3{16q}S_r\le R_r\le16q^2S_r.
\]
Also,
\[
  \|V_r\|_\infty\le2\lambda_r\le8q\frac{Q_r}{m_r}
  =8q\frac{\|V_r\|_2}{\sqrt{R_r}}.
\]
The omitted mass in group $r$ is bounded by
\[
  \|W_r\|_1\le2q\lambda_rS_r\sum_{j\ge T}2^{-j}
  =4q2^{-T}\lambda_rS_r\le4q2^{-T}m_r.
\]
Summing gives
\[
  \|W\|_1\le\epsilon_T\|u\|_1,\quad
  \epsilon_T=4q2^{-T}.
\]
The selected support sizes grow as $S_{r+1}\ge qS_r$.
Hence
\[
  \sum_{l<r}S_l\le\frac{S_r}{q-1}
\]
and
\[
  \sum_{l<r}k_l\le\frac{Tq}{q-1}S_r.
\]
These are deterministic statements. Their constants
depend only on $q,T$ and are independent of coefficient
signs, the number of empty levels, and the dynamic
range.

\begin{remark}[A small example]
  Take $q=2$ and suppose the nonempty dyadic levels
  have support sizes $1,1,3,2,10$ at successive levels.
  The selected support records have sizes $1,3,10$: the
  second level is not selected because $1<2\cdot1$, the
  third is selected because $3\ge2\cdot1$, the fourth is
  not selected because $2<2\cdot3$, and the fifth is
  selected because $10\ge2\cdot3$. Thus the first two
  dyadic levels form the first group, the next two form
  the second, and the final level begins the third. The
  actual coefficient values and their signs play no role
  in this selection.
\end{remark}

%%%%%%%%%%%%%%%%%%%%
\subsection{One regular dyadic group after removing
  earlier sample
  spans}\label{projected-regular-cluster-lemma}
%%%%%%%%%%%%%%%%%%%%

\begin{lemma}\label{record-lemma-5-1}
  Fix $c_{\rm rank},C_{\rm reg}>0$. There are
  constants $C_{\rm low},c_{\rm up},c_{\rm proj}>0$,
  depending only on $c_{\rm rank},C_{\rm reg}$, for
  which the following holds. Let $H,J\subset[n]$ be
  disjoint deterministic supports. Let $P_H,Q_H$ project
  away from the left and right sample spans indexed by
  $H$. Suppose their residual ranks are at least
  $c_{\rm rank}d$.

  Let $u$ be deterministic and supported on $J$. Put
  \[
    m=\|u\|_1, \sigma=\|u\|_2,
    R=\frac{m^2}{\sigma^2}.
  \]
  Assume
  \[
    \|u\|_\infty\le C_{\rm reg}\frac\sigma{\sqrt R},
    C_{\rm low}\log d\le R\le c_{\rm up}d.
  \]
  Then, almost surely with respect to
  $\sigma(a_i,b_i:i\in H)$,
  \[
    \Pp\left(\|P_H\mathcal{B}uQ_H\|_*<c_{\rm proj}
    \sqrt d m\middle|\sigma(a_i,b_i:i\in H)\right)\le
    e^{-c_{\rm proj}dR}.
  \]
\end{lemma}

\begin{proof}
  Choose measurable orthonormal bases of the residual
  spaces. Conditional on the head sigma-field, the
  projected samples indexed by $J$ have independent
  standard Gaussian coordinates in those bases.

  Fix $\theta_0=1/2$. Choose $\gamma>0$ such that
  \[
    C_{\rm reg}\sqrt\gamma\le\theta_0<1/\sqrt2,
  \]
  and put $r=\lfloor\gamma R\rfloor$. Increasing
  $C_{\rm low}$ if needed gives $r\ge c\log d$, while
  decreasing $c_{\rm up}$ makes $r$ smaller than a fixed
  fraction of each residual rank.

  Partition each residual coordinate space into
  \[
    J_0=\left\lfloor\frac{\min\{\operatorname{rank}
    P_H,\operatorname{rank}Q_H\}}r\right\rfloor
  \]
  blocks of side $r$, and pair blocks with the same
  index. Then $J_0\ge c d/r$.

  For a paired block, write
  \[
    M_j=\frac1{\sqrt d}\sum_{i\in J}u_i
    g_{ij}k_{ij}^T = \frac{\sigma\sqrt r}{\sqrt d}Y_j,
  \]
  where
  \[
    Y_j=\frac1{\sqrt r}\sum_{i\in J}\frac{u_i}{\sigma}
    g_{ij}k_{ij}^T.
  \]
  The coefficient cap is
  \[
    \sqrt r\max_i\frac{|u_i|}{\sigma} \le
    C_{\rm reg}\sqrt{\frac rR} \le\theta_0.
  \]
  The arbitrary-center estimate
  \cref{record-lemma-2-1}, with $\theta=\theta_0$ and
  center zero, gives
  \[
    \Pp\left(\|M_j\|_*<\nu_{\rm sb}\sigma
    \frac{r^{3/2}}{\sqrt d}\middle|\sigma(a_i,b_i:i\in
    H)\right)\le p_r,
  \]
  where $p_r=e^{-c_{\rm sb}r^2}$. The events for
  different paired blocks are conditionally independent
  because they use disjoint coordinates of every current
  $a_i$ and $b_i$.

  Let $N_{bad}$ be the number of bad blocks. For all
  large $d$,
  \[
    \Pp(N_{bad}\ge J_0/2\mid\sigma(a_i,b_i:i\in H))
    \le \sum_{s\ge J_0/2}\binom{J_0}{s}p_r^s \le
    2^{J_0}p_r^{J_0/2}.
  \]
  Since $p_r=e^{-c_{\rm sb}r^2}$ and $J_0r\ge cd$, the
  last expression is at most
  $e^{-cJ_0r^2}\le e^{-cdr}$.

  On the complementary event,
  \cref{record-lemma-2-2} gives
  \[
    \|P_H\mathcal{B}uQ_H\|_* \ge
    \sum_{j=1}^{J_0}\|M_j\|_* \ge \frac{J_0}{2}
    \nu_{\rm sb}\sigma\frac{r^{3/2}}{\sqrt d} \ge
    c\sigma\sqrt{dr}.
  \]
  Since $r\asymp R$ and $\sigma\sqrt R=m$, the last
  quantity is at least $c\sqrt d m$. Also $dr\asymp dR$.
  This proves the lemma.
\end{proof}

%%%%%%%%%%%%%%%%%%%%
\subsection{The small-support prefix and the point
  where span removal
  stops}\label{sparse-prefix-and-span-cutoff}
%%%%%%%%%%%%%%%%%%%%

Let $j_{\rm pre}$ be the largest group index for which
\[
  k_{\rm pre}:=\sum_{r\le j_{\rm pre}}k_{r}\le s_{sp}.
\]
If there is no such selected group, put
$j_{\rm pre}=0$, $k_{\rm pre}=0$, and $V_{0}=0$.
Otherwise, put $V_{0}=\sum_{r\le j_{\rm pre}}V_{r}$.
On $E_{sp}$,
\[
  \|\mathcal{B}V_{0}\|_{*}\ge c_{s}^{2}\sqrt{d}\|V_{0}\|_{1}.
\]
If selected group $j_{\rm pre}+1$ exists, then
\[
  S_{j_{\rm pre}+1}\ge c_{q,T}\frac{d}{\log d}.
\]
Indeed, first suppose $k_{\rm pre}>s_{sp}/2$. The
cumulative support estimate gives
$k_{\rm pre}\le C_{q,T}S_{j_{\rm pre}}$, so geometric
growth gives
$S_{j_{\rm pre}+1}\ge qS_{j_{\rm pre}}\ge
c_{q,T}s_{sp}$. If $k_{\rm pre}\le s_{sp}/2$,
maximality gives $k_{j_{\rm pre}+1}>s_{sp}/2$,
while $k_{j_{\rm pre}+1}\le TqS_{j_{\rm pre}+1}$.
The same conclusion follows. This also covers
$j_{\rm pre}=0$.

The deterministic dyadic-group estimates imply
\[
  \|V_{r}\|_{\infty}\le8q\frac{\|V_{r}\|_{2}}
  {\sqrt{R_{r}}}.
\]
Apply \cref{record-lemma-5-1} with $c_{\rm rank}=1/4$
and $C_{\rm reg}=8q$, and let $C_{\rm low},c_{\rm up}$
be the resulting effective-support thresholds. Fix
\[
  0<\rho<\min\left\{\frac18,
  \frac{c_{\rm up}}{32q^{2}}\right\}.
\]
For all sufficiently large $d$, this gives
\[
  h-\rho d\ge\frac{d}{4}.
\]
Starting after the small-support prefix, retain
selected groups only while
\[
  \sum_{l\le r}k_{l}\le\rho d.
\]
Before each retained selected group, both residual
ranks are at least $h-\rho d$. Moreover,
$S_{r}\le k_{r}\le\rho d$, so
\[
  R_{r}\le16q^{2}S_{r}\le16q^{2}\rho d<c_{\rm up}d.
\]
Every retained selected group outside the small-support
prefix also has
\[
  R_{r}\asymp_{q}S_{r}\ge c_{q,T}\frac{d}{\log d}
  \ge C_{\rm low}\log d
\]
for all sufficiently large $d$. Thus
\cref{record-lemma-5-1} applies with the displayed fixed
parameters.

If $r_{*}$ is the first selected group not treated by
span removal, then
\[
  S_{r_{*}}\ge c_{q,T,\rho}d.
\]
Indeed, the support through $r_{*}$ exceeds $\rho d$,
while
\[
  \sum_{l\le r_{*}}k_{l}\le Tq\sum_{l\le r_{*}}S_{l}\le
  \frac{Tq^{2}}{q-1}S_{r_{*}}.
\]
Since $S_{r}\le|\operatorname{supp}u|\le Kd$ and the
selected support sizes grow by $q$, the number of
remaining groups satisfies
\[
  J_{\rm late}\le1+\log_{q}
  \left(\frac{C_{q}KT}{\rho}\right).
\]


%%%%%%%%%%%%%%%%%%%%
\subsection{Combining nuclear-norm witnesses on
  orthogonal subspaces}\label{orthogonal-polar-aggregation}
%%%%%%%%%%%%%%%%%%%%

For each retained group we construct a matrix of operator
norm at most one that attains the nuclear norm of that
group's projected contribution. The left and right support
spaces for different groups are orthogonal. The witnesses
can therefore be added without increasing the operator norm
above one.

For a matrix $M$, define
\[
  \operatorname{pol}(M)=M(M^TM)^{\dagger/2},
\]
where $A^{\dagger/2}$ acts by $\lambda^{-1/2}$ on every
positive eigenspace of $A$ and by zero on its kernel. Thus
$\|\operatorname{pol}(M)\|_{op}\le1$ and
$\langle\operatorname{pol}(M),M\rangle=\|M\|_*$. Let
$L_0,R_0$ project onto the left and right sample spans of
the sparse prefix, put $M_0=\mathcal{B}V_0$, and define
$Z_0=\operatorname{pol}(M_0)$. Then $Z_0=L_0Z_0R_0$; set
$Z_0=0$ when $M_0=0$.

For a selected group $r$ treated by span removal, let
$P_{r-1},Q_{r-1}$ project away from every retained sample
span preceding $r$, and let $P_r,Q_r$ also remove the
current sample spans. Define
\[
  D_r=P_{r-1}-P_r, E_r=Q_{r-1}-Q_r.
\]
Nestedness makes $D_r,E_r$ orthogonal projections, and
different increments are pairwise orthogonal. Moreover,
\[
  P_{r-1}\mathcal{B}V_rQ_{r-1} =
  D_rP_{r-1}\mathcal{B}V_rQ_{r-1}E_r.
\]
Indeed, nestedness gives $P_rP_{r-1}=P_r$. For every
current index $i$, one also has $P_ra_i=0$. Therefore
\[
  D_rP_{r-1}a_i=(P_{r-1}-P_r)P_{r-1}a_i=P_{r-1}a_i.
\]
The same calculation on the right gives
$E_rQ_{r-1}b_i=Q_{r-1}b_i$.

Put $M_r=P_{r-1}\mathcal{B}V_rQ_{r-1}$ and
$Z_r=\operatorname{pol}(M_r)$. Then
\[
  Z_r=D_rZ_rE_r, \|Z_r\|_{op}\le1.
\]
The prefix support projections and all increments are
mutually orthogonal. Hence, for
\[
  Z=Z_0+\sum_{r\text{ captured}}Z_r,
\]
the same calculation as in \cref{record-lemma-2-2} gives
$Z^TZ\preceq I$ and therefore
\[
  \|Z\|_{op}\le1.
\]
If $k<r$, every left sample in $V_k$ lies in the span
removed before increment $D_r$. Thus
\[
  \langle Z_r,\mathcal{B}V_k\rangle=0.
\]
All earlier cross terms vanish pathwise. Orthogonal
projection and the canonical partial polar map can be
chosen as Borel functions of the samples.

%%%%%%%%%%%%%%%%%%%%
\subsection{Controlling cross terms from later
  coefficient groups}\label{future-cross-term-supermartingale}
%%%%%%%%%%%%%%%%%%%%

The witness for an earlier group also pairs with later
groups. These terms have conditional mean zero but are
not bounded almost surely. We use their exact conditional
moment generating function. The quantity $B_{mart}$ below
is the reciprocal radius on which that moment generating
function is finite; it is not an almost-sure bound on an
increment.

Fix deterministic $u$ and its dyadic decomposition. Let
$\mathcal{F}_{k-1}$ contain all samples in the
small-support prefix and in retained groups preceding
group $k$. Let
\[
  Z_{<k}=Z_0+\sum_{\substack{r<k\\r\text{ captured}}}Z_r
\]
and define
\[
  X_k=\langle Z_{<k},\mathcal{B}V_k\rangle.
\]
Then $Z_{<k}$ is $\mathcal{F}_{k-1}$-measurable and the
current cluster samples are independent of
$\mathcal{F}_{k-1}$.

For independent $a,b\sim N(0,I_h)$ and fixed $Z$,
\[
  \E e^{\xi a^TZb}=\det(I-\xi^2ZZ^T)^{-1/2}
\]
when $|\xi|\|Z\|_{op}<1$. To verify it, condition on $a$,
integrate in $b$, and then use the Gaussian
quadratic-form formula.

Independence across coefficient indices gives
\[
  \E(e^{tX_k}\mid\mathcal{F}_{k-1}) = \prod_{i\in
  \operatorname{supp}V_k} \det\left(I-\frac{t^2(V_k)_i^2}
  {d}Z_{<k}Z_{<k}^T\right)^{-1/2}.
\]
If
\[
  |t|\frac{\|V_k\|_\infty\|Z_{<k}\|_{op}}{\sqrt d}\le
  \frac1{\sqrt2},
\]
then every eigenvalue in the logarithm is at most
$1/2$. Since $-\log(1-x)\le2x$ on $[0,1/2]$,
\[
  \E(e^{tX_k}\mid\mathcal{F}_{k-1}) \le
  \exp\left(\frac{t^2Q_k\|Z_{<k}\|_F^2}{d}\right).
\]
Thus the conditional MGF estimate holds for both signs
of $t$.

Every polar witness has squared Frobenius norm at most
its supporting rank. Orthogonality and the cumulative
support bound give
\[
  \|Z_{<k}\|_F^2 \le \sum_{l<k}k_l \le C_{q,T}S_k.
\]
The deterministic dyadic-group estimates give
\[
  Q_k\le C_q\frac{m_k^2}{S_k}
\]
and
\[
  \|V_k\|_\infty\le C_q\frac{m_k}{S_k}.
\]
Put $M=\sum_km_k=\|V\|_1$. Let
\[
  v_k=\frac{CQ_k\|Z_{<k}\|_F^2}{d}.
\]
The predictable variance sum has the deterministic upper
bound
\[
  \sum_kv_k \le \frac{C_{q,T}}d\sum_km_k^2 \le
  \frac{C_{q,T}}dM^2=:V_{mart}.
\]
Every cluster after the sparse prefix has $S_k\ge
c_{q,T}d/\log d$. Therefore the reciprocal MGF radius is
\[
  B_{mart}:=\max_k\frac{\|V_k\|_\infty}{\sqrt d} \le
  C_{q,T}\frac{M\log d}{d^{3/2}}.
\]
After the homogeneous normalization $M=1$, these are
exactly $V_{mart}\le C_{q,T}/d$ and $B_{mart}\le
C_{q,T}\log d/d^{3/2}$. The second quantity is an
MGF-radius parameter. It is not an almost-sure increment
bound.

For $0\le\lambda\le c/B_{mart}$, iteration of the
conditional MGF bound first gives
\[
  \E\exp\left(-\lambda\sum_kX_k-\lambda^2\sum_kv_k\right)
  \le1.
\]
Since $\sum_kv_k\le V_{mart}$ pathwise, it follows that
\[
  \E\exp\left(-\lambda\sum_kX_k\right)\le
  e^{\lambda^2V_{mart}}.
\]
Markov's inequality gives
\[
  \Pp\left(\sum_kX_k\le-x\right) \le
  \inf_{0\le\lambda\le c/B_{mart}}
  \exp(-\lambda x+\lambda^2V_{mart}).
\]
Choose
\[
  \lambda=\min\left\{\frac{x}{2V_{mart}},\frac{c}
  {2B_{mart}}\right\}.
\]
Then
\[
  \Pp\left(\sum_kX_k\le-x\right) \le
  \exp\left[-c\min\left\{\frac{x^2}{V_{mart}},
  \frac{x}{B_{mart}}\right\}\right].
\]
At $x=a_{\rm mart}M\sqrt d$, where $a_{\rm mart}>0$ will
be chosen below,
\[
  \frac{x^2}{V_{mart}}\ge c_{q,T}d^2,
  \frac{x}{B_{mart}}\ge c_{q,T}\frac{d^2}{\log d}.
\]
Thus
\[
  \Pp\left(\sum_kX_k<-a_{\rm mart}M\sqrt d\right) \le
  \exp\left(-c_{q,T}\frac{d^2}{\log d}\right).
\]

%%%%%%%%%%%%%%%%%%%%
\subsection{When one remaining dyadic block has
  macroscopic support}\label{macroscopic-late-driver}
%%%%%%%%%%%%%%%%%%%%

If the groups treated above carry less than half of the
coefficient mass, then one of the remaining dyadic blocks
has support comparable to $d$. We condition on every
sample outside that block. The other contributions then
form an arbitrary measurable matrix center, while the
selected block remains independent and can be handled by
the arbitrary-center small-ball estimate.

\begin{lemma}\label{record-lemma-9-1}
  Fix $c_0,K>0$. Let $S\subset[n]$, $s=|S|$, suppose
  \[
    c_0d\le s\le Kd,
  \]
  and suppose
  \[
    \lambda\le|u_i|\le2\lambda
  \]
  for $i\in S$. Let $C$ be measurable with respect to
  \[
    \mathcal{G}_S=\sigma(a_i,b_i:i\notin S).
  \]
  Then there is a numerical constant $c_{\rm ld}>0$ such
  that, for all sufficiently large $d$ (with the starting
  dimension allowed to depend on $c_0,K$), almost surely
  with respect to $\mathcal{G}_S$,
  \[
    \Pp\left(\|C+\mathcal{B}u_S\|_*<c_{\rm ld}\sqrt d
    \|u_S\|_1\sqrt{\frac{\min\{s,d\}}s}\middle|
    \mathcal{G}_S\right) \le e^{-c_{\rm ld}d\min\{s,d\}}.
  \]
\end{lemma}

\begin{proof}
  Put $\sigma=\|u_S\|_2$ and choose
  \[
    r=\lfloor\gamma\min\{s,d\}\rfloor
  \]
  with a fixed numerical $\gamma>0$ small. Since
  $s\ge c_0d$, after increasing the starting dimension
  if necessary,
  \[
    r\ge\frac\gamma2\min\{s,d\}.
  \]
  Also, since $\sigma\ge\lambda\sqrt s$,
  \[
    \sqrt r\frac{\|u_S\|_\infty}{\sigma} \le
    2\sqrt\gamma<1/\sqrt2.
  \]
  Partition both coordinate spaces into $J_0\ge cd/r$
  paired $r$ by $r$ blocks. Conditional on
  $\mathcal{G}_S$, every block of $C$ is fixed and the
  selected coefficient block uses independent Gaussian
  coordinates in distinct paired coordinate blocks. For
  each paired block, \cref{record-lemma-2-1} shows that
  the conditional probability that its nuclear norm is
  smaller than
  \[
    c\sigma\frac{r^{3/2}}{\sqrt d}
  \]
  is at most $e^{-cr^2}$. Let $N_{good}$ denote the
  number of paired blocks whose nuclear norm is at least
  this threshold. The binomial calculation in
  \cref{record-lemma-5-1} gives
  \[
    \Pp(N_{good}<J_0/2\mid\mathcal{G}_S)\le
    e^{-cJ_0r^2}\le e^{-cdr}.
  \]
  On the complementary event, the pinching estimate
  \cref{record-lemma-2-2} gives
  \[
    \|C+\mathcal{B}u_S\|_* \ge c\frac dr\sigma
    \frac{r^{3/2}}{\sqrt d} = c\sigma\sqrt{dr}.
  \]
  Flatness gives $\sigma\ge\|u_S\|_1/(2\sqrt s)$. Since
  $r\asymp\min\{s,d\}$, the last display proves the
  asserted scale and exponent.
\end{proof}

Suppose now that less than half of $M=\|V\|_1$ lies in
the small-support prefix and the groups treated by span
removal. The remaining groups carry more than $M/2$.
There are at most $J_{late}$ of them, so one remaining
group has mass at least $M/(2J_{late})$. Its base dyadic
bin has mass at least a $1/(4q)$ fraction of its group
mass, because the group upper bound is $4q\lambda_rS_r$.

Call this base bin $S$. It is deterministic once the
coefficient vector is fixed. It satisfies
\[
  c_{q,T,\rho}d\le|S|\le Kd.
\]
Reveal $\mathcal{G}_S$. The center
\[
  C=\mathcal{B}(V-u_S)
\]
is $\mathcal{G}_S$-measurable, while the samples in the
selected block remain independent of $\mathcal{G}_S$.
The macroscopic-block estimate \cref{record-lemma-9-1},
applied with $c_0=c_{q,T,\rho}$ and the preceding fixed
$K$, therefore gives, almost surely with respect to
$\mathcal{G}_S$,
\[
  \Pp\left(\|\mathcal{B}V\|_*<\frac{c_{\rm ld}\sqrt d}
  {8q\sqrt{K_+}J_{late}}M\middle|\mathcal{G}_S\right)
  \le e^{-c_{q,T,\rho,K}d^2}.
\]

%%%%%%%%%%%%%%%%%%%%
\subsection{Combining the two alternatives for one fixed
  coefficient vector}\label{fixed-vector-combination}
%%%%%%%%%%%%%%%%%%%%

By homogeneity, normalize $\|u\|_1=1$. Let $T$ be
specified below and put $M=\|V\|_1\ge1-\epsilon_T$.

Let $M_{cap}$ be the mass in the small-support prefix
and the groups treated by span removal. If $M_{cap}\ge
M/2$, the sparse and projected lower bounds give a
diagonal polar pairing at least
\[
  c\sqrt dM_{cap}\ge c\sqrt dM.
\]
By \cref{orthogonal-polar-aggregation}, the aggregate
witness has operator norm at most one and all backward
terms vanish. The supermartingale estimate in
\cref{future-cross-term-supermartingale} gives
\[
  \Pp\left(\sum_kX_k<-cM\sqrt d\right)\le
  \exp\left(-c_{q,T}\frac{d^2}{\log d}\right).
\]
Choose $a_{\rm mart}$ to be half the diagonal constant.
Then
\[
  \|\mathcal{B}V\|_*\ge c\sqrt dM.
\]
For each projected group $r$, let $F_r$ be the event
that the lower bound from \cref{record-lemma-5-1} fails.
Every such $R_r$ outside the small-support prefix is at
least $c_{q,T}d/\log d$, and there are at most
$1+\log_q(Kd)$ groups. The tower property and a union
bound give
\[
  \Pp\left(\bigcup_rF_r\right)\le C\log
  d\exp\left(-c_{q,T}\frac{d^2}{\log d}\right).
\]
If $M_{cap}<M/2$, let $F_{late}$ be the event that the
following lower bound fails:
\[
  \|\mathcal{B}V\|_* \ge \frac{c}
  {\sqrt{K_+}[1+\log_q(C_qK_+T/\rho)]}\sqrt dM.
\]
The conditional estimate in
\cref{macroscopic-late-driver}, followed by the tower
property, gives
\[
  \Pp(F_{late})\le e^{-c_{q,T,\rho,K}d^2}.
\]
Combining the branches, for every fixed deterministic
$u$,
\[
  \Pp\left(\{\|\mathcal{B}V\|_*<c_{ret}(K,T)\sqrt d
  \|V\|_1\}\cap E_{sp}\right) \le C\log
  d\exp\left(-c_{q,T,K}\frac{d^2}{\log d}\right),
\]
where
\[
  c_{ret}(K,T) \ge \frac{c}
  {\sqrt{K_+}[1+\log_q(C_qK_+T/\rho)]}.
\]
On $E_{atom}$,
\[
  \|\mathcal{B}W\|_*\le C_A\epsilon_T\sqrt d\|u\|_1.
\]
Therefore
\[
  \|\mathcal{B}u\|_* \ge \left[c_{ret}(K,T)(1-
  \epsilon_T)-C_A\epsilon_T\right]\sqrt d\|u\|_1.
\]
Now fix $q=2$ and $\rho$ as above. Set
\[
  T=\left\lceil T_0+2\log_2(eK_+)\right\rceil.
\]
The lower constants in the small-support-prefix and
projected-group arguments depend on $q$ and $\rho$, but
not on $T$, and the macroscopic-block scale constant
$c_{\rm ld}$ is numerical. The truncation depth affects
the cumulative support factor, the exceptional exponent,
the starting dimension, and the number of remaining
groups. Thus the retained lower-bound scale loses only
the factor $1+\log K_+$ displayed below.

The ratio
\[
  \frac{(1+\log(K_+T))}{K_+^{3/2}}
\]
is bounded on $K_+\ge1$. Hence $T_0$, depending only on
the atom constant, can be chosen so that
\[
  C_A\epsilon_T\le\frac14c_{ret}(K,T)
\]
for every fixed $K$. With $T_0$ increased if needed,
$\epsilon_T\le1/2$. Define $c_{\rm ms}(A)>0$ to be a
sufficiently small constant depending only on the
single-summand bound for aspect parameter $A$. It follows
that
\[
  \|\mathcal{B}u\|_* \ge \frac{c_{\rm ms}(A)}
  {\sqrt{K_+}(1+\log K_+)}\sqrt d\|u\|_1.
\]
If $F_u$ denotes the event that this inequality fails,
then, for every fixed deterministic $u$,
\[
  \Pp(F_u\cap E_{sp}\cap E_{atom})\le C\log
  d\exp\left(-c_{q,T,K}\frac{d^2}{\log d}\right).
\]

%%%%%%%%%%%%%%%%%%%%
\subsection{Uniformity over every support and
  coefficient vector}\label{whole-class-deterministic-union}
%%%%%%%%%%%%%%%%%%%%

Normalize $\|u\|_1=1$. Denote the fixed-vector constant
obtained in \cref{fixed-vector-combination} by
\[
  c^{fix}_{K,A}=\frac{c_{\rm ms}(A)}
  {\sqrt{K_+}(1+\log K_+)}.
\]
For each support $S$ with $|S|\le Kd$, take a
deterministic $\delta$-net of its $\ell_1$ unit sphere
in $\ell_1$ distance, with
\[
  \delta=\frac{c^{fix}_{K,A}}{4C_A}.
\]
A volumetric packing argument gives at most
\[
  (1+4/\delta)^{|S|}
\]
points. Thus
\[
  \log N_{coef}\le Kd\log(1+4/\delta)=O_{A,K}(d).
\]
If $Kd\le n/2$, then
\[
  \log N_{supp} \le Kd\log\left(\frac{en}{Kd}\right)
  +O_K(d)=O_{A,K}(d\log d).
\]
If $Kd>n/2$, then either $n=O_K(d)$ and
$2^n=e^{O_K(d)}$, or one may use the same upper bound
after replacing $Kd$ by $\min\{Kd,n\}$. In every case,
\[
  \log N_{supp}=O_{A,K}(d\log d).
\]
Each deterministic net point determines its dyadic
decomposition. Absolute dyadic locations, gaps between
selected support sizes, group partitions, and stopping
indices are deterministic functions of that point.
Projections and nuclear-norm witnesses are Borel
functions of the revealed Gaussian samples and are
handled by conditional probability. Thus the support and
coefficient nets account for all entropy terms.

The union-bound exponent is
\[
  -c'_{K,A}\frac{d^2}{\log d}+C_{A,K}d\log d
  \longrightarrow-\infty.
\]
Indeed, if $E_{net}$ denotes validity at every net
point, then
\[
  \Pp(E_{net}^c\cap E_{sp}\cap E_{atom}) \le
  \exp\left[-c'_{K,A}\frac{d^2}{\log d}+C_{A,K}d\log
  d\right].
\]
If $u_0$ is a net point with $\|u-u_0\|_1\le\delta$,
then on $E_{atom}$,
\[
  \|\mathcal{B}(u-u_0)\|_* \le C_A\sqrt d\delta \le
  \frac14c^{fix}_{K,A}\sqrt d.
\]
Both $u$ and $u_0$ have $\ell_1$-norm one, so the net
lower bound extends with a fixed factor loss. Combining
the last three probability estimates gives
\[
  \Pp\left(\inf_{\substack{u\ne0\\
  |\operatorname{supp}u|\le Kd}}\frac{\|\mathcal{B}u\|_*}
  {\sqrt d\|u\|_1}<\frac12c^{fix}_{K,A}\right)\to0.
\]
Renaming the constant proves \cref{main-theorem-6-1}
and hence \cref{the-clipped-head-theorem}.
   
\end{aiblock}

%%%%%%%%%%%%%%%%%%%%
\section{How the bilinear block enters the centered
  adjoint map}\label{off-diagonal-small-balls-and-product-state-continuity}
%%%%%%%%%%%%%%%%%%%%

\begin{aiblock}
 
Let $h=\lfloor d/2\rfloor$. Split each sample into
independent coordinate halves $a_i,b_i\in\mathbb R^h$; if $d$
is odd, the middle coordinate is discarded. Define
\[
  \mathcal{B}y=\frac1{\sqrt d}\sum_i y_i a_ib_i^T.
\]
The identity matrix has zero off-diagonal block, and
rectangular compression is nuclear-norm contractive, so
this projection preserves the validity of the lower bound.
Hence
\[
  \|\mathcal{T}_d^*y\|_*\ge\|\mathcal{B}y\|_*
\]
for $y\perp\1_n$.

\begin{lemma}[Arbitrary-center determinant small
  ball]\label{main-lemma-4-1}
  Fix $0<\theta<1/\sqrt2$. Uniformly over deterministic
  $z\in\mathbb R^n$ satisfying
  \[
    \|z\|_2=1, \|z\|_\infty\le\frac\theta{\sqrt d},
  \]
  and deterministic $N\in\mathbb R^{h\times h}$,
  \[
    \Pp\left(\left\|\frac1{\sqrt d}\sum_i z_i
    a_ib_i^T-N\right\|_F\le t\sqrt d\right)
    \le(C_\theta t)^{h^2}, 0<t<1.
  \]
\end{lemma}

\begin{proof}
  Apply \cref{prop:det-frobenius-small-ball} with
  $r=h$. The coefficient cap is valid because
  \[
    \sqrt h\|z\|_\infty\le\theta\sqrt{h/d}<\theta.
  \]
  Moreover,
  \[
    \frac1{\sqrt d}\sum_i z_ia_ib_i^T=\sqrt{h/d}Y_z.
  \]
  Thus a Frobenius radius $t\sqrt d$ on the left becomes
  a radius $t(d/h)\sqrt h$ for $Y_z$. Since $d/h\le3$ for
  $d\ge2$, the factor is absorbed into $C_\theta$. The
  exact coordinate-block statement used later is
  \cref{cor:det-scaled-block}.
\end{proof}

\begin{lemma}[Operator norm bound for the bilinear
  measurement matrix]\label{main-lemma-4-2}
  If $n=O(d^2)$, then
  \[
    \|\mathcal{B}\|_{2\mathbin{\to}F}=O_{\Pp}(\sqrt d).
  \]
\end{lemma}

\begin{proof}
  Use row vectorization, so that
  $\operatorname{vec}_r(ab^T)=a\otimes b$. If $V$ is the
  $h^2$ by $n$ matrix with columns $a_i\otimes b_i$, then
  \[
    \operatorname{vec}_r(\mathcal{B}y)=\frac1{\sqrt d}Vy.
  \]
  Consequently,
  \[
    \|\mathcal{B}\|_{2\mathbin{\to}F}^2 =\frac1d
    \left\|\sum_i(a_i\otimes b_i)(a_i\otimes
    b_i)^T\right\|_{op}.
  \]
  Introduce independent real copies $a_i',b_i'$ and, for
  $s,t\in\{-1,1\}$, put
  \[
    z_{i,s,t}=\frac{(a_i+\mathrm i s a_i')\otimes(b_i+
    \mathrm i t b_i')}{2h}.
  \]
  Write
  \[
    \phi_{i,s}=\frac{a_i+\mathrm i s a_i'}{\sqrt{2h}},
  \]
  and
  \[
    \psi_{i,t}=\frac{b_i+\mathrm i t b_i'}{\sqrt{2h}}.
  \]
  Then $z_{i,s,t}=\phi_{i,s}\otimes\psi_{i,t}$. Every
  complex coordinate of either factor has independent
  real and imaginary parts of variance $1/(2h)$. Hence
  \[
    \E\|\phi_{i,s}\|_2^2=\E\|\psi_{i,t}\|_2^2=1.
  \]
  Thus, for each fixed $(s,t)$, the vectors $z_{i,s,t}$
  are independent tensor products of two complex Gaussian
  vectors in $\mathbb C^h$. Ambainis, Harrow, and
  Hastings call these \emph{two-factor complex Gaussian
  product states}; this is source terminology, and the
  normalization above is exactly theirs.

  To verify the PSD domination, put
  \[
    u_i=a_i\otimes b_i,
  \]
  \[
    v_i=a_i\otimes b_i',
  \]
  \[
    r_i=a_i'\otimes b_i,
  \]
  and
  \[
    q_i=a_i'\otimes b_i'.
  \]
  Then
  \[
    2h z_{i,s,t}=u_i+\mathrm i t v_i+\mathrm i s
    r_i-stq_i.
  \]
  Since $\sum_s s=\sum_t t=\sum_{s,t}st=0$,
  summing the four Hermitian rank-one
  projectors cancels all mixed terms and
  gives
  \[
    \sum_{s,t}z_{i,s,t}z_{i,s,t}^* =\frac1{h^2}
    \left(u_iu_i^T+v_iv_i^T+r_ir_i^T+q_iq_i^T\right).
  \]
  Therefore
  \[
    (a_i\otimes b_i)(a_i\otimes b_i)^T \preceq
    h^2\sum_{s,t}z_{i,s,t}z_{i,s,t}^*.
  \]
  In the notation of
  \cite{AmbainisHarrowHastings2012} with $k=2$,
  local dim $h$, fix $x_0>\sup n/h^2$ to get
  fixed aspect $N_{\rm aug}/h^2=x_0$ for Cor~3,
  put $N_{\rm aug}=x_0h^2$, and for each
  $(s,t)$ augment by $N_{\rm aug}-n$ i.i.d.\
  states; adding PSD terms only increases
  $\|\cdot\|_{op}$, so the augmented
  $\widehat M_{s,t}^{\rm aug}$ dominates the
  original sum. Denote its normalized version
  by $M_{N_{\rm aug},h,2}^{\rm norm}$
  ($\E\|\phi\|_2^2=1$) for the concentration
  below.


  The distinction between these two ensembles matters in
  the concentration formula. Corollary 3 gives, at the
  fixed aspect $p/h^2=x_0$,
  \[
    \E\|M_{N_{\rm aug},h,2}^{\rm norm}\|_{op} \le
    (1+\sqrt{x_0})^2+o(1).
  \]
  Equation (11) of Ambainis, Harrow, and Hastings states
  that, for fixed $0<\epsilon\le1$ and $\Delta>0$,
  \[
    \begin{aligned}
      &\Pp\left(
      \left|\|\widehat M_{s,t}^{\rm aug}\|_{op}
      -\E\|M_{N_{\rm aug},h,2}^{\rm norm}\|_{op}\right|
      \ge \epsilon\E\|M_{N_{\rm aug},h,2}^{\rm norm}\|
      _{op}+\Delta \right)\\
      &\hspace{4em}\le4N_{\rm aug}
      e^{-h\epsilon^2/16}+2e^{-(h-1)\Delta^2/8}.
    \end{aligned}
  \]
  In particular, by fixing $\epsilon$ and $\Delta$,
  there are constants $C,c>0$ such that
  \[
    \Pp\left( \left\|\sum_{i=1}^n
    z_{i,s,t}z_{i,s,t}^*\right\|_{op}>C \right) \le
    e^{-ch+O(\log h)}.
  \]
  The four sign ensembles are dependent, but a union over
  four events needs no independence. On their common
  event, summing the PSD domination over $i$ and using
  vectorization gives
  \[
    \|\mathcal{B}\|_{2\mathbin{\to}F}^2 \le\frac{h^2}{d}
    \sum_{s,t} \left\|\sum_i
    z_{i,s,t}z_{i,s,t}^*\right\|_{op} \le
    C\frac{h^2}{d}.
  \]
  Since $h=\lfloor d/2\rfloor$, this is at most $Cd$.
  Hence
  \[
    \Pp\left(\|\mathcal{B}\|_{2\mathbin{\to}F}>C\sqrt d
    \right) \le e^{-cd+O(\log d)}.
  \]
  This proves the asserted tightness and a stronger
  fixed-constant high-probability form.
\end{proof}
   
\end{aiblock}


 
%%%%%%%%%%%%%%%%%%%%
\section{The uniform estimate for
  small-coordinate
  vectors}\label{the-localized-conditional-tube-theorem}
%%%%%%%%%%%%%%%%%%%%

\begin{aiblock}
 
For $R>0$, let
\[
  \mathcal B_*(R)=\{N\in\mathbb R^{h\times h}:
  \|N\|_*\le R\}.
\]
Let $B_F$ be the Frobenius unit ball in
$\mathbb R^{h\times h}$.

The nuclear-ball covering lemma
Lemma~\ref{lem:det-nuclear-cover} gives
\[
  \log N(\mathcal B_*(R),sB_F)\le C\frac{R^2h}{s^2}.
\]
At $R=cqd$ and $s=\zeta q\sqrt d$, the right side is
$C(c/\zeta)^2d^2$. Homogeneity removes $q$ exactly. This
is essential because a scale-dependent low-rank cover
would instead introduce an unnecessary $d^2\log(1/q)$
cost.

\begin{proposition}\label{prop:small-coordinate-restatement}
  Fix $\alpha_0<1/4$, $0<\theta<1/\sqrt2$, and positive
  constants $A,B,\beta,K,c_0,C_0,L$. Suppose $n\le
  \alpha_0d^2$ and work on
  \[
    E_d(L)=\{\|\mathcal B\|_{2\to F}\le L\sqrt d\}.
  \]
  Let $F_{d,L}$ be the event that the simultaneous
  conclusion below fails. There are $c,\gamma>0$ such
  that, for all sufficiently large $d$,
  \[
    \Pp(F_{d,L}\cap E_d(L))\le e^{-\gamma d^2}.
  \]
  The conclusion is the following. Simultaneously for
  every $c_0d^{-1/2}\le q\le C_0$, if $I\subset[n]$ and
  $U,V$ satisfy
  \[
    |I|\le Kd,\quad\supp U\subset I,\quad\supp
    V\subset I^c,
  \]
  \[
    \|U\|_1\le Aq\sqrt d,\quad\beta q\le\|V\|_2\le Bq,
  \]
  and
  \[
    \sqrt d\|V\|_\infty\le\theta\|V\|_2,
  \]
  then
  \[
    \|\mathcal B(U+V)\|_*\ge cqd.
  \]
\end{proposition}

\begin{proof}
  This is the restriction of \cref{main-proposition-5-1}
  to $c_0d^{-1/2}\le q\le C_0$. That proposition proves
  the stronger statement simultaneously for every $q>0$,
  with the same event-level probability bound.
\end{proof}
\end{aiblock}


 
%%%%%%%%%%%%%%%%%%%%
\section{The uniform estimate for the large-coordinate
  part}\label{the-uniform-estimate-for-the-large-coordinate-part}
%%%%%%%%%%%%%%%%%%%%

\begin{aiblock}

This theorem is the global output of
\cref{expert-multiscale-module}. It contains no flatness
assumption and is simultaneous over supports and
coefficients chosen from the same data. Its dependence on
the support-size constant $K$ is displayed because it
will be needed when $K=\eta^{-2}$ in the next section.

\begin{theorem}\label{the-clipped-head-theorem}
  Fix $0<K,A<\infty$ and assume $n\le Ad^2$. Put
  $K_+=\max\{K,1\}$. There is $c_{K,A}>0$ such that,
  uniformly over every such sequence $n=n_d$,
  \[
    \Pp\left( \inf_{\substack{u\ne0\\|\supp u|\le Kd}}
    \frac{\|\mathcal Bu\|_*}{\sqrt d\|u\|_1}<c_{K,A}
    \right)\longrightarrow0.
  \]
  Equivalently, with probability tending to one,
  uniformly over all $u$ with $|\supp u|\le Kd$,
  \[
    \|\mathcal Bu\|_*\ge c_{K,A}\sqrt d\|u\|_1.
  \]
  Moreover, there is a constant $c_{\rm ms}(A)>0$,
  depending only on $A$, such that
  \[
    c_{K,A}\ge\frac{c_{\rm ms}(A)}{\sqrt{K_+}(1+\log
    K_+)}.
  \]
\end{theorem}

\begin{proof}
  The dyadic coefficient decomposition, projected block
  estimate, orthogonal witness aggregation, conditional
  exponential estimate for future cross terms,
  macroscopic last-block alternative, and deterministic
  union bound are proved in
  \cref{expert-multiscale-module}.
\end{proof}

 \end{aiblock}
%%%%%%%%%%%%%%%%%%%%
\section{Combining the large and small
  coordinates}\label{exact-clipping-and-the-mixed-modulus}
%%%%%%%%%%%%%%%%%%%%

\begin{aiblock}
The next norm is an infimal convolution of a scaled
$\ell_1$ norm and a scaled Euclidean norm. Its minimizing
decomposition automatically separates a vector into a part
supported on at most $d/\eta^2$ coordinates and a remaining
part with a coordinatewise cap. This deterministic fact is
the exact interface between the two preceding
probabilistic estimates.

For $0<\eta<1$, define
\[
  N_\eta(y)=\inf_{y=u+v} \left[\eta\sqrt d\|u\|_1+d\|v\|_2
  \right].
\]
Minimize $\eta\sqrt d\|y-v\|_1+d\|v\|_2$ and put $u=y-v$.
If $v\ne0$, subgradient optimality gives, coordinatewise,
\[
  \frac{dv_i}{\|v\|_2}\in\eta\sqrt d\partial|y_i-v_i|.
\]
Thus
\[
  |v_i|\le\frac{\eta\|v\|_2}{\sqrt d},
\]
with equality and matching sign whenever $u_i\ne0$. If
$I=\operatorname{supp}u$, then
\[
  \|v\|_2^2\ge\sum_{i\in I}v_i^2=|I|\frac{\eta^2\|v\|_2^2}{d},
\]
so $|I|\le d/\eta^2$. If $v=0$, optimality gives an
$s\in\partial\|y\|_1$ satisfying $\eta\sqrt d\|s\|_2\le
d$. Since $|s_i|=1$ on $\operatorname{supp}y$, the same
support bound follows.
\end{aiblock}
\begin{aiblock}
For $v\ne0$, put
\[
  I=\operatorname{supp}u, U=y_I=u+v_I, V=y_{I^c}=v_{I^c}.
\]
Then $U,V$ have disjoint supports. Moreover,
\[
  \eta\sqrt d\|U\|_1\le\eta\sqrt d\|u\|_1+d\|v\|_2=
  N_\eta(y),
\]
where the first inequality uses
$\|v_I\|_1\le\sqrt{|I|}\|v\|_2$, and $\|V\|_2\le\|v\|_2\le
N_\eta(y)/d$. If $\|y\|_2=1$ and
\[
  q_y=\frac{N_\eta(y)}d,
\]
then $q_y\le1$ by choosing $v=y$. Conversely, for every
$y=u+v$,
\[
  \eta\sqrt d\|u\|_1+d\|v\|_2\ge\eta\sqrt d(\|u\|_2+
  \|v\|_2)\ge\eta\sqrt d,
\]
so
\[
  \frac\eta{\sqrt d}\le q_y\le1,
\]
\[
  \|U\|_1\le\frac{q_y\sqrt d}{\eta}, \|V\|_2\le q_y.
\]
If $\|V\|_2\ge\beta q_y$, then
\[
  \sqrt d\|V\|_\infty\le\frac\eta\beta\|V\|_2.
\]
On $I$, the matching-sign relation gives
\[
  \eta\sqrt d\|v_I\|_1=d\frac{\|v_I\|_2^2}{\|v\|_2}.
\]
Consequently,
\[
  \eta\sqrt d\|U\|_1 =N_\eta(y)-d\frac{\|V\|_2^2}
  {\|v\|_2} \ge N_\eta(y)-d\|V\|_2.
\]

\begin{theorem}[Mixed-norm lower bound for the adjoint
  map]\label{main-theorem-7-1}
  Assume $\limsup n/d^2<1/4$. For every $\delta>0$, there
  are fixed constants $\eta_\delta,c_\delta>0$ and events
  $G_{d,\delta}$ such that
  \[
    \limsup_d\Pp(G_{d,\delta}^c)\le\delta
  \]
  and, on $G_{d,\delta}$,
  \[
    \|\mathcal{T}_d^*y\|_*\ge c_\delta N_{\eta_\delta}(y)
  \]
  for every $y\perp\1_n$.
\end{theorem}

\begin{proof}
  By homogeneity, let $\|y\|_2=1$. Tightness in
  \cref{main-lemma-4-2} gives, for fixed $\delta$, a fixed
  $L_\delta$ such that
  \[
    \limsup_d\Pp(E_d(L_\delta)^c)\le\delta.
  \]
  For the multiscale theorem, use the fixed aspect bound
  $A_0=1$, since $n\le\alpha_0d^2<d^2$. Its quantitative
  constant gives, for every $0<\eta<1$,
  \[
    \frac{c_{\eta^{-2},A_0}}{\eta} \ge\frac{c}{1+\log
    (1/\eta)}.
  \]
  Choose $\eta=\eta_\delta$ so small that
  \[
    2CL_\delta\eta <\frac12\frac{c}{1+\log(1/\eta)}
    (1-2\eta).
  \]
  This is possible because $\eta[1+\log(1/\eta)]\to0$.
  Now set
  \[
    K=\eta^{-2},
  \]
  \[
    \beta=2\eta,
  \]
  and
  \[
    c_H(\eta)=\frac{c_{K,A_0}}{\eta}.
  \]
  Then $\eta/\beta=1/2<1/\sqrt2$, and
  \cref{the-clipped-head-theorem} gives
  \[
    \|\mathcal{B}U\|_* \ge c_{K,A_0}\sqrt d\|U\|_1.
  \]
  On $E_d(L_\delta)$,
  \[
    \|\mathcal{B}z\|_* \le\sqrt h\|\mathcal{B}z\|_F \le
    C L_\delta d\|z\|_2.
  \]
  The chosen inequality remains true with the actual
  $c_H(\eta)$ in place of its lower bound. The absorption
  uses the rescaled large-coordinate constant
  $c_H=c_{K,A_0}/\eta$, not the unscaled constant
  $c_{K,A_0}$.

  If $\|V\|_2\ge\beta q_y$,
  \cref{prop:small-coordinate-restatement} applies with
  the fixed parameters
  \[
    A=2/\eta,\ B=1,\ \beta=2\eta,\ K=\eta^{-2},\
    \theta=1/2,\ c_0=\eta,\ C_0=1,\ L=L_\delta,
  \]
  and gives $\|\mathcal{B}(U+V)\|_*\ge cN_\eta(y)$.

  If $\|V\|_2<\beta q_y$, then
  \[
    \|\mathcal{B}U\|_* \ge c_H(\eta)(1-\beta)N_\eta(y),
  \]
  while
  \[
    \|\mathcal{B}V\|_* \le CL_\delta\beta N_\eta(y).
  \]
  The parameter inequality absorbs this error and leaves
  half the head contribution. The case $v=0$ is purely
  sparse and follows directly from
  \cref{the-clipped-head-theorem}.

  Let $E_{\rm large}(K)$ be the simultaneous
  large-coordinate event in
  \cref{the-clipped-head-theorem}, and let
  $F_{d,L_\delta}$ be the bad event in
  \cref{prop:small-coordinate-restatement} with the
  parameters displayed above. Define
  \[
    G_{d,\delta}=E_d(L_\delta)\cap E_{\rm large}(K)\cap
    F_{d,L_\delta}^c.
  \]
  Then
  \[
    \Pp(G_{d,\delta}^c)\le\Pp(E_d(L_\delta)^c)+
    \Pp(E_{\rm large}(K)^c)+\Pp(F_{d,L_\delta}\cap
    E_d(L_\delta)).
  \]
  The first term has limiting upper bound at most
  $\delta$, the second tends to zero, and the third is
  at most $e^{-c_\delta d^2}$. Taking the minimum of the
  two branch constants proves the theorem on
  $G_{d,\delta}$.
\end{proof}

\end{aiblock}

%%%%%%%%%%%%%%%%%%%%
\section{Exact correction and preservation of positive
  definiteness}\label{quotient-exactification-and-sat-below-one-quarter}
%%%%%%%%%%%%%%%%%%%%


\begin{aiblock}
 
The dual norm of $N_\eta$ on the full coefficient space
is
\[
  N_\eta^*(e)=\max\left\{ \frac{\|e\|_\infty}
  {\eta\sqrt d}, \frac{\|e\|_2}{d} \right\}.
\]
Indeed, the polar of the infimal-convolution unit ball
is the intersection of the $\ell_\infty$ ball of radius
$\eta\sqrt d$ and the Euclidean ball of radius $d$.
Restriction to $\1_n^\perp$ can only reduce the dual
support function, so the displayed formula remains a
valid upper bound there.

Almost surely the rank-one matrices $x_ix_i^T$ are
linearly independent whenever $n\le d(d+1)/2$. This rank
event is the nonvanishing of a polynomial, and it is
nonempty because the matrices
\[
  e_je_j^T\text{ and }(e_j+e_k)(e_j+e_k)^T\text{ for
  }1\le j<k\le d
\]
span $\operatorname{Sym}_d$: subtracting the two
relevant diagonal matrices from the second type recovers
$e_je_k^T+e_ke_j^T$. Thus $\mathcal{A}_d$ is surjective
and $\mathcal{T}_d$ is surjective onto $\1_n^\perp$.

\begin{lemma}[Finite-dimensional quotient
  duality]\label{lem:quotient-duality}
  Let $T:X\to Y$ be a surjective linear map between
  finite-dimensional normed spaces. Then, for every
  $e\in Y$,
  \[
    \inf_{Tx=e}\|x\|_X =\sup_{\substack{y^*\in Y^*\\
    \|T^*y^*\|_{X^*}\le1}} |\langle e,y^*\rangle|.
  \]
  The infimum is attained.
\end{lemma}

\begin{proof}
  The left side is the quotient norm induced by
  $X/\ker T$. Its dual norm is
  $y^*\mapsto\|T^*y^*\|_{X^*}$. Duality of a
  finite-dimensional norm with its polar gives the
  equality. Attainment follows by intersecting the
  closed affine fiber $T^{-1}(e)$ with a sufficiently
  large compact norm ball.
\end{proof}

Define
\[
  \kappa_d(S)=\inf\{\|H\|_{op}:\cT_d(H)=-\PiZero\rvec(S)\}.
\]
For $e=\PiZero\rvec(S_d^*)$,
Lemma~\ref{lem:quotient-duality} gives
\[
  \kappa_d(S_d^*) =\inf_{\mathcal{T}_d(H)=-e}\|H\|_{op}
  =\sup_{\substack{y\perp\1_n\\\|\mathcal{T}_d^*y\|_*
  \le1}} |\langle e,y\rangle|.
\]
On $G_{d,\delta}$,
\[
  \kappa_d(S_d^*) \le c_\delta^{-1} \max\left\{
  \frac{\|\PiZero\rvec(S_d^*)\|_\infty}{\eta_\delta
  \sqrt d}, \frac{\|\PiZero\rvec(S_d^*)\|_2}{d}
  \right\}.
\]
For every fixed $\delta$, \cref{main-proposition-3-1}
makes the right side tend to zero in probability.
Therefore, for every $r>0$,
\[
  \limsup_d\Pp(\kappa_d(S_d^*)>r)\le\delta.
\]
Letting $\delta\downarrow0$ after the $d$ limit gives
$\kappa_d(S_d^*)\xrightarrow{\Pp}0$.

The quotient infimum is attained because the affine
constraint is closed and an operator-norm ball is compact
in finite dimension. Choose an attainer $H_d$; no
measurable selection is needed for the pathwise
existence conclusion. Since $S_d^*\succeq\lambda_-I$,
\[
  Q_d=S_d^*+H_d\succeq\frac{\lambda_-}2I
\]
with probability tending to one. Also
$\PiZero\mathcal{A}_d(Q_d)=0$, so
$\mathcal{A}_d(Q_d)=c_d\1_n$. Positivity gives $c_d>0$.
Then
\[
  \widehat S_d=\frac d{c_d}Q_d
\]
is positive definite and belongs to $\mathcal{F}_d$.
This proves the feasible assertion in \cref{thm:main}.

\end{aiblock}




% %%%%%%%%%%%%%%%%%%%% -Part-
% \part{Infeasibility above the
%   threshold}\label{part-infeasibility-above-the-threshold}
% %%%%%%%%%%%%%%%%%%%%



%%%%%%%%%%%%%%%%%%%%
\section{Leave-one-row interpolation for diffuse
  positive semidefinite
  matrices}\label{expert-leave-one-row-module}
%%%%%%%%%%%%%%%%%%%%


\begin{humanblock}


This section contains the full comparison argument before
it is used in the proof of infeasibility. The phrase
leave-one-row interpolation means the following concrete
construction. During the replacement of row $m$, we
condition on all rows except row $m$ and define one Gibbs
measure from those rows. That conditional measure is
independent of both possible choices for the current row.
The comparison then follows from an exact third-order
interpolation identity.

\end{humanblock}
\begin{aiblock}

 % BEGIN INLINED FILE: leave_one_row_module.tex
%%%%%%%%%%%%%%%%%%%%
\subsection{Uniform lower bound for diffuse
  matrices}\label{structured-diffuse-theorem}
%%%%%%%%%%%%%%%%%%%%

This section proves \cref{thm:diffuse-restatement}. We
begin by restating it with the local notation used in the
comparison argument. The theorem is uniform over all
diffuse positive semidefinite matrices and all vectors in
one fixed low-dimensional subspace. The subspace may be
random only after conditioning has made it deterministic
and independent of the Gaussian rows; this point is checked
explicitly at the end of the section in
Section \ref{conditional-nuisance-spaces-in-the-peeling-argument}.



\begin{theorem}\label{main-theorem-10-1}
  Fix $\epsilon>0$, $0<B<\infty$, and $0<C_1<\infty$.
  There are constants
  \[
    0<\tau\le1, c>0, c'>0,
  \]
  depending only on $\epsilon,B,C_1$, with the following
  property.
  
  Let $w_1,\ldots,w_n$ be independent $N(0,I_p)$
  vectors, and put
  \[
    Z_i=\frac{w_iw_i^T-I_p}{\sqrt2}.
  \]
  Let $L\subseteq\mathbb R^n$ be a deterministic
  linear subspace of dimension $m_L$. Assume
  \[
    n-m_L\ge\left(\frac14+\epsilon\right)p^2
  \]
  and
  \[
    n\le C_1p^2.
  \]
  Then, for all sufficiently large $p$,
  \[
    \Pp\left( \inf_{\substack{D\succeq0,\ \|D\|_F=1,\
    \|D\|_{op}\le\tau\\ v\in L,\ \|v\|_1\le Bn}}
    \frac1n\sum_{i=1}^n|\langle Z_i,D\rangle-v_i|<c
    \right) \le e^{-c'p^2}.
  \]
  The constants and the probability estimate are uniform
  over the orientation of deterministic $L$. The result is
  not simultaneous over subspaces chosen from the same
  $w_i$ without conditioning. The conditional form used in
  the ellipsoid proof is stated in
  \cref{conditional-nuisance-spaces-in-the-peeling-argument}.
\end{theorem}

%%%%%%%%%%%%%%%%%%%%
\subsection{Coordinates and the Gaussian comparator}\label{coordinates-and-the-gaussian-comparator}
%%%%%%%%%%%%%%%%%%%%
Equip $\operatorname{Sym}_p$ with the Frobenius inner product. For the standard basis $e_1, \ldots, e_p$, define
\[
    Q_{aa} = e_a e_a^\top
\]
and, when $a < b$,
\[
    Q_{ab} = \frac{e_a e_b^\top + e_b e_a^\top}{\sqrt{2}}.
\]
These matrices form an orthonormal basis of $\operatorname{Sym}_p$.

\begin{lemma} \label{lem:infeasibility_above:wick}
    For $1 \leq i \leq n$, let $w_1, \ldots, w_n$ be independent $N(0,I_p)$ vectors and let
    \[
        Z_i = (w_i w_i^\top - I_p) / \sqrt{2}.
    \]
    For any $A, B \in \operatorname{Sym}_p$ and $1 \leq i \leq n$,
    \[
        \E \langle Z_i, A \rangle\langle Z_i, B \rangle = \langle A, B \rangle.
    \]
\end{lemma}
 

 
\begin{proof}

    To lighten notation we drop the subscript $i$, $1 \leq i \leq n$, from $Z$ and $w$, since the argument holds for any $i$.
    
    Instead, in this proof the subscript refers to the specific component of $w$. For example, $w_a$, $1 \leq a \leq p$, refers to the $a$-th entry of the vector $w$ which is $N(0, 1)$.

    The coordinates of $Z = (w w^\top - I_p) / \sqrt{2}$ are
    \[
        \langle Z, Q_{aa} \rangle = \frac{w_a^2 - 1}{\sqrt{2}}
    \]
    and
    \[
        \langle Z, Q_{ab} \rangle = w_a w_b.
    \]
    
    For any $A \in \operatorname{Sym}_p$,
    \[
        \langle Z, A \rangle = \frac{1}{\sqrt{2}}(w^\top A w - \operatorname{Tr} A).
    \]
    
    For $w \sim N(0, I_p)$, Wick's identity yields
      
    \[
        \E[w_a w_b w_c w_d] = \delta_{ab} \delta_{cd} + \delta_{ac} \delta_{bd} + \delta_{ad} \delta_{bc},
    \]
    which equals $3$ when $a = b = c = d$, $1$ when indices pair up, $0$ otherwise.
    Now,
    \[
        w^\top A w \, w^\top B w = \sum_{a,b,c,d} A_{ab} B_{cd} w_{a} w_{b} w_{c} w_{d}.
    \]
    Taking expectation, one gets
  
    \begin{align*}
        &\E[w^\top A w \, w^\top B w] \\
            &\quad = \sum_{a,b,c,d} A_{ab} B_{cd} \delta_{ab} \delta_{cd}
            +\sum_{a,b,c,d}A_{ab} B_{cd} \delta_{ac} \delta_{bd}
            +\sum_{a,b,c,d}A_{ab} B_{cd} \delta_{ad} \delta_{bc} \\
            &\quad = \operatorname{Tr} A \operatorname{Tr} B
            + \operatorname{Tr}(A^\top B)
            + \operatorname{Tr}(A B) \\
            &\quad = \operatorname{Tr} A \operatorname{Tr} B + 2 \operatorname{Tr}(A B).
    \end{align*}
    Since $\E[w^\top A w] = \operatorname{Tr} A$ and $\E[w^\top B w] = \operatorname{Tr} B$,
    \begin{align*}
        \E[(w^\top A w - \operatorname{Tr} A) (w^\top B w - \operatorname{Tr} B)]
            &= \E[w^\top A w \, w^\top B w] - \operatorname{Tr} A \, \operatorname{Tr} B \\
            &= 2 \operatorname{Tr}(AB).
    \end{align*}
    Therefore,
    \[
        \E \langle Z, A \rangle\langle Z, B \rangle
            = \frac{1}{2} \E[(w^\top A w - \operatorname{Tr} A)(w^\top B w - \operatorname{Tr} B)]
            = \operatorname{Tr}(AB).
    \]

\end{proof}

Thus, if we define $G = \sum_{a\le b} g_{ab} \, Q_{ab}$ where $g_{ab}$ are independent $N(0,1)$, then
\[
    \E[G] = 0
    \qquad \text{and} \qquad
    \E \langle G, A \rangle\langle G, B \rangle = \operatorname{Tr}(AB),
\]
i.e., $G$ is a covariance-matched isonormal Gaussian version of $Z$ on $\operatorname{Sym}_p$.

From now on, we write $G_1, \dots, G_n$ for the copies of $G$ corresponding to $w_i$, $1 \leq i \leq n$, and use single-index subscript $\{Q_a\}_{1 \leq a \leq p (p + 1) / 2}$ to enumerate the basis of $\operatorname{Sym}_p$.



\begin{lemma} \label{lem:infeasibility_above:gaussian_width}
    Let
    \[
        \mathcal P_p = \{D \in \operatorname{Sym}_p : D \succeq 0 \text{~and~} \|D\|_F = 1\}.
    \]
    Then its Gaussian width, $w(\mathcal P_p)$, satisfies
    \[
        w(\mathcal P_p) \le \sqrt{\frac{p (p + 1)}{4}}.
    \]
\end{lemma}

\begin{proof}
    Recall that the Gaussian width of $\mathcal P_p$ is defined as
    \[
        w(\mathcal P_p) = \E \sup_{D \in \mathcal P_p} \langle H, D \rangle
    \]
    where $H$ is a symmetric isonormal Gaussian operator.
    Now,
    \[
        H = H_{+} - H_{-},\;
        H_{\pm} \succeq 0,\;
        H_{+} H_{-} = 0,\;
        \|H\|_{F}^{2} = \|H_{+}\|_{F}^{2} + \|H_{-}\|_{F}^{2}.
    \]
    Therefore,
    \[
        \sup_{D \in \mathcal P_p} \langle H, D \rangle \le \|H_+\|_F
    \]
    and 
    \[
        \E\|H\|_{F}^{2} = \E\|H_{+}\|_{F}^{2} + 
        \E\|H_{-}\|_{F}^{2}.
    \]
    
    Now, since the law of $H$ is invariant under $H \mapsto -H$,
    \[
    \E\|H_+\|_F^2=\frac12\E\|H\|_F^2=\frac{p(p+1)}4.
    \]
    
    By Jensen's inequality, the Gaussian width satisfies
    \[
    w(\mathcal P_p) \le \sqrt{\frac{p (p + 1)}{4}}.
    \]

\end{proof}


\begin{lemma} \label{lem:infeasibility_above:projected_rows}
    Let $G_1, \dots, G_n$ be mutually independent isonormal Gaussian operators on $\operatorname{Sym}_p$.
    Let $L \subseteq \mathbb R^n$ be a deterministic linear subspace of dimension $m_L$.
    Fix a deterministic retained row set $J \subseteq [n]$ with $|J| = m$. Let $L_J = \{v_J : v \in L\}$.
    If $U_J$ is deterministic and has orthonormal rows spanning $L_J^\perp$, then $U_J G_J$ is an $r_J$-row isonormal Gaussian operator, where
    \[
        r_J = m - \dim L_J \ge m - m_L.
    \]
\end{lemma}

\begin{proof}
    Since $L$ is deterministic, so is $L_J$, and $\dim L_J \le \dim L = m_L$ by the rank-nullity theorem.
    Hence
    \[
        r_J = m - \dim L_J \ge m - m_L.
    \]
    If $r_J = 0$ there is nothing else to prove, so assume $r_J \geq 1$.
    
    Recall that $\{Q_a\}_a$ is an orthonormal basis of $\operatorname{Sym}_p$ and write $\xi_{ia} = \langle G_i, Q_a\rangle = G_i(Q_a)$.
    Since $G_1, \dots, G_n$ are independent isonormal, then $\xi_{ia}$ are independent $N(0, 1)$ in $a$, with
    \[
        \E[\xi_{ia} \xi_{jb}] = \delta_{ij} \delta_{ab}.
    \]
    
    Put $G_J = (G_j)_{j \in J}$. Then
    \[
        (U_J G_J)_r = \sum_{j \in J} U_{rj} G_j,
        \qquad
        (U_J \xi)_{ra} = \sum_{j \in J} U_{rj} \xi_{ja},
    \]
    and
    \begin{align*}
        \E[(U_J \xi)_{ra} (U_J \xi)_{sb}] &=\sum_{j, k} U_{rj} U_{sk} \E[\xi_{ja} \xi_{kb}] \\
        &=\sum_j U_{rj} U_{sj} \delta_{ab} = (U_J U_J^\top)_{rs} \delta_{ab} = \delta_{rs}\delta_{ab},
    \end{align*}
    since rows are orthonormal.
    
    Jointly Gaussian with identity covariance implies that $(U_J \xi)_{ra}$ are independent $N(0,1)$ in $r$, $a$, and the claim follows.
\end{proof}




 

%%%%%%%%%%%%%%%%%%%%
\subsection{Gaussian escape inequality}\label{ssec:infeasibility_above:gaussian_escape}
%%%%%%%%%%%%%%%%%%%%

For completeness, the exact Gaussian escape inequality used here follows from the Gaussian min-max comparison.
Recall that
\[
    \mathcal P_p = \{D \in \operatorname{Sym}_p : D \succeq 0 \text{~and~} \|D\|_F = 1\}.
\]

\begin{lemma} \label{lem:infeasibility_above:gaussian_escape}
    Let $G_{1}, \dots, G_{r}$ be independent isonormal Gaussians on $\operatorname{Sym}_p$, and let $A : \operatorname{Sym}_p \to \mathbb R^{r}$ be a random linear operator defined by
    \[
        (A D)_{j} = \langle G_{j}, D \rangle \qquad \text{for $j = 1, \ldots, r$}.
    \]
    Let $g_{r} \sim N(0, I_{r})$ be independent of $A$.
    
    Then for $t > 0$
    \[
        \mathbb P\Bigl(\inf_{D \in \mathcal P_{p}} \|A D\|_{2} < \E \|g_{r}\|_{2} - w(\mathcal P_{p}) - t\Bigr) \le e^{-t^{2} / 2}.
    \]
\end{lemma}

\begin{proof}
    For any $O : \operatorname{Sym}_p \to \mathbb R^{r}$ define
    \[
        f(O) = \inf_{D \in \mathcal P_{p}} \|O D\|_{2}.
    \]
    It is not hard to see that
    \[
        f(O) = \inf_{D \in \mathcal P_{p}} \sup_{y \in S^{r - 1}} \langle y, O D \rangle,
    \]
    where
    \[
        S^{r - 1} = \{y \in \mathbb R^{r} : \|y\|_{2} = 1\},
    \]
    so $f$ is $1$-Lipschitz.

    Let $h$ be isonormal on $\operatorname{Sym}_p$, independent of $A$ and $g_r$. On $\mathcal P_{p}\times S^{r - 1}$ set
    \[
        X_{D,y} = \langle g_{r}, y \rangle + \langle h, D \rangle,
        \qquad
        Y_{D,y} = \langle y, A D \rangle = \sum_{j = 1}^r y_{j} \langle G_{j}, D \rangle.
    \]
    Both are centred Gaussian, continuous in $(D, y)$ for fixed $A$, so finite Theorem A \cite[Theorem A]{Gordon1988} extends to compact sets as in \cite[Cor 1.2]{Gordon1988}.
    Theorem A states that if
    \begin{align*}
        \E|X_{ij} - X_{ik}|^{2} &\le \E|Y_{ij} - Y_{ik}|^{2}, \\
        \E|X_{ij} - X_{\ell k}|^{2} &\ge \E|Y_{ij} - Y_{\ell k}|^{2} \quad \text{for $i \ne \ell$},
    \end{align*}
    then $\E\min_{i}\max_{j} X_{ij} \le \E\min_{i}\max_{j} Y_{ij}$.
    Equal variances are not required, only increments. In our case $i = D$, $j = y$.

    Put $a = \langle y, y' \rangle \in [-1,1]$, $b = \langle D, D' \rangle_{F} \in [0,1]$ for $D, D' \succeq 0$.
    For $\|D\|_{F} = 1$, $\|y\|_{2} = 1$,
    \begin{align*}
        \E\langle G_{j}, D \rangle \langle G_{j}, D' \rangle &= b, \\
        \E(y_{j} \langle G_{j}, D \rangle - y'_{j} \langle G_{j}, D' \rangle)^{2} &= y_{j}^{2} + y_{j}'^{2} - 2 b y_{j} y'_{j}.
    \end{align*}
    Summing in $j$,
    \begin{align*}
        \E|Y_{D, y} - Y_{D', y'}|^{2} &= 2 (1 - a b), \\
        \E|X_{D, y} - X_{D', y'}|^{2} &= \|y - y'\|_{2}^{2} + \|D - D'\|_{F}^{2} = 2 (1 - a) + 2 (1 - b).
    \end{align*}
    Hence
    \[
        \E|X - X'|^{2} - \E|Y - Y'|^{2} = 2 (1 - a) (1 - b) \ge 0
    \]
    with equality if $D' = D$ ($b = 1$) or $y' = y$ ($a = 1$).
    Thus, conditions of Theorem A hold, and
    \[
        \E \min_{D} \max_{y} X_{D, y} \le \E \min_{D} \max_{y} Y_{D, y}.
    \]

    Now $\max_{y \in S^{r - 1}} \langle g_{r}, y \rangle = \|g_{r}\|_{2}$, $\max_{y} \langle y, A D \rangle = \|A D\|_{2}$, so
    \[
        \max_{y}X_{D,y}= \|g_{r}\|_{2} + \langle h, D \rangle,
        \qquad
        \max_{y}Y_{D,y}= \|AD\|_{2},
    \]
    therefore
    \[
        \min_{D} \max_{y} X_{D,y} = \|g_{r}\|_{2} + \inf_{D}\langle h, D \rangle,
        \qquad
        \min_{D} \max_{y} Y_{D,y} = f(A).
    \]
    With $w(\mathcal P_{p}) = \E\sup_{D}\langle h, D\rangle$ and $-h\stackrel{d}=h$,
    \[
        \E\inf_{D}\langle h,D\rangle=-w(\mathcal P_{p})
    \]
    and $g_{r} \perp h$, hence
    \[
        \E f(A)
            = \E \min_{D} \max_{y} Y \ge \E\min_{D} \max_{y} X
            = \E \|g_{r}\|_{2} - w(\mathcal P_{p}).
    \]

    Since $f$ is $1$-Lipschitz, by Gaussian concentration \cite[Theorem 3.2]{Gordon1988},
    \[
        \Pp\bigl(f(A) < \E f(A) - t\bigr) \le e^{-t^{2}/2} \qquad \text{for $t > 0$}.
    \]
    Therefore,
    \[
        \Pp\Bigl(\inf_{D}\|A D\|_{2} < \E\|g_{r}\|_{2} - w(\mathcal P_{p}) - t\Bigr) \le e^{-t^{2}/2}.
    \]
\end{proof}



%%%%%%%%%%%%%%%%%%%%
\subsection{A robust Gaussian loss margin}\label{a-robust-gaussian-loss-margin}
%%%%%%%%%%%%%%%%%%%%

\begin{lemma} \label{lem:infeasibility_above:dimensionality_bounds}
    Let $G_1, \dots, G_n$ be mutually independent isonormal Gaussian operators on $\operatorname{Sym}_p$.
    Fix $\epsilon > 0$ and $0 < C_1 < \infty$.
    Let $L \subseteq \mathbb R^n$ be a deterministic linear subspace of dimension $m_L$, and assume that
    \[
        n - m_L \ge a_0 p^2
        \qquad
        \text{and}
        \qquad
        n \le C_1 p^2,
    \]
    where
    \[
        a_0 = \frac{1}{4} + \epsilon.
    \]
    Then, there exists $\gamma = \gamma(\epsilon) > 0$, such that except on an event of probability at most
    \[
        \exp\left(-\frac{\gamma^2 p^2}{4}\right),
    \]
    the inequality
    \[
        \inf_{D \in \mathcal P_p} \inf_{v \in L} \|(\langle G_i, D\rangle - v_i)_{i \in J}\|_2 \ge\gamma p
    \]
    holds simultaneously for all retained row sets $J \subseteq [n]$ with $|J| = m$.
\end{lemma}

\begin{proof}
    Set
    \[
        \Delta = \sqrt{\frac{1}{4} + \frac{\epsilon}{2}} - \frac{1}{2}.
    \]
    Choose $0 < \theta < 1/2$ small enough so that
    \[
        \theta \, C_1 \le \frac{\epsilon}{2}
        \qquad \text{and} \qquad
        C_1 H_2(\theta) \le \frac{1}{8}\left(\frac{\Delta}{4}\right)^2,
    \]
    where $H_2(t) = -t \log(t) - (1 - t) \log(1 - t)$. This is possible because $H_2(t) \to 0$ as $t \downarrow 0$. Define
    \[
        \gamma = \frac{\Delta}{4}, k = \lfloor\theta n\rfloor, m = n - k.
    \]
    
    Lemma~\ref{lem:infeasibility_above:projected_rows}, the dimensional assumptions, and the choice of $\theta$ give
    \[
        r_J = m - \dim L_J \ge n - m_L - k \ge \left(\frac{1}{4} + \frac{\epsilon}{2}\right)p^2.
    \]
    
    Since, $\|\cdot\|_{2}$ is $1$-Lipschitz, so $\operatorname{Var}\|g_{r_J}\|_{2} \le 1$ and $\E\|g_{r_J}\|_{2}^{2} = r_J$, whence $\E\|g_{r_J}\|_{2} \ge \sqrt{r_J - 1}$.
    Consequently, combining the bounds above with Lemma~\ref{lem:infeasibility_above:gaussian_width}, for all sufficiently large $p$,
    \[
        \E\|g_{r_J}\|_2 - w(\mathcal P_p) \ge 2 \gamma p.
    \]
    
    Using $t = \gamma p$ in Lemma~\ref{lem:infeasibility_above:gaussian_escape} gives
    \[
        \Pp\left( \inf_{D \in \mathcal P_p}\inf_{v \in L} \|(\langle G_i, D\rangle - v_i)_{i \in J}\|_2 < \gamma p \right) \le e^{-\gamma^2 p^2 / 2}.
    \]
    
    There are at most
    \[
        \binom{n}{k} \le e^{n H_2(\theta)}
    \]
    possible retained sets. Moreover,
    \[
        n H_2(\theta) \le C_1 p^2 H_2(\theta) < \frac{\gamma^2 p^2}{4}.
    \]
    Hence, except on an event of probability at most
    \[
        \exp\left(-\gamma^2 p^2 / 4\right),
    \]
    the inequality
    \[
        \inf_{D \in \mathcal P_p} \inf_{v \in L} \|(\langle G_i, D\rangle - v_i)_{i \in J}\|_2 \ge\gamma p
    \]
    holds simultaneously for all $J$ with $|J| = m$.
\end{proof}


Following the notation used in the proof, define the fixed coordinate threshold
\[
    a = \frac{\gamma}{2\sqrt{C_1}}.
\]

Fix $D \in \mathcal P_p$ and $v \in L$. If at most $k$ residual coordinates
\[
    r_i = \langle G_i, D\rangle - v_i
\]
satisfied $|r_i| \ge a$, one could select $m = n - k$ coordinates with $|r_i| < a$. On this retained set,
\[
    \|r_J\|_2 < a \sqrt n \le \frac{\gamma p}{2},
\]
contradicting the simultaneous lower bound. Thus more than $k$ coordinates satisfy $|r_i| \ge a$. Since $k + 1 > \theta n$, their fraction is larger than $\theta$.

Choose a fixed smooth function $\chi : [0,\infty) \to [0,1]$ such that $\chi(s) = 0$ for $s \le 1/4$ and $\chi(s) = 1$ for $s \ge 1$, and put
\begin{align}
    \ell(t) = \chi(t^2/a^2). \label{eq:infeasibility_above:smoothing}
\end{align}

Then $\ell \in C^3(\mathbb R; [0,1])$, $\ell(t)=1$ when $|t| \ge a$, and
\[
    \ell(t) \le C_\ell|t|, \qquad C_\ell = \frac{2}{a}.
\]
On the preceding Gaussian event,
\begin{align}
    \inf_{D \in \mathcal P_p} \inf_{v \in L} \frac{1}{n} \sum_{i = 1}^n \ell(\langle G_i, D\rangle - v_i) \ge \theta. \label{eq:infeasibility_above:smooth_lower_bound}
\end{align}
This is the required robust Gaussian bounded-loss margin. Every parameter in it depends only on $\epsilon$ and $C_1$.


%%%%%%%%%%%%%%%%%%%%
\subsection{\texorpdfstring{The empirical
  Frobenius-to-$\ell_1$ event}{The empirical
  Frobenius-to-l1 event}}\label{the-empirical-frobenius-to-ell_1-event}
%%%%%%%%%%%%%%%%%%%%

In the previous section we established a lower bound for the worst-case scenario of $\frac{1}{n} \sum_{i = 1}^n \ell(\langle G_i, D\rangle - v_i)$.
However, $G_i$'s are only covariance-matched isonormal Gaussian versions of $Z_i$'s.
To start bridging that gap, in this section we provide a uniform $O(1)$ bound on $\frac{1}{n} \sum_{i = 1}^n |\langle Z_i, D \rangle|$.

\begin{lemma} \label{lem:infeasibility_above:average_z_bound}
    Define
    \[
        \Lambda_p = \sup_{\substack{A \in \operatorname{Sym}_p \\ \|A\|_F = 1}} \frac{1}{n} \sum_{i = 1}^n |\langle Z_i, A \rangle|.
    \]
    Then there exist $A_* = A_*(\epsilon) < \infty$ and a universal constant $c_2 > 0$, such that
    \[
        \Pp(\Lambda_p > A_*) \le 2 \exp(-c_2 \, a_0 \, p^2).
    \]
\end{lemma}


\begin{proof}
    For a fixed $A \in \operatorname{Sym}_p$ with $\|A\|_F=1$, Hanson-Wright inequality gives
    \[
        \Pp(|\langle Z_i, A \rangle| > t) \le 2 \exp[-c_0 \min(t^2, t)],
    \]
    where $c_0$ is a universal constant since $\|A\|_F = 1$ implies that $\|A\|_{op} \le 1$.
    Recall that the subexponential norm is defined as
    \[
        \|X\|_{\psi_1} = \inf\{t > 0 : \E\exp(|X| / t) \le 2\}.
    \]
    The preceding tail bound implies $\|\langle Z_i, A \rangle\|_{\psi_1} \le C_0$ uniformly in $A$, and from Lemma~\ref{lem:infeasibility_above:wick} and Jensen one has $\E|\langle Z_i, A \rangle| \le 1$.
    
    Let $Y_{i}(A) = |\langle Z_{i}, A \rangle|$. To lighten notation we drop the explicit dependence of $Y_i(A)$ on $A$ and just write $Y_i$.
    Then $\bar Y_{i} = Y_{i} - \E Y_{i}$ is centered subexponential with $\|\bar Y_{i}\|_{\psi_{1}} \le 2 C_{0}$, since centering at most doubles the $\psi_1$-norm: $\|Y - \E Y\|_{\psi_1} \le 2 \|Y\|_{\psi_1}$ because $|EY| \le \|Y\|_{\psi_1}$.
    Therefore, Bernstein's inequality \cite[Theorem 2.8.1]{Vershynin2018} gives
    \[
        \Pp\Bigl(\frac{1}{n} \sum_{i = 1}^n \bar Y_{i} > s\Bigr) \le 2 \exp[-c_{1} n \min(s^2, s)]
    \]
    with $c_{1} = c_{1}(C_{0})$. Since $Y_1, \dots, Y_n$ are i.i.d.\ and $\E Y_{1} \le 1$,
    \[
        \Bigl\{\frac{1}{n} \sum_{i = 1}^n Y_{i} > 1 + s \Bigr\} \subset
        \Bigl\{\frac{1}{n} \sum_{i = 1}^n \bar Y_{i} > s + 1 - \E Y_{1} \Bigr\} \subset
        \Bigl\{\frac{1}{n} \sum_{i = 1}^n \bar Y_{i} > s \Bigr\}.
    \]
    Hence
    \[
        \Pp\left(\frac{1}{n} \sum_{i = 1}^n| \langle Z_i, A\rangle| > 1 + s\right) \le 2 \exp[-c_1 n \min(s^2, s)].
    \]
    
    Let $d_p = p (p + 1) / 2$. A Frobenius $1/2$-net $\mathcal N$ of the unit sphere in $\operatorname{Sym}_p$ can be chosen with
    \[
        |\mathcal N| \le 5^{d_p}.
    \]
    By the assumption of Lemma~\ref{lem:infeasibility_above:dimensionality_bounds}, $n \ge a_0p^2$ and $a_0 \geq 1/4$, so
    \[
        \frac{d_p}{n} \le \frac{p (p + 1)}{2 a_0 p^2} \le \frac{1}{a_0} \le 4,
    \]
    hence there exists a universal constant $c_2 > 0$ and a large enough $s_0$ depending only on $\epsilon$ via $a_0$, such that
    \[
        c_1 \min(s_0^2, s_0) - \frac{d_p}{n} \log(5)
            \ge c_1 \min(s_0^2, s_0) - \frac{\log(5)}{a_0}
            \ge c_2 > 0.
    \]
    The union bound then gives
    \[
        \Pp\left(\max_{A_0 \in \mathcal N} \frac{1}{n} \sum_{i = 1}^n |\langle Z_i, A_0 \rangle| > 1 + s_0 \right)
            \le 2 e^{-c_2 n}
            \le 2 e^{-c_2 a_0 p^2}.
    \]
    For any Frobenius-unit $A$, choose $A_0 = A_0(A) \in \mathcal N$ with $\|A - A_0\|_F \le 1/2$. By homogeneity,
    \begin{align*}
        \Lambda_p
            &\le \sup_{\substack{A \in \operatorname{Sym}_p \\ \|A\|_F = 1}} \frac{1}{n} \sum_{i = 1}^n \Big(|\langle Z_i, A_0 \rangle| + |\langle Z_i, A - A_0 \rangle|\Big) \\
            &\le \max_{A_0' \in \mathcal N} \frac{1}{n} \sum_{i = 1}^n |\langle Z_i, A_0' \rangle| + \frac{1}{2} \Lambda_p,
    \end{align*}
    thus, since $\Lambda_p < \infty$ a.s.,
    \[
        \{\Lambda_p > 2(1 + s_0)\} \subset\Big\{ \max_{A_0 \in \mathcal N} \frac{1}{n} \sum_{i = 1}^n |\langle Z_i, A_0 \rangle| > 1 + s_0 \Big\}.
    \]
    Taking $A_* = 2 (1 + s_0)$, depending only on $\epsilon$ through $s_0$, one gets
    \[
        \Pp(\Lambda_p > A_*) \le 2e^{-c_2 a_0 p^2}.
    \]
\end{proof}



%%%%%%%%%%%%%%%%%%%%
\subsection{Discretization and an approximation to the minimum}\label{deterministic-nets-and-free-energy-inequalities}
%%%%%%%%%%%%%%%%%%%%

Section~\ref{a-robust-gaussian-loss-margin} gives a uniform positive lower bound for the Gaussian model, and \cref{lem:infeasibility_above:average_z_bound} gives that, with high probability, $\Lambda_p \le A_*$. Both are infima/suprema over a continuum with a non-smooth extremum.
This section replaces the continuum by a finite collection and the extremum by a smooth average suitable for interpolation.
We show that both the discretization and smoothing errors can be made small, both a small fraction of the Gaussian margin with fixed constants.

For the moment keep $\tau \in (0,1]$ symbolic, and let
\[
    \mathcal K_\tau = \{D \in {\rm Sym}_p : D \succeq 0, \|D\|_F = 1, \|D\|_{op} \le \tau\}.
\]
If this class is empty, \cref{main-theorem-10-1} is immediate. Assume below that it is nonempty, forcing $\tau \geq 1 / \sqrt{p}$.

Fix $0 < B < \infty$ and let
\[
    \mathcal V = L\cap\{v : \|v\|_1 \le B n\}
\]
be the radius-$(B n)$ ball in the $m_L$-dimensional normed space $(L,\|\cdot\|_1)$.

\begin{lemma} \label{lem:infeasibility_above:nets-entropy}
    For $0 < r \le 1$, $0 < s \le B$ there is a maximal $r$-net $\mathcal N_D \subset {\mathcal K}_\tau$ in Frobenius norm with centers in ${\mathcal K}_\tau$ and an $(s n)$-net $\mathcal N_L \subset \mathcal V$ in $\ell_1$-norm intrinsic to $L$ such that for $\mathcal C = \mathcal N_D \times \mathcal N_L$,
    \begin{align*}
        \frac{\log(|\mathcal C|)}{n} \le R_*,
    \end{align*}
    where
    \[
        R_* = \frac{1}{a_0} \log\left(\frac{3}{r}\right) + \log \Bigl(1 + \frac{2 B}{s}\Bigr).
    \]
\end{lemma}

\begin{proof}
    ${\mathcal K}_\tau$ lies in the Frobenius unit ball of dimension $d_p$.
    A maximal $r$-net has disjoint $r / 2$-balls, volume comparison gives
    \[
        |\mathcal N_D| (r / 2)^{d_p} \le (1 + r / 2)^{d_p},
    \]
    i.e.,
    \[
        \log(|\mathcal N_D|) \le d_p \log(1 + 2 / r) \le d_p \log(3 / r) \qquad \text{for $r \le 1$}.
    \]
    
    $\mathcal V$ is a $(B n)$-ball in $(L , \|\cdot\|_1)$, covered by $(1 + 2 B / s)^{m_L}$ balls of radius $(s n)$, thus
    \[
        \log(|\mathcal N_L|) \le m_L \log\left(1 + \frac{2 B}{s}\right).
    \]
    
    Therefore,
    \begin{align*}
        \log(|\mathcal C|) \le d_p \log\left(\frac{3}{r}\right) + m_L \log\Bigl(1 + \frac{2 B}{s}\Bigr).
    \end{align*}
    
    Since $m_L \leq n$ and \cref{lem:infeasibility_above:dimensionality_bounds} implies that $d_p \le p^2$ and $n \ge a_0 p^2$, the claim follows.
\end{proof}



A candidate $u = (R_u, v_u) \in \mathcal C$ has row shift $c_{iu} = -(v_u)_i$.
For a row ensemble $X=(X_1,\ldots,X_n)$, define the total bounded loss
\[
    H_X(u) = \sum_{i = 1}^n \ell(\langle X_i, R_u \rangle + c_{iu})
\]
and its normalized log-sum-exp approximation
\[
    F_X = -\frac{1}{\beta n} \log\left[\frac{1}{|\mathcal C|} \sum_{u \in \mathcal C} e^{-\beta H_X(u)} \right].
\]
\begin{lemma} \label{lem:infeasibility_above:net_approximation}
    Let $H_* = \min_u H_X(u)$, $L_j = \|\ell^{(j)}\|_\infty$, choose the mesh radii
    \[
        r \le \min\left\{1, \frac{\theta}{64 L_1 A_*}\right\}
        \qquad \text{and} \qquad
        s \le \min\left\{B, \frac{\theta}{64 L_1}\right\},
    \]
    and $\beta \ge (16 R_*) / \theta$.
    Then
    \[
        \frac{H_*}{n} \le F_X \le \frac{H_*}{n} + \frac{\theta}{16}.
    \]
\end{lemma}

\begin{proof}
    We have
    \[
        \frac{e^{-\beta H_*}}{|\mathcal C|} \le \frac{1}{|\mathcal C|} \sum_u e^{-\beta H_X(u)} \le e^{-\beta H_*}.
    \]
    Taking $-(\beta n)^{-1}$ times the logarithm proves the exact inequalities
    \[
        \frac{H_*}{n} \le F_X \le \frac{H_*}{n} + \frac{\log(|\mathcal C|)}{\beta n}.
    \]
    Then
    \[
        \frac{\log|\mathcal C|}{\beta n} \le \frac{\theta}{16},
    \]
    and the claim follows.
\end{proof}

The order is important: first choose the Gaussian margin and $\ell$, then the two fixed mesh radii, and then the fixed $\beta$. The covering estimate is uniform in $\tau$. The telescoping argument in \cref{telescoping-concentration-and-continuum-extension} chooses $\tau$ after $\beta$, and the net is then instantiated inside that chosen class.

%%%%%%%%%%%%%%%%%%%%
\subsection{One conditional Gibbs law shared by both current-row distributions} \label{a-common-leave-one-row-cavity-law}
%%%%%%%%%%%%%%%%%%%%

Let $G_1, \ldots, G_n$ be independent isonormal Gaussian rows, independent of $Z_1, \ldots, Z_n$. Define
\[
    X^{(m)} = (Z_1, \ldots, Z_m, G_{m + 1}, \ldots, G_n) \qquad \text{for $0 \le m \le n$}.
\]
Then $X^{(n)}=Z$, $X^{(0)}=G$, and
\[
    \E[F_Z] - \E[F_G] = \sum_{m = 1}^n \E[F_{X^{(m)}} - F_{X^{(m - 1)}}].
\]
Fix $m$ and set
\[
    \mathcal F_{-m} = \sigma(Z_1, \ldots, Z_{m - 1}, G_{m + 1}, \ldots, G_n).
\]
The rows $Z_m$ and $G_m$ are mutually independent and are both independent of $\mathcal F_{-m}$. Every hybrid row other than the $m$-th is $\mathcal F_{-m}$-measurable.
Define the \emph{row-deleted} Gibbs distribution on the net $\mathcal C = \mathcal N_D\times\mathcal N_L$.
For $u = (R_u, v_u) \in \mathcal C$, $c_{iu} = -(v_u)_i$,
\[
    \nu_{-m}(u) = \frac{e^{-\beta \sum_{i \ne m} \ell(\langle X_i, R_u \rangle + c_{iu})}}{\sum_{v \in \mathcal C} e^{-\beta \sum_{i \ne m} \ell(\langle X_i, R_v \rangle + c_{iv})}}.
\]
Because the sum excludes $i=m$, $\nu_{-m}$ is $\mathcal F_{-m}$-measurable.

The other hybrid rows are identical on both sides of the replacement. Therefore the same conditional law governs both replacements; no Gibbs object depending on either current row is used in the conditioning step.

Let
\[
    H_{-m}(u) = \sum_{i \ne m} \ell(\langle X_i, R_u \rangle + c_{iu}),
\]
where $X_i$ is the common $X_i^{(m)} = X_i^{(m - 1)}$.
Then
\[
    \sum_u e^{-\beta H_X(u)} = \left(\sum_v e^{-\beta H_{-m}(v)}\right) \cdot \sum_u \nu_{-m}(u) e^{-\beta\ell(\langle X_m, R_u \rangle + c_{mu})}
\]

Hence,
\[
    F_X = C_{-m} + \phi_{\nu}(X_m)
\]
where
\[
    C_{-m} = -\frac{1}{\beta n} \log\sum_v e^{-\beta H_{-m}(v)} + \frac{1}{\beta n} \log|\mathcal C|
\]
is $\mathcal F_{-m}$-measurable and independent of $X_m$, and
\[
    \phi_{\nu}(X) = -\frac{1}{\beta n} \log\sum_{u \in \mathcal C} \nu_{-m}(u) e^{-\beta\ell(\langle X, R_u \rangle + c_{mu})}.
\]

$C_{-m}$ cancels in the difference:
\[
    F_{X^{(m)}} - F_{X^{(m - 1)}} = \phi_{\nu}(Z_m) - \phi_{\nu}(G_m).
\]
Therefore,
\[
    \E[F_{X^{(m)}}-F_{X^{(m-1)}}] = \E\big[ \E[\phi_{\nu}(Z_m)-\phi_{\nu}(G_m)\mid\mathcal F_{-m}]\big].
\]

Conditionally on $\mathcal F_{-m}$, $\nu_{-m}$ is deterministic, $Z_m$ and $G_m$ are independent of it and have matching first two moments. The $m$-th summand is now a one-row comparison problem.

%%%%%%%%%%%%%%%%%%%%
\subsection{Exact derivatives and the independent-draw expansion}\label{exact-derivatives-and-the-full-replica-expansion}
%%%%%%%%%%%%%%%%%%%%

For a generic current row $X$, define
\[
    q_u = \langle X, R_u \rangle + c_{mu}.
\]
Let $\mu_X$ be the full Gibbs law obtained by tilting $\nu_{-m}$ with $e^{-\beta\ell(q_u)}$. Relative to the leave-one-row law, it is
\[
    \frac{d\mu_X}{d\nu_{-m}}(u) = \frac{e^{-\beta \ell(q_u)}}{\sum_v \nu_{-m}(v) e^{-\beta \ell(q_v)}}.
\]
Write
\[
    \langle f \rangle_X = \sum_{u \in \mathcal C} f(u) \mu_X(u)
\]
and
\[
    \operatorname{Cov}_X(f, g) = \langle f g \rangle_X - \langle f \rangle_X \langle g \rangle_X.
\]
For a row direction $A \in \operatorname{Sym}_p$, put
\begin{align*}
    a_A(u) &= \ell'(q_u) \langle A, R_u \rangle, \\
    b_{AB}(u) &= \ell''(q_u) \langle A, R_u \rangle \langle B, R_u \rangle, \\
    c_{ABC}(u) &= \ell'''(q_u) \langle A, R_u \rangle \langle B, R_u \rangle \langle C, R_u \rangle.
\end{align*}
For every row-dependent observable $f$,
\[
    D\langle f\rangle_X[B] = \langle D f[B] \rangle_X - \beta\operatorname{Cov}_X(f, a_B).
\]
The first term is differentiation inside, the second is differentiation of the Gibbs weight itself.
Differentiating gives
\begin{align*}
  D\phi_\nu(X)[A] &= \frac{1}{n} \langle a_A \rangle_X, \\
  D^2\phi_\nu(X)[A,B] &= \frac{1}{n} \left( \langle b_{AB} \rangle_X - \beta[\langle a_A a_B \rangle_X - \langle a_A \rangle_X \langle a_B \rangle_X] \right).
\end{align*}
Let $\sum_{cyc}$ denote the three cyclic choices of $(A, B, C)$. Then
\[
  n D^3\phi_\nu(X)[A, B, C] = \langle c_{ABC} \rangle_X - \beta\sum_{cyc} \operatorname{Cov}_X(b_{A B}, a_C) + \beta^2\kappa_{3, X}(a_A, a_B, a_C),
\]
where
\[
    \kappa_{3,X}(a, b, c)
        = \langle abc \rangle_X
        - \langle ab \rangle_X \langle c \rangle_X
        - \langle ac \rangle_X \langle b \rangle_X
        - \langle bc \rangle_X \langle a \rangle_X
        + 2 \langle a \rangle_X \langle b \rangle_X \langle c \rangle_X.
\]
The $\ell'''$ group is the one-independent-draw from the same Gibbs probability measure term
\[
    \langle c_{ABC}\rangle_X.
\]
For each of the three cyclic choices, the $\ell'' \ell'$ group is
\[
  \operatorname{Cov}_X(b_{A B}, a_C) = \langle b_{A B} a_C \rangle_X -\langle b_{A B} \rangle_X \langle a_C \rangle_X,
\]
which is respectively a one-independent-draw term and a two-independent-draw term.
The $(\ell')^3$ group has one one-independent-draw term, $\langle a_A a_B a_C\rangle_X$, three two-independent-draw terms
\[
    -\langle a_A a_B \rangle_X \langle a_C \rangle_X,
    -\langle a_A a_C \rangle_X \langle a_B \rangle_X,
    -\langle a_B a_C \rangle_X \langle a_A \rangle_X,
\]
and the three-independent-draw term
\[
    2 \langle a_A \rangle_X \langle a_B \rangle_X \langle a_C \rangle_X.
\]


Since $0 \le \ell \le 1$ by \eqref{eq:infeasibility_above:smoothing},
\[
    e^{-\beta} \le \frac{d\mu_X}{d\nu_{-m}} \le e^\beta.
\]
For product measures (one, two, and three independent draws) the exact product bounds are:
\[
    e^{-k \beta} \le \frac{d\mu_X^{\otimes k}}{d\nu_{-m}^{\otimes k}} \le e^{k \beta} \qquad \text{for $k = 1, 2, 3$}.
\]

If we replace an expectation under $\mu_X$ by one under $\nu_{-m}$ we pay at most $e^{k\beta}$.
Only after these replacements do we take absolute values and use
\[
    |\ell'|\le L_1, |\ell''|\le L_2, |\ell'''|\le L_3.
\]
After this change of measure, independent samples are drawn from $\nu_{-m}$, which is ${\mathcal F}_{-m}$-measurable and independent of both $Z_m$ and $G_m$. Its $R_u$'s are fixed given ${\mathcal F}_{-m}$.

Then the scalar part of $D^3$ can be bounded purely by constants. The $R$-dependent part $\langle A,R_{u_1} \rangle \langle B, R_{u_2} \rangle \langle C, R_{u_3}\rangle$ remains, but now $u_i$ are independent of $X$.

Counting absolute values: $Cov$ gives 2 terms per cyclic, $\kappa_3$ gives $1+1+1+1+2=6$ terms in absolute value. Hence the total absolute scalar mass is at most
\[
    \frac{1}{n} \big[L_3 + 6 \beta L_1 L_2 + 6 \beta^2 L_1^3\big].
\]


%%%%%%%%%%%%%%%%%%%%
\subsection{An exact integration-by-parts identity for centered rank-one rows}\label{exact-second-chaos-identity-and-interpolation}
%%%%%%%%%%%%%%%%%%%%

The next identity is the only point where the non-Gaussian distribution of $Z=(ww^\top - I_p)/\sqrt2$ enters.
We derive it in coordinates rather than invoking a named interpolation principle. The Hessian term matches the Gaussian row exactly; the remaining error contains the third derivative and is small when all candidate matrices have small operator norm.

For $1\le a, b \le p$, define
\[
    S_{ab} = \frac{e_a e_b^\top + e_b e_a^\top}{2}
\]
and
\[
    B_a(w) = \frac{e_a w^\top + w e_a^\top}{\sqrt{2}}.
\]
Then
\[
    \frac{\partial Z}{\partial w_a} = B_a(w)
\]
and
\[
    \frac{\partial^2 Z}{\partial w_a \partial w_b} = \sqrt{2} S_{ab}.
\]
Let $M \in \operatorname{Sym}_p$ be independent of $w$, let $s\ge0$ be deterministic, and set $Y = M + s Z$.
Consider $\Phi\in C^3(\operatorname{Sym}_p)$ with bounded derivatives, and identify $D\Phi$ with the symmetric gradient and $D^2\Phi$, $D^3\Phi$ with the symmetric bi/tri-linear forms.

Let
\[
    f_{ab}(w) = \langle D\Phi(Y), S_{ab} \rangle.
\]
Note that since $D^j\Phi$ is bounded for $j = 1, 2, 3$, it follows that $f_{ab}$ is bounded, and $\partial_a\partial_b f_{ab}$ is at most quadratic in $w$.
Since $M$ is independent of $w$, the requirements of conditional Stein's lemma, given $M$, are met; and by the law of total expectation,
\[
    \E[(w_a w_b - \delta_{ab}) f_{ab}(w)] = \E[\partial_a \partial_b f_{ab}(w)].
\]

Since
\[
  \langle Z, D\Phi(Y) \rangle = \frac{1}{\sqrt{2}} \sum_{a,b} (w_a w_b - \delta_{ab}) \langle D\Phi(Y), S_{ab} \rangle,
\]
the chain rule gives, before summing over $a, b$,
\[
    \partial_a\partial_b f_{ab}(w) = s^2 D^3\Phi(Y)[B_a(w), B_b(w), S_{ab}] + s \sqrt{2} D^2\Phi(Y)[S_{ab}, S_{ab}].
\]
Indeed, $\partial_a Y = s B_a(w)$ and $\partial_a B_b(w) = \sqrt{2} S_{ab}$. Multiplication by $1/\sqrt{2}$ and summation over the ordered pairs $(a, b)$ therefore gives
\[
    \E\langle Z, D\Phi(M + s Z) \rangle = s \, \E\Delta\Phi(M + s Z) + s^2 \, \E\mathfrak T_w(D^3\Phi(M + sZ)),
\]
where
\[
    \Delta\Phi(Y) = \operatorname{Tr} D^2 \Phi(Y) = \sum_{a \leq b} D^2\Phi(Y)[Q_{ab}, Q_{ab}]
\]
and
\[
    \mathfrak T_w(T) = \frac{1}{\sqrt{2}} \sum_{a, b = 1}^p T[B_a(w), B_b(w), S_{ab}].
\]
The Hessian term is exactly the Frobenius Laplacian.
Indeed, $S_{aa} = Q_{aa}$, while $S_{ab} = Q_{ab}/\sqrt{2}$ for $a \ne b$, and the two ordered pairs $(a, b)$ and $(b, a)$ give one copy of $D^2\Phi[Q_{ab}, Q_{ab}]$.

For matrices $A,B,C \in \operatorname{Sym}_p$, the ordered pure tensor $T = A \otimes B \otimes C$ is the trilinear form
\[
    T[U,V,W] = \langle A, U \rangle \langle B, V \rangle \langle C, W \rangle.
\]
For this tensor,
\[
    \mathfrak T_w(A \otimes B \otimes C) = \sqrt{2} w^\top A C B w.
\]
This follows from
\[
    \langle A, B_a(w) \rangle = \sqrt{2}(A w)_a
\]
and $\langle C, S_{ab} \rangle = C_{ab}$.



For $\Phi=\phi_\nu$, finiteness of the candidate family and boundedness of the first three derivatives of $\ell$ give deterministic bounds on $D^j\Phi$ for $1 \le j \le 3$.
The Gaussian variables appearing below have finite moments of every order. Hence, on every interval $[\delta, 1 - \delta]$, the difference quotients and
integration-by-parts integrands admit an integrable dominating function, which justifies differentiating under the expectation. The bound obtained below is uniform in $\delta$ and justifies passage to the endpoints by dominated convergence.

Now set
\[
    Y_t = \sqrt{t} \, G + \sqrt{1 - t} \, Z.
\]
Gaussian integration by parts gives
\[
    \E\langle G, D\Phi(Y_t) \rangle = \sqrt{t} \, \E\Delta\Phi(Y_t).
\]
Using the second-chaos identity with $s = \sqrt{1 - t}$ yields
\[
    \frac d{dt}\E\Phi(Y_t) = -\frac{\sqrt{1 - t}}{2} \, \E\mathfrak T_w(D^3\Phi(Y_t)).
\]
The right side is uniformly bounded and tends to zero as $t \uparrow 1$. The identity can first be integrated on $[\delta, 1 - \delta]$ and then  $\delta \downarrow 0$. Thus both endpoints are integrable, and no fourth derivative is used.



%%%%%%%%%%%%%%%%%%%%
\subsection{Bounding the third-derivative error for one replaced row}\label{fixed-triple-contraction-and-the-one-row-bound}
%%%%%%%%%%%%%%%%%%%%

Let $R_1, R_2, R_3$ be symmetric and satisfy
\[
    \|R_j\|_F \le 1, \|R_j\|_{op} \le \tau \le 1.
\]
For every ordering,
\[
    |\operatorname{Tr}(R_1 R_2 R_3)| \le \|R_1\|_F \|R_2 R_3\|_F \le \tau
\]
and
\[
  \|R_1 R_2 R_3\|_F \le \|R_1\|_{op} \|R_2\|_{op} \|R_3\|_F \le \tau^2.
\]
The symmetric part $Q$ of the ordered product obeys the same bounds. Since
\[
    \E|w^\top Q w| \le |\operatorname{Tr}Q| + [\E|w^\top Q w - \operatorname{Tr}Q|^2]^{1/2},
\]
and
\[
    \E|w^\top Q w - \operatorname{Tr}Q|^2 = 2\|Q\|_F^2,
\]
we obtain
\[
    \E|w^\top Q w| \le \tau + \sqrt{2} \tau^2 \le (1 + \sqrt{2}) \tau.
\]
Combining this with the termwise change-of-measure bounds and the tensor contraction gives a universal $C_2$ such that
\[
    \E|\mathfrak T_w(D^3\phi_\nu(Y_t))| \le \frac{C_2 e^{3\beta}\tau}{n} [L_3 + \beta L_1 L_2 + \beta^2 L_1^3].
\]
Integrating the interpolation and absorbing the constant factors into $C_2$ gives the normalized one-row estimate
\[
    |\E\phi_\nu(Z) - \E\phi_\nu(G)| \le \frac{C_2 e^{3\beta} \tau}{n} [L_3 + \beta L_1 L_2 + \beta^2 L_1^3].
\]
The factor $1/n$ comes from the free-energy normalization. The estimate is uniform in the leave-one-row law, the finite family, the shifts, $p$, and $n$.

%%%%%%%%%%%%%%%%%%%%
\subsection{Telescoping, concentration, and continuum extension}\label{telescoping-concentration-and-continuum-extension}
%%%%%%%%%%%%%%%%%%%%

Summing the $n$ hybrid-row estimates yields
\[
    |\E F_Z - \E F_G| \le C_2 e^{3\beta} \tau [L_3 + \beta L_1 L_2 + \beta^2 L_1^3].
\]
Choose $\tau \le 1$ after $\beta$ so that the right side is at most $\theta / 16$. Then for $p \ge \tau^{-2}$, $1/\sqrt p \le \tau$ and ${\mathcal K}_\tau \neq \emptyset$. So a $p$-independent $\tau$ works for all sufficiently large $p$.

On the Gaussian robust event, every candidate has $H_G(u) / n \ge \theta$. Since $F_G \ge \min_u H_G(u) / n$ and $F_G \ge 0$ always,
\[
    \E F_G \ge \theta[1 - e^{-\gamma^2 p^2 / 4}].
\]
For all sufficiently large $p$, this is at least $15 \theta / 16$. Hence
\[
    \E F_Z \ge \frac{7 \theta}{8}.
\]
Replacing one row changes every $H_X(u)$ by at most one.
The log-sum-exp map therefore changes $F_X$ by at most $1 / n$. The bounded-differences inequality \cite{McDiarmid1989} gives
\[
    \Pp(|F_Z - \E F_Z| > t) \le 2 e^{-2 n t^2}.
\]
With $t = \theta / 8$,
\[
    \Pp\left(F_Z < \frac{3 \theta}{4}\right) \le 2 e^{-n \theta^2 / 32}.
\]
The free-energy inequality and the entropy choice imply, on the complement of this event,
\[
    \min_{u \in \mathcal C} \frac{H_Z(u)} n \ge F_Z - \frac{\log|\mathcal C|}{\beta n} \ge \frac{11 \theta}{16}.
\]
Now take any $D \in \mathcal K_\tau$ and $v\in L$ with $\|v\|_1 \le B n$. Choose net points $R$ and $v_0$ satisfying
\[
    \|D - R\|_F \le r, \|v - v_0\|_1 \le s n.
\]
If $D = R$ set $A = 0$, otherwise put $A = (D - R) / \|D - R\|_F$ so $\|A\|_F = 1$ and $D - R = \|D - R\|_F A$. Then it follows from the definition of $\Lambda_p$ that
\[
    \frac{1}{n} \sum_{i = 1}^n |\langle Z_i, D - R \rangle| = \|D - R\|_F \frac{1}{n} \sum_{i = 1}^n |\langle Z_i, A\rangle| \le \|D - R\|_F \Lambda_p \le r \Lambda_p
\]
on $\{\Lambda_p \le A_*\}$, with the convention $0 \cdot \Lambda_p = 0$. In particular on $\{\Lambda_p \le A_*\}$,
\[
    \frac{1}{n} \sum_{i = 1}^n \big(|\langle Z_i, D - R \rangle| + |v_i - (v_0)_i|\big) \le \Lambda_p r + s \le A_*r + s.
\]
Therefore the Lipschitz property of $\ell$ gives
\[
    \frac{1}{n} \sum_{i = 1}^n |\ell(\langle Z_i, D \rangle - v_i) - \ell(\langle Z_i, R \rangle - (v_0)_i)| \le L_1(A_* r + s) \le \frac{\theta}{32},
\]
hence
\[
    \inf_{\substack{D \in \mathcal K_\tau \\ v \in L, \|v\|_1 \le B n}} \frac{1}{n}\sum_{i = 1}^n \ell(\langle Z_i, D \rangle - v_i) \ge \frac{\theta}{2}.
\]
Since $\ell(t) \le C_\ell |t|$, the conclusion of \cref{main-theorem-10-1} holds with
\[
    c = \frac\theta{2 C_\ell}.
\]
For large $p$, the total failure probability is at most
\[
    2 e^{-n \theta^2 / 32} + 2 e^{-c_2' n} \le e^{-c'p^2}
\]
after decreasing $c' > 0$, because $n \ge a_0 p^2$ and $n \le C_1 p^2$.

This proves \cref{main-theorem-10-1}.
% TODO: remove reference to future \cref{lem:cavity-comparison} or re-structure.

%%%%%%%%%%%%%%%%%%%%
\subsection{Conditional nuisance-space reduction}\label{conditional-nuisance-spaces-in-the-peeling-argument}
%%%%%%%%%%%%%%%%%%%%

Fix a deterministic subspace $E_0\subseteq\mathbb R^d$ of
dimension $j\le K$, and split each sample as
\[
  x_i=g_i+w_i\in E_0\oplus E_0^\perp.
\]

Suppose that
\[
S_0=H_0+\rho D_0,
\]
where $H_0$ is supported on $E_0$, $D_0$ is supported on
$E_0^\perp$,
\[
D_0\succeq0,\qquad
\|D_0\|_F=1,\qquad
\|D_0\|_{op}\le\tau,
\]
and $\rho\ge\delta$.


Let
\[
  \mathcal G=\sigma(g_1,\ldots,g_n).
\]
Conditional on $\mathcal G$, the vectors $w_i$ are still
independent $N(0,I_p)$ vectors with $p=d-j$. Define
\[
  \mathcal Q_w(A)_i=\frac{w_i^TAw_i-\operatorname{Tr}A}
  {\sqrt2} =\langle Z_i,A\rangle.
\]
For $A\in\operatorname{Sym}(E_0)$, also define
\[
  \mathcal Q_g(A)_i=\frac{g_i^TAg_i-\operatorname{Tr}A}
  {\sqrt2}.
\]
Now define
\[
  L(E_0,g)=\operatorname{span}\{\1_n\}+\{\mathcal Q_g(A):
  A\in\operatorname{Sym}(E_0)\}.
\]
This is a $\mathcal G$-measurable subspace. After
conditioning, it is a deterministic subspace independent
of all tail rows, and
\[
  \dim L(E_0,g)\le1+\frac{j(j+1)}2\le1+\frac{K(K+1)}2.
\]
For every realization of $g$ for which the dimensional
inequality holds, define
\[
  \mathcal E_{\rm str}(g)=\left\{ \inf_{\substack{D\succeq0,\
  \|D\|_F=1,\ \|D\|_{op}\le\tau\\ v\in L(E_0,g),\
  \|v\|_1\le Bn}} \frac1n\|\mathcal Q_w(D)-v\|_1\ge c
  \right\}.
\]
Then \cref{main-theorem-10-1} and the empirical
continuity estimate give
\[
  \Pp\left((\mathcal E_{\rm str}(g)\cap\{\Lambda_w\le
  A_*\})^c\mid\mathcal G\right)\le3e^{-c'p^2},
\]
where
\[
  \Lambda_w=\sup_{\|A\|_F=1}\frac1n\|\mathcal Q_w(A)\|_1.
\]
The bound is independent of the orientation of the
realized nuisance subspace. Integrating over $g$
preserves the same bound. Since the event in
\cref{main-theorem-10-1} takes an infimum over the full
classes of $D$ and $v$, both may be measurable functions
of the tail rows. This is why the nuisance vector may be
selected after observing the tail.


The proof of \cref{main-theorem-10-1} also gives the
joint tail-continuity event
\[
  \frac1n\|\mathcal Q_w(A)\|_1\le A_*\|A\|_F
\]
for every symmetric $A$, with conditional failure probability at most
$e^{-cp^2}$. In particular,
\[
  \frac1n\|\mathcal Q_w(D_0)\|_1\le A_*.
\]


For the decomposition $S_0=H_0+\rho D_0$, define
\[
  \bar a
  =
  \frac1{\rho n}\1_n^T\mathcal Q_d(S_0),
  \qquad
  v
  =
  -\mathcal Q_g(H_0/\rho)+\bar a\1_n.
\]
Then
\[
  \mathcal Q_w(D_0)-v
  =
  \rho^{-1}\PiZero\mathcal Q_d(S_0).
\]

Assume $B>A_*+1$, and decrease $c$ if necessary so that
$0<c\le1$. On
\[
  \mathcal E_{\rm str}(g)\cap\{\Lambda_w\le A_*\},
\]
if
\[
  \frac1n\|\PiZero\mathcal Q_d(S_0)\|_1<\rho c,
\]
then
\[
  \frac1n\|\mathcal Q_w(D_0)-v\|_1<c.
\]
Moreover,
\[
  \frac1n\|v\|_1
  \le
  \frac1n\|\mathcal Q_w(D_0)\|_1
  +
  \frac1n\|\mathcal Q_w(D_0)-v\|_1
  <
  A_*+1
  <
  B.
\]
This contradicts the definition of
$\mathcal E_{\rm str}(g)$. Therefore, on the same good
event,
\[
  \frac1n\|\PiZero\mathcal Q_d(S_0)\|_1
  \ge
  \rho c
  \ge
  \delta c.
\]

Thus the conditional diffuse margin holds uniformly for
every deterministic $E_0$ of dimension at most $K$, with
conditional failure probability at most $3e^{-c'p^2}$. 

\end{aiblock}


%%%%%%%%%%%%%%%%%%%%
\section{The comparison estimate extracted from the
  preceding
  calculation}\label{the-comparison-estimate-extracted-from-the-preceding-calculation}
%%%%%%%%%%%%%%%%%%%%

\begin{aiblock}
 
\begin{lemma}[Finite-family leave-one-row
  comparison]\label{lem:cavity-comparison}
  Let $Z=(ww^T-I_p)/\sqrt2$, where $w\sim N(0,I_p)$, and
  let $G$ be the isonormal Gaussian vector on $\Sym_p$.
  Let $\mathcal C$ be a nonempty finite set. For each
  $u\in\mathcal C$, let $R_u\in\Sym_p$ be deterministic with
  $\|R_u\|_F\le1$ and $\|R_u\|_{op}\le\tau$, and let
  $c_{iu}$ be deterministic. 
  For every fixed $\beta>0$, if
  $0\le\ell\le1$ is $C^3$ and
  \[
    F_X=-\frac1{\beta n}\log\left[
    \frac1{|\mathcal C|}\sum_{u\in\mathcal C}
    \exp\left(-\beta\sum_{i=1}^n
    \ell(\langle X_i,R_u\rangle+c_{iu})\right)
    \right],
  \]
  then
  \[
    |\E F_Z-\E F_G| \le Ce^{3\beta}\tau \left(L_3+\beta
    L_1L_2+\beta^2L_1^3\right),
  \]
  where $L_j=\|\ell^{(j)}\|_\infty$. The constant is
  uniform in $p$, $n$, $\mathcal C$, and the shifts.
\end{lemma}

\begin{proof}
  This is exactly the finite-family comparison established in
  \cref{expert-leave-one-row-module}, especially the coordinate
  Stein identity, the leave-one-row Gibbs construction, the
  derivative expansion, and the telescoping row replacement.
\end{proof}

\end{aiblock}


%%%%%%%%%%%%%%%%%%%%
\section{Applying the comparison uniformly over diffuse
  matrices and nuisance
  vectors}\label{applying-the-comparison-uniformly-over-diffuse-matrices-and-nuisance-vectors}
%%%%%%%%%%%%%%%%%%%%

\begin{aiblock}
 
\begin{theorem}\label{thm:diffuse-restatement}
  Fix $\epsilon,B,C_1>0$. There are $\tau,c,c'>0$ such that
  the following holds. Let $L\subset\R^n$ be
  deterministic, with dimension $m_L$, and assume
  \[
    n-m_L\ge\left(\frac14+\epsilon\right)p^2, \qquad
    n\le C_1p^2.
  \]
  Let $w_1,\ldots,w_n$ be independent $N(0,I_p)$ vectors,
  and put
  \[
    Z_i=\frac{w_iw_i^T-I_p}{\sqrt2}, \qquad 1\leq i\leq n.
  \]
  Then, for all sufficiently large $p$,
  \[
    \Pp\left[ \inf_{\substack{D\succeq0,\\ \|D\|_F=1,\\
    \|D\|_{op}\le\tau\\ v\in L,\\ \|v\|_1\le Bn}}
    \frac1n\sum_i|\langle Z_i,D\rangle-v_i|<c \right]\le
    e^{-c'p^2}.
  \]
% Maybe adding a human input defining the set M=\{D\succeq
  The constants are uniform over the orientation of
  deterministic $L$.
\end{theorem}

\begin{proof}
  The Gaussian robust-loss margin, finite-family
  comparison, leave-one-row construction, exact
  third-derivative expansion, continuum extension, and
  conditional nuisance-space form are proved in
  \cref{expert-leave-one-row-module}.
\end{proof}

\end{aiblock}



%%%%%%%%%%%%%%%%%%%%
\section{From diffuse matrices to the full positive
  semidefinite
  cone}\label{fixed-rank-peeling-and-unsat-above-one-quarter}
%%%%%%%%%%%%%%%%%%%%

\begin{aiblock}
 
Peeling Lemma~\ref{main-lemma-11-2}: $S=H+R$, $\rank H\le K$:
either $\|R\|_F\le\delta$ (low-rank by $c_*$)
or $S=H+\rho D$ $\tau$-diffuse (diffuse by
Grassmann net and Theorem~\ref{thm:diffuse-restatement}).

We first record a global empirical-continuity
estimate.
For every fixed $0<c_{\rm asp}\le C_{\rm asp}<\infty$,
define, whenever $c_{\rm asp}d^2\le n\le C_{\rm asp}d^2$,
\[
  \Lambda_d=\sup_{\|A\|_F=1}\frac1n\|\mathcal Q_d(A)\|_1.
\]
There are constants $C_{\rm cont}<\infty$ and
$c_{\rm cont}>0$, depending only on $c_{\rm asp},C_{\rm asp}$,
such that
\[
  \Pp(\Lambda_d>C_{\rm cont})\le e^{-c_{\rm cont}n}.
\]
The fixed-net argument in
\cref{expert-leave-one-row-module} proves this estimate.

\begin{lemma}[Universal fixed-rank lower
  bound]\label{main-lemma-11-1}
  Fix $0<c_{\rm asp}\le C_{\rm asp}<\infty$. There is a
  numerical $c_*>0$ such that, for every fixed $K$, there
  is $c_K>0$ for which, whenever $c_{\rm asp}d^2\le n\le
  C_{\rm asp}d^2$,
  \[
    \Pp\left( \inf_{\substack{\operatorname{rank}A\le
    K\\\|A\|_F=1}} \frac1n\|\PiZero\mathcal{Q}_d(A)\|_1<c_*
    \right)\le e^{-c_Kd^2}.
  \]
  Only the starting dimension may depend on fixed $K$.
\end{lemma}

\begin{proof}
  Pair consecutive observations and put
  \[
    Y_j(A)=\mathcal{Q}_d(A)_{2j-1}-\mathcal{Q}_d(A)_{2j}.
  \]
  For $\|A\|_F=1$,
  \[
    \E Y_j(A)^2=2.
  \]
  The Gaussian quadratic fourth-moment formula gives
  \[
    \E\mathcal{Q}_d(A)_1^4=3+12\operatorname{Tr}(A^4)
    \le15,
  \]
  and hence $\E Y_j(A)^4\le36$. Paley-Zygmund and Chernoff
  give, for fixed $A$,
  \[
    \Pp\left(\frac1m\sum_{j=1}^m|Y_j(A)|<\frac1{72}\right)
    \le e^{-cm}, m=\lfloor n/2\rfloor.
  \]
  Here is the covering estimate used to make the preceding
  bound uniform. For $1\le r\le K$, write
  \[
    A=U\operatorname{diag}(\lambda)U^T,
  \]
  where $U\in\operatorname{St}(d,r)$ and
  $\|\lambda\|_2=1$. A maximal $\varepsilon/4$-separated
  subset of the Stiefel manifold in Frobenius distance has
  cardinality at most
  \[
    \left(1+\frac{8\sqrt r}{\varepsilon}\right)^{dr},
  \]
  by the usual disjoint-ball volume comparison in
  $\mathbb R^{dr}$, and the unit sphere in $\mathbb R^r$
  has an $\varepsilon/4$-net of cardinality at most
  $(1+8/\varepsilon)^r$. For two such parameter pairs,
  \[
    \|U\operatorname{diag}(\lambda)U^T
    -V\operatorname{diag}(\mu)V^T\|_F \le2\|U-V\|_F+
    \|\lambda-\mu\|_2.
  \]
  Consequently the Frobenius unit sphere of rank-at-most-
  $K$ symmetric matrices has an $\varepsilon$-net of
  cardinality
  \[
    \sum_{r=1}^K\left(\frac{C_K}{\varepsilon}\right)^{dr+r}.
  \]
  At fixed resolution its entropy is $O_K(d)$, whereas the
  fixed-object failure exponent is $\Theta(d^2)$. On the
  continuity event the net extension costs $2C_{\rm cont}
  \|A-A_0\|_F$. Finally,
  \[
    \|\PiZero\mathcal{Q}_d(A)\|_1\ge\sum_{j=1}^m|Y_j(A)|.
  \]
  Choosing a numerical resolution gives a numerical $c_*$
  independent of fixed $K$.
\end{proof}

\begin{lemma}[Finite decomposition into large eigenvalues
  and a diffuse
  remainder]\label{main-lemma-11-2}
  For fixed $0<\tau,\delta<1$, choose
  \[
    K=\left\lceil\frac{2\log(1/\delta)}{\tau^2}\right\rceil+1.
  \]
  Every PSD Frobenius-unit $S$ has one of two
  decompositions.

  The low-rank branch is
  \[
    S=H+R, H,R\succeq0, HR=0, \operatorname{rank}H\le K,
    \|R\|_F\le\delta.
  \]
  The diffuse branch is
  \[
    S=H+\rho D, H,D\succeq0, HD=0, \operatorname{rank}H\le
    K,
  \]
  \[
    \rho>\delta, \|D\|_F=1, \|D\|_{op}\le\tau.
  \]
\end{lemma}

\begin{proof}
  Start with $R_0=S$. If $\|R_j\|_F\le\delta$, stop in
  the first branch. If $\|R_j\|_{op}\le\tau\|R_j\|_F$,
  stop in the second. Otherwise remove a largest
  eigenvalue and put it in $H$. Then
  \[
    \|R_{j+1}\|_F^2<(1-\tau^2)\|R_j\|_F^2.
  \]
  After $K$ steps the tail is at most $\delta$, so the
  alternatives are exhaustive.
\end{proof}

\paragraph{Replacing the random leading eigenspace by a
  deterministic net.} The leading eigenspace of $S$ is
random. To make the conditioning in
\cref{conditional-nuisance-spaces-in-the-peeling-argument}
uniform we replace $\operatorname{Gr}(j,d)$ by a
deterministic net and rotate the approximant to the true
subspace with a near-identity orthogonal map.

\begin{lemma}[Grassmann $\zeta$-net and near-identity
  rotation]\label{lem:grassmann-net}
  Let $0<\zeta<1/2$ and equip $\operatorname{Gr}(j,d)$
  with $d_{\rm Gr}(E,F)=\|P_E-P_F\|_{op}$, where $P_E$ is
  the orthogonal projector onto $E$. For $j\ge1$,
  $\operatorname{Gr}(j,d)$ has a $\zeta$-net of
  cardinality at most
  \[
    \left(1+\frac{4\sqrt j}{\zeta}\right)^{dj},
  \]
  and for $j=0$ the singleton is a net. If
  $d_{\rm Gr}(E,F)\le\zeta$, there exists $O\in O(d)$
  with
  \[
    OE=F\quad\text{and}\quad\|O-I\|_{op}\le\sqrt2\,\zeta.
  \]
\end{lemma}

\begin{proof}
  Choose a maximal $\zeta/2$-separated set of orthonormal
  $j$-frames in Frobenius distance. Disjoint Frobenius
  balls of radius $\zeta/4$ around these frames lie in the
  Euclidean ball of radius $\sqrt j+\zeta/4$ in
  $\R^{dj}$, which proves the displayed cardinality.
  Maximality makes the set a $\zeta/2$-net, and
  \[
    \|UU^{T}-VV^{T}\|_{op}\le2\|U-V\|_{F}
  \]
  shows that the corresponding subspaces form a $\zeta$-
  net.

  For the last claim, rotate the principal two-planes from
  $E$ to $F$ and act as the identity on the remaining
  directions. If $\theta_{\max}$ is the largest principal
  angle, then
  \[
    \sin\theta_{\max}=\|P_E-P_F\|_{op},\qquad
    \|O-I\|_{op}=2\sin(\theta_{\max}/2) \le\sqrt2\sin
    \theta_{\max},
  \]
  since $\zeta<1/2$ implies $\theta_{\max}<\pi/2$ and
  $\cos(\theta_{\max}/2)\ge1/\sqrt2$.
\end{proof}

\begin{theorem}[Uniform conic lower
  bound]\label{main-theorem-11-3}
  For fixed $1/4<\alpha_-\le\alpha_+<\infty$, there are
  $c,C>0$ such that, for all sufficiently large $d$,
  whenever
  \[
    \alpha_-d^2\le n\le\alpha_+d^2,
  \]
  one has
  \[
    \Pp\left[ \inf_{\substack{S\succeq0\\\|S\|_F=1}}
    \frac1n\|\PiZero\mathcal{Q}_d(S)\|_1<c \right]\le
    e^{-Cd^2}.
  \]
\end{theorem}

\begin{proof}
  Set
  \[
    \epsilon_0=\frac12\left(\alpha_--\frac14\right)\text{
    and }C_1=4\alpha_+.
  \]
  Let $C_{\rm cont}$ be the global continuity constant for
  the aspect interval $[\alpha_-,\alpha_+]$, let $c_*$ be
  the fixed-rank constant, and let $A_*=A_*(\epsilon_0,C_1)$
  be the tail-continuity constant proved in
  \cref{expert-leave-one-row-module}. Choose
  $0<\delta<1/2$ so that
  \[
    g_{\rm lr}=c_*\sqrt{1-\delta^2}-2C_{\rm cont}\delta>0,
  \]
  and set $B=A_*+2$. Apply \cref{thm:diffuse-restatement}
  with parameters $(\epsilon_0,B,C_1)$, and denote its
  margin and probability constants by $c_{\rm str}$ and
  $c'_{\rm str}$. Let $\tau$ be its operator-norm cap, and
  then choose $K=K(\tau,\delta)$ by
  \cref{main-lemma-11-2}.

  By the preceding Grassmann cover and rotation statement,
  the rotation below satisfies $\|O-I\|_{op}\le C_O\zeta$
  with $C_O=\sqrt2$. Put $C_{\rm rot}=2C_OC_{\rm cont}$
  and choose a fixed $\zeta>0$ such that
  \[
    C_{\rm rot}\frac\zeta\delta\le\min\left\{
    \frac{c_{\rm str}}4,\frac12\right\}.
  \]
  Finally choose
  \[
    0<c_{\rm cone}<\min\left\{\frac{g_{\rm lr}}2,
    \frac{\delta c_{\rm str}}4,\frac\delta2\right\}.
  \]
  In the low-rank branch, if $n^{-1}\|\PiZero\mathcal{Q}_d
  (S)\|_1<c_{\rm cone}$, then
  \[
    c_*\|H\|_F \le\frac1n\|\PiZero\mathcal{Q}_d(H)\|_1
    \le\frac1n\|\PiZero\mathcal{Q}_d(S)\|_1+2C_{\rm cont}
    \|R\|_F.
  \]
  Since $\|H\|_F\ge\sqrt{1-\delta^2}$, this gives
  \[
    c_*\sqrt{1-\delta^2}<c_{\rm cone}+2C_{\rm cont}\delta<
    c_*\sqrt{1-\delta^2},
  \]
  a contradiction.

  In the diffuse branch, let $j=\operatorname{rank}H$ and
  $E=\operatorname{range}H$. For each $0\le j\le K$, fix
  a deterministic $\zeta$-net of $\operatorname{Gr}(j,d)$
  before observing the samples. Choose the nearest point
  $E_0$ and an orthogonal $O$ such that
  \[
    OE=E_0\text{ and }\|O-I\|_{op}\le C_O\zeta.
  \]
  Put
  \[
    S_0=OSO^T=H_0+\rho D_0.
  \]
  Then $H_0$ is supported on $E_0$, $D_0$ on $E_0^\perp$,
  and the canonical rotation obeys the following sharper
  PSD estimate. If $S=\sum_j\lambda_jv_jv_j^T$ with
  $\lambda_j\ge0$, then
  \[
    \operatorname{Tr}(SOSO^T) =\sum_{j,k}\lambda_j\lambda_k
    (v_j^TOv_k)^2 \ge\cos^2(\theta_{\max})\sum_j\lambda_j^2,
  \]
  because the symmetric part of the canonical
  principal-plane rotation is at least
  $\cos(\theta_{\max})I$. Consequently,
  \[
    \|OSO^T-S\|_F^2 =2\operatorname{Tr}(S^2)-2\operatorname
    {Tr}(SOSO^T) \le2\sin^2(\theta_{\max})\|S\|_F^2.
  \]
  Since $\|S\|_F=1$ and $\sin\theta_{\max}\le\zeta$, this
  gives
  \[
    \|S_0-S\|_F\le C_O\zeta\text{ and }\frac{\|S_0-S\|_F}
    {\rho}\le C_O\frac\zeta\delta.
  \]
  For fixed $E_0$, write $x_i=g_i+w_i$ along
  $E_0\oplus E_0^\perp$ and define
  \[
    \mathcal G_{E_0}=\sigma(g_1,\ldots,g_n).
  \]
  Conditional on $\mathcal G_{E_0}$, the $w_i$ are
  independent $N(0,I_p)$ vectors with $p=d-j$. For
  $A\in\operatorname{Sym}(E_0)$ and
  $D\in\operatorname{Sym}(E_0^\perp)$, define
  \[
    \mathcal Q_g(A)_i=\frac{g_i^TAg_i-\operatorname{Tr}A}
    {\sqrt2}\text{ and } \mathcal Q_w(D)_i=
    \frac{w_i^TDw_i-\operatorname{Tr}D}{\sqrt2}.
  \]
  Then
  \[
    \mathcal Q_d(H_0+\rho D_0)=\mathcal Q_g(H_0)+\rho
    \mathcal Q_w(D_0).
  \]
  Define
  \[
    L(E_0,g)=\operatorname{span}\{\1_n\}+\{\mathcal Q_g(A):
    A\in\operatorname{Sym}(E_0)\}.
  \]
  After conditioning, this is a deterministic subspace
  independent of the tail samples, and
  \[
    \dim L(E_0,g)\le1+\frac{j(j+1)}2.
  \]
  The tail dimension is $p=d-j=d-O_K(1)$ and the nuisance
  dimension is $O_K(1)$. For all sufficiently large $d$,
  \[
    n-\dim L(E_0,g)\ge\alpha_-d^2-1-\frac{K(K+1)}2\ge
    \left(\frac14+\epsilon_0\right)p^2,
  \]
  and, because $p\ge d/2$,
  \[
    n\le\alpha_+d^2\le4\alpha_+p^2=C_1p^2.
  \]
  Put
  \[
    \bar a=\frac1{\rho n}\1_n^T\mathcal Q_d(S_0)\text{ and
    } v=-\mathcal Q_g(H_0/\rho)+\bar a\1_n.
  \]
  The block identity and the definition of $\PiZero$ give
  \[
    \mathcal Q_w(D_0)-v=\rho^{-1}\PiZero\mathcal Q_d(S_0).
  \]
  On the global continuity event, $\|\PiZero z\|_1\le2
  \|z\|_1$ and the rotation estimate give
  \[
    \frac1n\|\mathcal Q_w(D_0)-v\|_1 \le C_{\rm rot}
    \frac\zeta\delta+\frac1{\delta n}\|\PiZero\mathcal
    Q_d(S)\|_1.
  \]
  Conditional on $\mathcal G_{E_0}$, define
  \[
    \mathcal E_{\rm str}=\left\{ \inf_{\substack{D\succeq0,\
    \|D\|_F=1,\ \|D\|_{op}\le\tau\\ v\in L(E_0,g),\
    \|v\|_1\le Bn}} \frac1n\|\mathcal Q_w(D)-v\|_1\ge
    c_{\rm str} \right\}.
  \]
  Then \cref{thm:diffuse-restatement} and the
  tail-continuity estimate from
  \cref{expert-leave-one-row-module} give, for a constant
  $c_{\rm cond}>0$ depending only on the fixed parameters,
  \[
    \Pp\left((\mathcal E_{\rm str}\cap\{\Lambda_w\le
    A_*\})^c\mid\mathcal G_{E_0}\right)\le3e^{-c_{\rm
    cond}p^2},
  \]
  where
  \[
    \Lambda_w=\sup_{\|D\|_F=1}\frac1n\|\mathcal Q_w(D)\|_1.
  \]
  The bound is uniform in the realized nuisance subspace.

  Suppose that $n^{-1}\|\PiZero\mathcal Q_d(S)\|_1<
  c_{\rm cone}$. On the preceding good events,
  \[
    \frac1n\|\mathcal Q_w(D_0)-v\|_1 \le C_{\rm rot}
    \frac\zeta\delta+\frac{c_{\rm cone}}\delta \le
    \frac{c_{\rm str}}2.
  \]
  Moreover,
  \[
    \frac1n\|v\|_1 \le\frac1n\|\mathcal Q_w(D_0)\|_1+
    \frac1n\|\mathcal Q_w(D_0)-v\|_1 \le A_*+1<B.
  \]
  Thus $D_0$ and $v$ lie in the classes prohibited by
  \cref{thm:diffuse-restatement}, a contradiction.
  Uniformity over both classes permits $D_0$ and $v$ to be
  selected after observing the tail rows.

  The deterministic Grassmann nets constructed above
  satisfy
  \[
    \sum_{j=0}^K N(\operatorname{Gr}(j,d),\zeta) \le1+
    \sum_{j=1}^K \left(1+\frac{4\sqrt j}{\zeta}\right)^{dj}
    \le(K+1)e^{C_{K,\zeta}d}.
  \]
  Consequently, for all sufficiently large $d$,
  \[
    (K+1)e^{C_{K,\zeta}d}3e^{-c_{\rm cond}(d-K)^2}
    \le e^{-c_{\rm prob}d^2}
  \]
  for a fixed $c_{\rm prob}>0$. Adding the
  global-continuity and fixed-rank exceptional events
  proves the theorem with margin $c=c_{\rm cone}$ and a
  positive failure exponent depending only on $\alpha_-$
  and $\alpha_+$.

  If $T\in\mathcal F_d$, then $T\ne0$. The normalization
  $S=T/\|T\|_F$ satisfies $S\succeq0$, $\|S\|_F=1$, and
  $\PiZero\mathcal Q_d(S)=0$, contradicting the margin. 
    If $\liminf n/d^2>1/4$ but the aspect ratio is
  unbounded, choose a fixed $\alpha_0$ with
  \[
    \frac14<\alpha_0<\min\left\{\frac12,
    \liminf_d\frac{n}{d^2}\right\}.
  \]
  For all sufficiently large $d$,
  $\lfloor\alpha_0d^2\rfloor\le n(d)$. After
  discarding finitely many initial dimensions,
  which does not affect limits, retain
  deterministically the first
  $\lfloor\alpha_0d^2\rfloor$ observations. 
  By antitone monotonicity
  (Section~\ref{deterministic-framework}),
  $n_1\le n_2\Rightarrow
  \cF_{n_2}\subseteq\cF_{n_1}$, so
  $\cF_{n(d)}\subseteq
  \cF_{\lfloor\alpha_0d^2\rfloor}$; hence
  infeasibility with
  $\lfloor\alpha_0d^2\rfloor$ observations implies
  infeasibility with $n(d)$ observations.
 
  This proves Theorem~\ref{thm:main}(ii) for
  unbounded aspect from
  \cref{main-theorem-11-3}.
\end{proof}
% END INLINED FILE: expert_body_main.tex

\end{aiblock}

%%%%%%%%%%%%%%%%%%%%
\section{Conclusion}\label{conclusion}
%%%%%%%%%%%%%%%%%%%%

\begin{humanblock}


Part I constructs an exact positive definite fit whenever
the aspect ratio stays a fixed distance below $1/4$. Part
II proves that the centered kernel contains no nonzero
positive semidefinite matrix whenever the aspect ratio
stays a fixed distance above $1/4$. These are precisely the
two assertions of Theorem~\ref{thm:main}. The proof contains no
statement about the critical sequence $n/d^2\to1/4$.

\end{humanblock}

% \section*{Statement on AI Use}

% \begin{humanblock}

% Meta AI played a substantial collaborative role in the research leading
% to this manuscript. The proof was developed through many rounds of
% human-directed interaction rather than a single autonomous AI run.
% Meta AI proposed proof strategies, derived intermediate lemmas and
% calculations, explored alternative approaches, produced candidate
% arguments, responded to mathematical objections, and revised or
% replaced arguments when gaps were identified. Many proposed arguments
% did not survive further checking and were rejected, substantially
% changed, or replaced.

% Review and collaboration with other mathematicians were essential
% throughout this process. The author discussed the developing proof
% with reviewers and collaborators, shared drafts for examination,
% received mathematical objections and suggestions, and used this
% feedback to identify parts of the argument requiring further work.
% In several cases, this feedback led to new rounds of mathematical
% analysis, clarification, and revision.

% The author selected the problem, set the research goals, decided
% which proof directions to pursue or abandon, evaluated competing
% arguments, and determined which parts of the proof required further
% investigation. Meta AI outputs were not accepted simply because it described an argument as correct or complete. AI-generated
% claims of correctness were treated as research suggestions rather
% than mathematical verification.

% Meta AI also contributed to the preparation and revision of the mathematical exposition. The discovery, derivation, criticism, organization, and writing of parts of this manuscript emerged from a collaborative process involving Meta AI, the author, and reviewers. Because these contributions were closely intertwined throughout the project, no numerical percentages of contribution are assigned.

% The author takes responsibility for all mathematical claims in this
% manuscript. The final arguments are presented on the basis of the
% author's judgment together with extensive human review and
% collaboration. AI-generated arguments or assessments of correctness
% are not treated as substitutes for mathematical verification.

% \end{humanblock}





\section{Acknowledgment} 
The authors would like to thank Ricky Simanjuntak and Joseph Brennan for discussions at various stages of the project, and Ricky Simanjuntak for his feedback on the final draft.

%==================== Bibliography ====================

\bibliographystyle{plain}
\bibliography{bib}


% %%%%%%%%%%%%%%%%%%%%
% \bibliographystyle{amsplain}
% % BEGIN INLINED BIBLIOGRAPHY
% \providecommand{\bysame}{\leavevmode\hbox to3em{\hrulefill}\thinspace}
% \providecommand{\MR}{\relax\ifhmode\unskip\space\fi MR }
% % \MRhref is called by the amsart/book/proc definition of \MR.
% \providecommand{\MRhref}[2]{%
%   \href{http://www.ams.org/mathscinet-getitem?mr=#1}{#2}
% }

% %\providecommand{\href}[2]{#2}
% %%%%%%%%%%%%%%%%%%%%


% \begin{thebibliography}{1}

% \bibitem{AmbainisHarrowHastings2012}
%   Andris Ambainis, Aram~W. Harrow, and Matthew~B. Hastings,
%   \emph{Random tensor theory: Extending random matrix
%   theory to mixtures of random product states},
%   Communications in Mathematical Physics \textbf{310}
%   (2012), no.~1, 25--74.

%   \bibitem{AizenmanSimsStarr2003}
%   Michael Aizenman, Robert Sims, and Shannon~L. Starr,
%   \emph{An extended variational principle for the
%   Sherrington--Kirkpatrick spin-glass model},
%   Physical Review B \textbf{68} (2003), 214403.

% \bibitem{BandeiraMaillard2025}
%   Afonso~S. Bandeira and Antoine Maillard, \emph{Exact
%   threshold for approximate ellipsoid fitting of random
%   points}, Electronic Journal of Probability \textbf{30}
%   (2025), 1--46.

% \bibitem{DavidsonSzarek2001}
%   Kenneth~R. Davidson and Stanislaw~J. Szarek,
%   \emph{Local operator theory, random matrices and {Banach}
%   spaces}, Handbook of the Geometry of Banach Spaces,
%   vol.~1, Elsevier, 2001, See Theorem II.13 for the
%   Gaussian singular-value estimate, pp.~317--366.

% \bibitem{MisiakiewiczWen2026}
%   Theodor Misiakiewicz and Garrett~G. Wen,
%   \emph{The sharp SAT/UNSAT phase transition
%   in random ellipsoid fitting},
%   arXiv:2608.10184 (2026).

% \bibitem{Fernique1974}
%   Xavier Fernique, \emph{Minorations des fonctions
%   al\'eatoires gaussiennes}, Annales de l'Institut Fourier
%   \textbf{24} (1974), no.~2, 61--66. See Section 2.1 for
%   Sudakov minoration.

% \bibitem{Gordon1988}
%   Yehoram Gordon, \emph{On {M}ilman's inequality and
%   random subspaces which escape through a mesh in
%   {$\mathbb R^n$}}, Geometric Aspects of Functional
%   Analysis, Lecture Notes in Mathematics, vol. 1317,
%   Springer, 1988, pp.~84--106.

% \bibitem{KaneDiakonikolas2023}
%   Daniel~M. Kane and Ilias Diakonikolas,
%   \emph{A nearly tight bound for fitting an
%   ellipsoid to Gaussian random points},
%   Proceedings of the Thirty-Sixth Conference on
%   Learning Theory, Proceedings of Machine Learning
%   Research \textbf{195} (2023), 3014--3028.

%   \bibitem{MaillardKunisky2024}
% Antoine Maillard and Dmitriy Kunisky,
% \emph{Fitting an ellipsoid to random points: predictions using the
% replica method},
% IEEE Transactions on Information Theory \textbf{70} (2024),
% no.~10, 7273--7296.

% \bibitem{McDiarmid1989}
%   Colin McDiarmid, \emph{On the method of bounded
%   differences}, Surveys in Combinatorics, 1989, London
%   Mathematical Society Lecture Note Series, vol. 141,
%   Cambridge University Press, 1989, pp.~148--188.

% \bibitem{RudelsonVershynin2013}
%   Mark Rudelson and Roman Vershynin,
%   \emph{{Hanson--Wright} inequality and {Sub-Gaussian}
%   concentration}, Electronic Communications in Probability
%   \textbf{18} (2013), no.~82, 1--9.

% \bibitem{SaundersonParriloWillsky2013}
%   James Saunderson, Pablo~A. Parrilo, and
%   Alan~S. Willsky,
%   \emph{Diagonal and low-rank decompositions and
%   fitting ellipsoids to random points},
%   Proceedings of the 52nd IEEE Conference on
%   Decision and Control (2013), 6031--6036.

%   \bibitem{Talagrand2011}
%   Michel Talagrand,
%   \emph{Mean Field Models for Spin Glasses, Volume I:
%   Basic Examples},
%   Ergebnisse der Mathematik und ihrer Grenzgebiete,
%   vol.~54, Springer, 2011.

% \bibitem{vaart}
% A.~W.~van der Vaart, \emph{Asymptotic Statistics},
%   Cambridge University Press, 1998.

% \bibitem{Vershynin2018}
%   Roman Vershynin, \emph{High-dimensional probability: An
%   introduction with applications in data science},
%   Cambridge University Press, 2018.



% \end{thebibliography}

% END INLINED BIBLIOGRAPHY

\end{document}
